\chapter{Delta-forms and harmonic tropicalizations}
\index{delta-form}\index{harmonic tropicalization}
\label{chap:delta-forms}

Ordinary CLD forms are smooth on each tropical chart.  This is exactly the
right class for Stokes' theorem and local differential calculus, but it is
too small for the curvature of a typical model metric.  A piecewise affine
weight has zero Hessian on the interiors of its cells; its curvature lies on
the walls.  Thus it is naturally a polyhedral current, not a smooth form.

There are two complementary enlargements of the CLD calculus.  Gubler and
K\"unnemann adjoin polyhedral integration currents and obtain
\emph{delta-forms}.  Gubler, Jell, and Rabinoff enlarge the permitted
coordinate maps from tropicalizations by units to tropicalizations by
harmonic functions and obtain \emph{weakly smooth forms}.  The first
enlargement repairs a coefficient problem; the second repairs a coordinate
and cohomology problem.  They overlap, but neither should be silently
identified with the other.

\Crefrange{sec:delta-tropical}{sec:delta-positivity} follow
\cite{GK17,GK19}.  Their global algebraic statements are made over an
algebraically closed, complete, nontrivially valued field, as in those
papers.  \Crefrange{sec:harmonic-trop}{sec:weak-smooth} use
\cite{GJR26} and work on good strict analytic spaces over a complete
nontrivially valued field.  The final sections explain the more recent local
psh theory of \cite{GR25PL,GR25Classical}.

\section{Orientation}
\label{sec:ch18-orientation}

\subsection{What is added, and to what}
\label{sec:ch18-what-is-added}

Everything proved in \crefrange{chap:cld-forms}{chap:cld-metrized-bundles}
lives inside two categories: the sheaf $A^{\bullet,\bullet}_X$ of
\emph{smooth} forms, whose local model is a smooth Lagerberg form pulled
back along a moment map, and the space of \emph{currents}, its topological
dual.  The first is an algebra; the second is not.  Between the two there is
nothing.

The present chapter constructs two intermediate categories, each designed to
carry an operation that neither of the extremes supports.

\begin{itemize}
\item The \emph{delta-forms} of Gubler and K\"unnemann keep the moment maps
 of \Cref{chap:analytic-tropicalization} but enlarge the coefficients:
 alongside a smooth Lagerberg form one is allowed a weighted polyhedral
 cycle in the same tropical chart, and the two kinds of coefficient are
 multiplied by wedge product and by stable tropical intersection
 respectively.  The resulting objects are still integrable, still satisfy
 Stokes' formula, and---this is the point---still have a product.  A
 piecewise affine metric weight now has an honest \emph{form} as its
 curvature, not merely a current.
\item The \emph{weakly smooth forms} of Gubler, Jell and Rabinoff keep the
 smooth Lagerberg coefficients but enlarge the charts: instead of
 tropicalizing by invertible analytic functions one tropicalizes by
 \emph{harmonic} piecewise linear functions.  Since $-\log|f|$ is harmonic
 for invertible $f$ but not conversely, this is a genuine enlargement, and
 it repairs a cohomological defect that no enlargement of coefficients can
 repair.
\end{itemize}

The two enlargements are orthogonal.  \Cref{sec:weak-smooth} records that
fact in the form of a warning rather than of a comparison theorem, because
no comparison theorem is available; a reader who assumes that
``$\delta$-form'' and ``weakly smooth form'' are two names for the same
enlargement will make false deductions in both directions.

\subsection{The five questions of the chapter}
\label{sec:ch18-five-questions}

\begin{enumerate}
\item\label{itm:ch18-plan:defect}
 \emph{Why is the smooth category too small?}  Not because piecewise affine
 functions fail to be smooth---that is a triviality---but because the
 curvature of a psh function can be a measure that \emph{no} smooth form
 can approximate in the relevant sense.  \Cref{sec:ch18-defect} makes this
 quantitative: on a one-dimensional balanced tropical space every smooth
 function has vanishing curvature atom at every vertex
 (\Cref{prop:ch18-smooth-zero-atom}), yet there are psh functions, obtained
 as decreasing limits of smooth psh functions, whose atom at a vertex is a
 prescribed positive number (\Cref{ex:ch18-gauge-atom}).  The curvature is
 therefore not merely non-smooth: it is invisible to the smooth sheaf.
\item\label{itm:ch18-plan:deltaforms}
 \emph{What replaces the smooth category?}
 \Crefrange{sec:delta-tropical}{sec:ch18-delta-currents} construct
 delta-preforms, delta-forms and delta-currents, prove that on a fixed open
 subset of an ambient vector space the associated polyhedral current
 determines a delta-preform
 (\Cref{prop:ch18-current-faithful}) and that the product descends to that
 normalization (\Cref{prop:ch18-product-descends}).
\item\label{itm:ch18-plan:positivity}
 \emph{What is positivity in the enlarged calculus, and does
 Poincar\'e--Lelong survive?}  \Cref{sec:delta-positivity} splits positivity
 into a facewise Hessian condition and an effectivity condition on the
 corner locus; \Cref{sec:ch18-delta-pl} proves the delta-form
 Poincar\'e--Lelong formula and checks its constants against
 \Cref{app:normalizations} and against
 \Cref{thm:cld-poincare-lelong}.
\item\label{itm:ch18-plan:harmonic}
 \emph{What is a harmonic tropicalization, and how functorial is it?}
 \Crefrange{sec:harmonic-trop}{sec:weak-smooth} define harmonic PL
 functions by residual numerical triviality, prove that they are stable
 under pullback and scalar extension
 (\Cref{prop:ch18-harmonic-functorial}), prove the balancing theorem for the
 canonical weights (\Cref{thm:harmonic-balancing}), and construct the
 resulting sheaf of weakly smooth forms.
\item\label{itm:ch18-plan:model}
 \emph{How does all of this compare with the model functions of
 \Cref{chap:metrics}?}  \Cref{sec:ch18-model-comparison} proves the
 comparison: a model function is a $\mathbf Q$-PL function, its Chern
 delta-form is the corner cycle of its skeletal restriction, and the weights
 of that corner cycle are the vertical intersection numbers that
 \Cref{thm:cl-model-properties} attaches to the same divisor.  The chapter
 therefore closes the loop back to \Cref{chap:metrics}, where the same
 numbers were computed by pure intersection theory.
\end{enumerate}

\subsection{Standing hypotheses}
\label{sec:ch18-hypotheses}

Different sections of this chapter genuinely work over different base
categories, and the differences are not cosmetic.  We fix the following.

\begin{itemize}
\item In \crefrange{sec:ch18-defect}{sec:delta-tropical} the ambient object
 is a weighted balanced rational polyhedral complex in a real vector space
 $N_{\mathbf R}=N\otimes_{\mathbf Z}\mathbf R$, or an abstract tropical
 space in the sense of \Cref{def:ch10-tropical-space}.  No field is present.
 Every statement there is elementary polyhedral geometry and is proved in
 full.
\item In \crefrange{sec:ch18-delta-analytic}{sec:delta-positivity} the base
 field $K$ is algebraically closed, complete and nontrivially valued, and
 $X$ is an algebraic variety over $K$.  This is the scope of
 \cite{GK17,GK19}, and we do not silently extend it: the very affine
 covering, the canonical torus embedding and the curve-lifting theorem all
 use it.
\item In \crefrange{sec:harmonic-trop}{sec:ch18-classical} the base field
 $K$ is complete and nontrivially valued but not assumed algebraically
 closed, and $X$ is a good strict $K$-analytic space.  This is the scope of
 \cite{GJR26,GR25PL,GR25Classical}.
\item In \Cref{sec:ch18-model-comparison} the base field is complete and
 nontrivially valued, and is discretely valued wherever a special fibre with
 integer multiplicities is used; this is stated again at the point of use.
\end{itemize}

Signs, calibrations and lattice measures are those of
\Cref{chap:tropical-integration}; in particular
$\dprimeddprime$ of a convex function is a positive $(1,1)$-current by
\Cref{thm:convex-ddc-positive}, integration uses the lattice measure of
\Cref{def:ch11-lattice-measure}, and the Poincar\'e--Lelong normalization is
\eqref{eq:appJ-cld-pl}.

\subsection{The CLD core and the later extensions}
\label{sec:ch18-core-vs-extension}

\Cref{part:cld} contains one theory and two supplements, and the reader
should be able to tell at any moment which one is in force.  We use the
following convention throughout the chapter, and record it here because it
governs the source notes at the end of the chapter.

\begin{center}
\small
\begin{tabular}{p{.24\textwidth}|p{.34\textwidth}|p{.34\textwidth}}
\textbf{Layer}&\textbf{Objects}&\textbf{Where proved in this book}\\ \hline
CLD core
 &smooth forms, currents, Poincar\'e--Lelong, Bedford--Taylor products,
  Chern currents
 &\Crefrange{chap:cld-forms}{chap:cld-metrized-bundles}, following
  \cite{CLD25}\\
Gubler--K\"unnemann extension
 &delta-preforms, delta-forms, delta-currents, tropical intersection,
  Chern delta-forms, local heights
 &\Crefrange{sec:delta-tropical}{sec:delta-positivity} of this chapter,
  following \cite{GK17,GK19}\\
Gubler--Jell--Rabinoff extension
 &harmonic PL functions, harmonic tropicalizations, weakly smooth forms,
  strong currents
 &\Crefrange{sec:harmonic-trop}{sec:weak-smooth} of this chapter, following
  \cite{GJR26}\\
Gubler--Rabinoff curve tests
 &semipositive PL functions, classically psh functions
 &\Cref{sec:ch18-classical} and \Cref{sec:ch15-classical}, following
  \cite{GR25PL,GR25Classical}\\
\end{tabular}
\end{center}

No theorem of one layer is used inside the proof of a theorem of another
unless the citation is displayed.  In particular the delta-form
Poincar\'e--Lelong formula of \Cref{thm:gk-poincare-lelong} is not used to
prove \Cref{thm:cld-poincare-lelong}; the implication runs the other way,
and \Cref{cor:ch18-pl-restricts} makes it explicit.

\section{The wall-curvature defect}
\label{sec:ch18-defect}

\subsection{The one-dimensional computation}
\label{sec:ch18-defect-onedim}

Consider a one-dimensional balanced tropical cycle \(C\) and a continuous
piecewise affine function \(\varphi\).  On every open edge,
\(\dprimeddprime\varphi=0\) as a smooth form.  At a vertex \(v\), however, integration
by parts gives
\begin{equation}
 \dprimeddprime[\varphi]
   =\sum_v\left(\sum_{e\ni v}m_e
             \frac{\partial\varphi}{\partial n_{v,e}}\right)\delta_v.
\label{eq:delta-graph-laplacian}
\end{equation}
Here \(m_e\) is the edge multiplicity and the derivative is the outgoing
slope.  The right-hand side is zero only when \(\varphi\) is harmonic.
Formula \eqref{eq:delta-graph-laplacian} is not new here: it is
\Cref{lem:ch15-vertex-atom} specialized to piecewise affine $\varphi$, whose
edgewise second derivative vanishes, and it is the graph Laplacian
\eqref{eq:graph-laplacian} when all weights are $1$.  What is new is the use
we make of it.

We shall call the number
\begin{equation}
\label{eq:ch18-atom}
 \operatorname{at}_v(\varphi)
 :=\sum_{e\ni v}m_e\,d_e\varphi(v)
\end{equation}
the \emph{atom} of $\varphi$ at $v$, $d_e\varphi(v)$ denoting the outgoing
slope of $\varphi$ along $e$ measured in the lattice length of
\Cref{def:ch11-lattice-measure}.  Thus
$\dprimeddprime[\varphi]=\sum_v\operatorname{at}_v(\varphi)\delta_v$ for
piecewise affine $\varphi$.

\subsection{Smooth functions have no atoms}
\label{sec:ch18-zero-atom}

The following elementary observation is the reason the whole chapter exists.
It is a restatement of \Cref{ex:ch10-smooth-at-vertex}, isolated as a
proposition because it is used four times below.

\begin{proposition}
\label{prop:ch18-smooth-zero-atom}
Let $C$ be a one-dimensional weighted rational polyhedral complex in
$N_{\mathbf R}$ which is balanced at the vertex $v$, and let $u$ be a smooth
function on a neighbourhood of $v$ in $|C|$, that is, the restriction to
$|C|$ of a smooth function $F$ defined on a neighbourhood of $v$ in
$N_{\mathbf R}$.  Then
\[
 \operatorname{at}_v(u)=0 .
\]
Consequently $\dprimeddprime[u]$ carries no mass at $v$, and if $u$ is
moreover affine on each edge through $v$ then $\dprimeddprime[u]=0$ near
$v$.
\end{proposition}

\begin{proof}
Translate so that $v=0$, and let $u_e\in N$ denote the primitive generator
of the edge $e$ at $v$, so that the lattice-length parameter along $e$ is
$s\mapsto su_e$ by \Cref{def:ch11-lattice-measure}.  The outgoing slope of
$u=F|_{|C|}$ along $e$ is then
\[
 d_eu(v)=\frac{\rd}{\rd s}\Big|_{s=0^+}F(su_e)
 =\bigl\langle \rd F(v),u_e\bigr\rangle ,
\]
by the chain rule; the derivative exists because $F$ is smooth on a
neighbourhood of $v$ in $N_{\mathbf R}$, and it does not depend on which
smooth ambient representative $F$ is chosen, since two representatives have
the same restriction to each edge.  Therefore, by linearity of
$\rd F(v)$,
\[
 \operatorname{at}_v(u)
 =\sum_{e\ni v}m_e\bigl\langle \rd F(v),u_e\bigr\rangle
 =\Bigl\langle \rd F(v),\ \sum_{e\ni v}m_eu_e\Bigr\rangle
 =\bigl\langle \rd F(v),0\bigr\rangle=0,
\]
the vanishing of $\sum_em_eu_e$ being the balancing condition of
\Cref{def:ch10-weighted-balanced} at the zero-dimensional face $\{v\}$.  The
last two assertions follow from \Cref{lem:ch15-vertex-atom}: the atom of
$\dprimeddprime[u]$ at $v$ is $\operatorname{at}_v(u)$, and the edgewise term
is $u''\,\rd s$, which vanishes when $u$ is affine on each edge.
\end{proof}

\begin{remark}[What the proposition does and does not say]
\label{rem:ch18-zero-atom-reading}
Three readings must be kept apart.
\begin{enumerate}
\item The statement is about the \emph{sheaf of smooth functions} on the
 tropical space $|C|$, as defined in \Cref{def:ch10-forms-on-support}: a
 smooth function near $v$ is by definition the restriction of one ambient
 smooth function.  It is not a statement about functions that happen to be
 $C^\infty$ on each edge; such a function has an arbitrary atom, as
 \eqref{eq:delta-graph-laplacian} shows.
\item On a tropical space with several charts the same conclusion holds,
 because by \Cref{thm:ch10-transition-transport}\,\cref{itm:ch10-tr-lattice}
 the local weighted star of a point is well defined up to integer-linear
 isomorphism, and both the atom and the balancing relation are computed in
 that star.  Balancing is part of the definition of a tropical space.
\item The proposition says nothing about a \emph{Berkovich} space, where the
 chart may be refined: a function that is not the pullback of a smooth
 function under one moment map may be the pullback of a smooth function
 under a finer one.  \Cref{sec:harmonic-trop} is precisely about the
 interplay between refinement and the vanishing of atoms; the objects there
 are harmonic, that is, of zero atom, and the question is whether they are
 charts.
\end{enumerate}
\end{remark}

\subsection{A psh function with a prescribed positive atom}
\label{sec:ch18-gauge}

We now produce, on the smallest interesting tropical space, a psh function
whose atom is a prescribed positive number and which is nevertheless a
decreasing limit of smooth psh functions.  By
\Cref{prop:ch18-smooth-zero-atom} the atom of each approximant is zero, so
the entire mass of the limiting curvature is created in the limit.  Every
step is carried to an explicit number.

\begin{example}[Minkowski gauge on the standard tropical line]
\label{ex:ch18-gauge-atom}
\index{tropical line!Minkowski gauge on}%
Let $N=\mathbf Z^2$ and let
\begin{equation}
\label{eq:ch18-gauge-rays}
 v_1=(1,0),\qquad v_2=(0,1),\qquad v_3=(-1,-1),
\end{equation}
and let $L\subset\mathbf R^2$ be the union of the three rays
$\mathbf R_{\geq0}v_i$, each with weight $1$.  This is the standard tropical
line of \Cref{ex:ch10-standard-line} in the coordinates in which it arises
analytically; it is balanced, since $v_1+v_2+v_3=0$.  Recall from
\Cref{ex:trop-line-in-torus} that $L=\Trop(U)$ for
$U=\{1+x+y=0\}\cap\mathbf G_{m,K}^2$, the tropical coordinates being the
valuations $(v(x),v(y))$.

\emph{The gauge.}  Let $T=\operatorname{conv}\{v_1,v_2,v_3\}$, a triangle
with $0$ in its interior, and let
\[
 g(z)=\inf\{\lambda>0:z\in\lambda T\}
\]
be its Minkowski gauge.  Then $g$ is convex, positively homogeneous of
degree one, and finite on all of $\mathbf R^2$.  Being positively
homogeneous and convex, $g$ is the support function of the polar body
$T^\circ=\{y:\langle y,z\rangle\leq1\ \text{for all }z\in T\}$; and since
$T$ is the convex hull of $v_1,v_2,v_3$,
\[
 T^\circ=\{y:\langle y,v_i\rangle\leq1,\ i=1,2,3\}
 =\{y:y_1\leq1,\ y_2\leq1,\ -y_1-y_2\leq1\},
\]
which is the triangle with vertices $(1,1)$, $(1,-2)$, $(-2,1)$.  Hence
\begin{equation}
\label{eq:ch18-gauge-maxplus}
 g(z)=\max\{z_1+z_2,\ z_1-2z_2,\ -2z_1+z_2\},
\end{equation}
a maximum of three integral linear forms: $g$ is an integral piecewise
affine, indeed max-plus polynomial, function on $\mathbf R^2$.

\emph{Its slopes on $L$.}  Each $v_i$ lies on $\partial T$, so $g(v_i)=1$
and, by homogeneity, $g(tv_i)=t$ for $t\geq0$.  We check this against
\eqref{eq:ch18-gauge-maxplus}:
\[
 \begin{aligned}
 g(v_1)&=\max\{1,1,-2\}=1,\\
 g(v_2)&=\max\{1,-2,1\}=1,\\
 g(v_3)&=\max\{-2,1,1\}=1 .
 \end{aligned}
\]
Each $v_i$ is primitive, so the lattice-length parameter along the $i$-th
ray is $t$, and the outgoing slope of $u:=g|_L$ at the origin along that ray
is $1$.  Therefore, by \eqref{eq:delta-graph-laplacian},
\begin{equation}
\label{eq:ch18-gauge-atom}
 \dprimeddprime[u]=\operatorname{at}_0(u)\,\delta_0=3\,\delta_0 .
\end{equation}
Note that $u$ is affine on each open ray, so there is no edgewise
contribution: the entire curvature is the atom.

\emph{It is not smooth.}  By \Cref{prop:ch18-smooth-zero-atom} the atom of a
smooth function vanishes, so $u$ is not smooth at $0$; indeed no smooth
function on $L$ has $\dprimeddprime$ equal to a nonzero multiple of
$\delta_0$, because a smooth function has smooth $\dprimeddprime$ and the
smooth $(1,1)$-forms on $L$ near $0$ are, by
\Cref{def:ch10-forms-on-support}, given cellwise by densities against the
lattice length, hence charge no point.

\emph{It is a decreasing limit of smooth psh functions.}  Fix a radial
mollifier $\rho$ on $\mathbf R^2$ and set $G_\varepsilon=g*\rho_\varepsilon$
and $u_\varepsilon=G_\varepsilon|_L$.  By
\Cref{prop:ch09-radial-mollification} the family $G_\varepsilon$ consists of
smooth convex functions and decreases to $g$ locally uniformly as
$\varepsilon\downarrow0$; by \Cref{thm:convex-ddc-positive},
$\dprimeddprime G_\varepsilon\geq0$ on $\mathbf R^2$.  Positivity of
$\dprimeddprime$ is preserved by affine pullback
(\Cref{prop:ch09-positivity-inclusions} and
\Cref{prop:ch09-pullback-functorial}), and each cell of $L$ is contained in
a line, so the restriction of $G_\varepsilon$ to each ray is a smooth convex
function of the lattice parameter; by the chart description of forms on
$|L|$ in \Cref{def:ch10-forms-on-support}, this says exactly that
$\dprimeddprime u_\varepsilon\geq0$ on $L$.  Hence each $u_\varepsilon$ is a \emph{smooth} psh
function on $L$, and $u_\varepsilon\downarrow u$ locally uniformly.  By
\Cref{def:ch16-psh-filtration}, $u$ is a locally approachable continuous psh
function on $L$.

\emph{Where the mass goes.}  By \Cref{prop:ch18-smooth-zero-atom},
$\operatorname{at}_0(u_\varepsilon)=0$ for every $\varepsilon>0$.  On the
other hand, homogeneity gives the exact value of $u_\varepsilon$ at the
vertex,
\begin{equation}
\label{eq:ch18-gauge-vertex-value}
 u_\varepsilon(0)=\int_{\mathbf R^2}g(-y)\rho_\varepsilon(y)\,\rd y
 =c_\rho\,\varepsilon,
 \qquad
 c_\rho:=\int_{\mathbf R^2}g(-z)\rho(z)\,\rd z>0 ,
\end{equation}
so $u_\varepsilon(0)\downarrow0=u(0)$ linearly in $\varepsilon$.  Fix
$R>\delta>0$ and let $B_r=\{z\in L:|z|\leq r\}$.  Since $t\mapsto
G_\varepsilon(tv_i)$ is convex, the measure $\dprimeddprime[u_\varepsilon]$
is nonnegative, and its mass on the part of the $i$-th ray between the
parameters $a<b$ is
\[
 \frac{\rd}{\rd t}\Big|_{t=b}G_\varepsilon(tv_i)
 -\frac{\rd}{\rd t}\Big|_{t=a}G_\varepsilon(tv_i)
 =:s_i^\varepsilon(b)-s_i^\varepsilon(a).
\]
As $\varepsilon\downarrow0$ we have $s_i^\varepsilon(t)\to1$ for every
$t>0$, because $g$ is affine with slope $1$ along the open ray and
$G_\varepsilon\to g$ locally uniformly, so that the derivatives of the
convex functions $G_\varepsilon(\cdot\,v_i)$ converge at every point of
differentiability of the limit.  Consequently
\[
 \bigl\|\dprimeddprime[u_\varepsilon]\bigr\|_{B_R\setminus B_\delta}
 =\sum_{i=1}^3\bigl(s_i^\varepsilon(R)-s_i^\varepsilon(\delta)\bigr)
 \longrightarrow 3-3=0,
\]
while, using $\sum_is_i^\varepsilon(0^+)=\operatorname{at}_0(u_\varepsilon)=0$,
\[
 \bigl\|\dprimeddprime[u_\varepsilon]\bigr\|_{B_\delta}
 =\sum_{i=1}^3\bigl(s_i^\varepsilon(\delta)-s_i^\varepsilon(0^+)\bigr)
 =\sum_{i=1}^3s_i^\varepsilon(\delta)
 \longrightarrow 3 .
\]
Thus the total mass on $B_R$ tends to $3$ and all of it concentrates at the
origin:
\begin{equation}
\label{eq:ch18-gauge-concentration}
 \dprimeddprime[u_\varepsilon]\longrightarrow 3\,\delta_0
 \qquad\text{weakly on }B_R,
\end{equation}
in agreement with \eqref{eq:ch18-gauge-atom} and with the continuity theorem
\Cref{thm:cld-bt-continuity}.
\end{example}

\begin{remark}[Reading the example]
\label{rem:ch18-gauge-reading}
The example is not a counterexample to anything: $u$ is psh, its curvature
is a positive measure, and the approximation theorem holds.  What it shows
is a structural fact about the \emph{target} of the curvature operator.  The
map $u\mapsto\dprimeddprime u$ sends smooth functions to smooth
$(1,1)$-forms, and on a one-dimensional balanced tropical space every smooth
$(1,1)$-form is absolutely continuous with respect to the lattice length.
The curvature of the psh function $u$ of \Cref{ex:ch18-gauge-atom} is purely
atomic.  So the closure of the smooth category under the operations we care
about---decreasing limits of potentials---does not stay inside the smooth
category on the curvature side, and the discrepancy is not a matter of
regularity but of \emph{support}: a whole $3$ units of mass have moved onto
a set of measure zero.

The remedy adopted in \Cref{chap:cld-currents} was to enlarge the target to
currents.  That works for the linear theory and fails for the nonlinear one,
since $\delta_0\wedge\delta_0$ is meaningless.  The remedy adopted in this
chapter is to enlarge the target to \emph{delta-forms}, in which
$3\delta_0$ is not a mere current but a form of bidegree $(1,1)$ with
polyhedral coefficient---the weighted cycle $3\{0\}$---which can be wedged
with other such coefficients by stable tropical intersection.  Both remedies
are needed; neither subsumes the other.
\end{remark}

\begin{remark}[Provenance]
\label{rem:ch18-gauge-provenance}
\Cref{prop:ch18-smooth-zero-atom} and \Cref{ex:ch18-gauge-atom} are
elementary and are recorded here for orientation only.  Neither is a theorem
of this book and no priority is claimed for either.  The vanishing of the
atom of a smooth function is the content of \Cref{ex:ch10-smooth-at-vertex};
the gauge computation is standard convex duality
(\Cref{app:convexity}); the concentration statement is
\Cref{thm:cld-bt-continuity} together with the elementary
one-dimensional convexity used above.  The pair is included because it makes
the abstract sentence ``the smooth category does not contain all
model-theoretic curvature'' into a computation with the number $3$ in it.
\end{remark}

\subsection{Higher dimensions: the corner locus}
\label{sec:ch18-corner-locus}

In higher dimensions, the same integration-by-parts calculation produces a
weighted polyhedral cycle of codimension one: the \emph{corner locus} of
\(\varphi\).  No nonzero ordinary smooth \((1,1)\)-form can be supported only
on that locus.  Passing merely from forms to all currents would retain the
object but lose a usable wedge product.  Delta-forms are the intermediate
class designed to keep both the wall current and tropical intersection.

The next example is the two-dimensional analogue of
\Cref{ex:ch18-gauge-atom}, and it is the one that shows why the coefficient
we must allow is a \emph{cycle} and not merely a measure.

\begin{example}[The corner locus of a tropical polynomial]
\label{ex:ch18-corner-max}
On $\mathbf R^2$ put $\varphi(x)=\max\{0,x_1,x_2\}$.  On each of the three
open regions
\[
 \{x_1<0,x_2<0\},\qquad\{x_1>0,x_1>x_2\},\qquad\{x_2>0,x_2>x_1\}
\]
the function $\varphi$ is affine, so $\dprimeddprime\varphi=0$ there.  The
corner locus, in the sense of \Cref{def:appG-corner-locus}, is the union of
the three rays
\[
 \sigma_1=\{x_1=0\geq x_2\},\quad
 \sigma_2=\{x_2=0\geq x_1\},\quad
 \sigma_3=\{x_1=x_2\geq0\},
\]
each of which we equip with the weight $1$; the fan
$C_\varphi=\sigma_1\cup\sigma_2\cup\sigma_3$ has primitive generators
$(0,-1)$, $(-1,0)$ and $(1,1)$, whose sum is $0$, so $C_\varphi$ is a
balanced one-dimensional cycle---again a standard tropical line, now in a
different coordinate presentation.

To see that the weight is $1$ and not something else, compute the jump of
$\rd\varphi$ across $\sigma_1$: on the side $x_1<0$ one has
$\rd\varphi=0$, and on the side $x_1>0$ one has $\rd\varphi=\rd x_1$; the
tangent lattice of $\sigma_1$ is $\mathbf Z(0,1)$, a lift of the primitive
normal pointing into the second region is $(1,0)$, and the jump in the
normal derivative is $1-0=1$.  The same computation on $\sigma_2$ gives
$1$.  On $\sigma_3$, whose tangent lattice is $\mathbf Z(1,1)$, subtract the
affine function $x_2$, which changes no corner weight by
\Cref{lem:ch18-corner-welldefined}\,\cref{itm:ch18-cc:affine}: near a generic
point of $\sigma_3$ the function becomes $\max\{0,x_1-x_2\}$, whose gradient
jumps by $\rd x_1-\rd x_2$ across the wall, and this pairs to $1$ with the
lift $(1,0)$ of the primitive normal, since $(a,b)\mapsto a-b$ identifies
$\mathbf Z^2/\mathbf Z(1,1)$ with $\mathbf Z$ and sends $(1,0)$ to $1$.
Hence
\begin{equation}
\label{eq:ch18-corner-max}
 \dprimeddprime[\varphi]=\delta_{C_\varphi},
\end{equation}
the integration current of \Cref{def:ch11-integration-current} over the
weighted fan $C_\varphi$.

Now observe two things.  First, $\delta_{C_\varphi}$ is a $(1,1)$-current
which is not represented by any smooth $(1,1)$-form: it charges a set of
Lebesgue measure zero.  Second---and this is the reason for the enlarged
calculus---the polyhedral subclass carries a multiplication that the
category of all currents does not.  On this fixed ambient space the current
does determine the collected cycle coefficients; the gain is algebraic
structure, not hidden linear data.  Wedging $\delta_{C_\varphi}$ with itself
is meaningless as a product of arbitrary currents, but the \emph{cycle}
$C_\varphi$ has a stable
self-intersection $C_\varphi\cdot C_\varphi$, which by
\Cref{thm:tropical-stable-intersection} is the $0$-cycle $1\cdot\{0\}$: the
tropical B\'ezout count of two lines.  Delta-forms are designed so that
\[
 \bigl(\dprimeddprime\varphi\bigr)^{\wedge2}=\delta_{C_\varphi\cdot C_\varphi}
 =\delta_0
\]
is a legitimate identity, and \Cref{ex:ch16-fan-p2} confirms that the
Bedford--Taylor product of \Cref{chap:cld-bedford-taylor} produces the same
measure $\delta_0$, of mass $1$.
\end{example}

\begin{remark}[Why not simply pass to currents]
\label{rem:ch18-why-not-currents}
The linear theory of \Cref{chap:cld-currents} already contains
$\delta_{C_\varphi}$, so nothing is missing there.  What is missing is a
product.  Currents form a module over smooth forms and nothing more:
\Cref{prop:ch14-no-pullback} records that even pullback fails, and the
one-variable obstruction $\delta_0\cdot\delta_0$ recalled in
\Cref{sec:ch16-obstruction} shows that no bilinear extension of the wedge
product to all currents exists.  The Bedford--Taylor recursion of
\Cref{chap:cld-bedford-taylor} circumvents the obstruction, but at the price
of carrying the potentials along: its output is a current attached to a
tuple of potentials, and \Cref{ex:ch16-cross-fan} shows that the output can
genuinely depend on the potential and not only on its curvature.  The
delta-calculus resolves the same difficulty from the other side.  It does
not multiply arbitrary currents.  Delta-preforms are polyhedral currents;
one computes their product from any polyhedral presentation, and
\Cref{prop:ch18-product-descends} proves that the answer is independent of
that presentation.
\end{remark}

\subsection{Model functions on a curve, with numbers}
\label{sec:ch18-defect-model}

\begin{example}[A model function]
\label{ex:ch18-model-function-defect}
Let \(\mathfrak X\) be a semistable formal curve and let \(D\) be a vertical
Cartier divisor.  Its model function is affine on each edge of the skeleton.
The current \(\dprimeddprime[\varphi_D]\) is the graph Laplacian, with the outgoing
slopes determined by degrees of \(\mathcal O(D)\) on special-fibre
components.  Unless all those degrees vanish, \(\dprimeddprime[\varphi_D]\) is not an
ordinary smooth form.
\end{example}

The general assertion of \Cref{ex:ch18-model-function-defect} is best read
against a case in which all the numbers are available.

\begin{example}[The chain of three components]
\label{ex:ch18-chain-numbers}
Take the semistable chain model $\widetilde{\mathcal C}$ of
\Cref{ex:metrics-cl-curve-chain}, whose special fibre is a chain
$C_0,E,C_2$ of projective lines with all multiplicities $1$, and whose
skeleton is the segment $\Gamma=[0,2]$ with the three divisorial points
$v_0,v_1,v_2$ at the integers.  Let $D=a_0C_0+a_1E+a_2C_2$ be a vertical
divisor.  By
\Cref{prop:metrics-divisor-model-function}\,(2), the model function
$\varphi_D$ takes the values $a_0,a_1,a_2$ at $v_0,v_1,v_2$ and is affine on
each of the two edges, all edge weights being $1$.  The three atoms of
\eqref{eq:ch18-atom} are
\begin{equation}
\label{eq:ch18-chain-atoms}
 \operatorname{at}_{v_0}(\varphi_D)=a_1-a_0,\quad
 \operatorname{at}_{v_1}(\varphi_D)=(a_0-a_1)+(a_2-a_1),\quad
 \operatorname{at}_{v_2}(\varphi_D)=a_1-a_2,
\end{equation}
so that
\[
 \dprimeddprime[\varphi_D]
 =(a_1-a_0)\delta_{v_0}
 +(a_0+a_2-2a_1)\delta_{v_1}
 +(a_1-a_2)\delta_{v_2},
\]
of total mass $0$, in accordance with \Cref{cor:ch15-degree-zero}.  These are
exactly the three numbers subtracted in \eqref{eq:metrics-cl-chain-measure},
where they appear as the difference between the twisted and untwisted
Chambert--Loir measures.  Concretely, for $a=(0,1,0)$ one finds
\[
 \dprimeddprime[\varphi_D]=\delta_{v_0}-2\delta_{v_1}+\delta_{v_2},
\]
and for $a=(0,0,1)$ one finds $-\delta_{v_1}+\delta_{v_2}$; neither is an
absolutely continuous measure on $\Gamma$, so neither is represented by a
smooth $(1,1)$-form.  For $a=(0,c,2c)$ with any $c$ the function
$\varphi_D$ is affine on $\Gamma$, all three atoms vanish except at the
endpoints, and one gets $c\,\delta_{v_0}-c\,\delta_{v_2}$: harmonic in the
interior, but still not smooth, because $\Gamma$ has boundary at $v_0$ and
$v_2$ and \Cref{prop:ch18-smooth-zero-atom} does not apply there.
\end{example}

\begin{example}[The blowup of $\mathbf P^1$, revisited]
\label{ex:ch18-p1-blowup-defect}
In \Cref{ex:metrics-p1-blowup} the model function of the exceptional divisor
of one blowup was computed to be the hat function
\eqref{eq:metrics-p1-blowup-hat} on the standard apartment, and its
Laplacian was found in \eqref{eq:metrics-p1-blowup-laplacian} to be
\[
 \dprimeddprime[\varphi_E]=\delta_{\zeta_{0,1}}-\delta_{x_E}.
\]
Here the tropical space is an interval with unit weights and two boundary
points, and $\varphi_E$ has slope $+1$ on $[0,1]$ and slope $0$ elsewhere.
The two atoms $+1$ and $-1$ are the smallest possible nonzero integers, and
they already exhibit the phenomenon of \Cref{sec:ch18-zero-atom}: the
curvature of the simplest nontrivial model function on the simplest possible
curve is a pair of point masses.  The delta-form of
\Cref{sec:ch18-chern-delta} attached to this metric is the weighted
$0$-cycle $\{\zeta_{0,1}\}-\{x_E\}$, and its total mass $0$ is the
degree-zero statement \Cref{cor:ch15-degree-zero}.
\end{example}

\section{Tropical cycles with smooth weights}
\label{sec:delta-tropical}

\subsection{Cycles, weights, and their integration currents}
\label{sec:ch18-cycles-weights}

Let \(N\) be a lattice, \(N_{\mathbf R}=N\otimes\mathbf R\), and let
\(\Omega\subset N_{\mathbf R}\) be open.  A pure \(d\)-dimensional
polyhedral complex \(C\) has a real weight \(m_\sigma\) on every maximal
face.  Constant weights give the usual tropical cycles.  Gubler--K\"unnemann
also allow weights that are smooth up to the faces of every maximal
polyhedron, with the corresponding weighted balancing condition.

The integration current \(\delta_C\) acts on a compactly supported
\((d,d)\)-superform by
\[
                \langle\delta_C,\alpha\rangle
                  =\sum_{\sigma\in C_d}
                     \int_\sigma m_\sigma\alpha.
\]
Changing the polyhedral subdivision does not change this current.

We record the last sentence as a lemma, because it is the first of the two
invariance statements that make the whole construction possible and because
its proof isolates the exact role of the lattice measure.

\begin{lemma}
\label{lem:ch18-subdivision-invariance}
Let $C$ be a weighted pure $d$-dimensional polyhedral complex in $\Omega$
and let $C'$ be a subdivision of $C$, with the induced weights
$m'_{\sigma'}=m_\sigma$ for $\sigma'\subseteq\sigma$ maximal.  Then
$\delta_{C'}=\delta_C$ as currents, and $C$ is balanced if and only if $C'$
is.
\end{lemma}

\begin{proof}
Fix a compactly supported $(d,d)$-form $\alpha$.  For each maximal face
$\sigma$ of $C$, the maximal faces $\sigma'$ of $C'$ contained in $\sigma$
cover $\sigma$ and meet pairwise in sets of dimension $<d$, hence of
$\mu_\sigma$-measure zero, $\mu_\sigma$ being the lattice measure of
\Cref{def:ch11-lattice-measure}.  Moreover $\sigma'$ and $\sigma$ have the
same tangent lattice $N_\sigma$, since $\sigma'$ has nonempty interior in
$\sigma$; therefore $\mu_{\sigma'}$ is the restriction of $\mu_\sigma$.
Additivity of the integral gives
$\sum_{\sigma'\subseteq\sigma}\int_{\sigma'}m_\sigma\alpha
=\int_\sigma m_\sigma\alpha$, and summing over $\sigma$ proves
$\delta_{C'}=\delta_C$.  This is
\Cref{prop:integral-subdivision-independent}.

For the second assertion, let $\tau'$ be a codimension-one face of $C'$.  If
$\tau'$ lies in the relative interior of a maximal face $\sigma$ of $C$,
exactly two maximal faces of $C'$ contain $\tau'$, with the same weight
$m_\sigma$ and opposite primitive normals, so the balancing sum at $\tau'$
is $m_\sigma(u+(-u))=0$ automatically.  If $\tau'$ lies in a codimension-one
face $\tau$ of $C$, then the maximal faces of $C'$ through $\tau'$ are in
weight-preserving, normal-preserving bijection with the maximal faces of $C$
through $\tau$, so the two balancing sums coincide.  This is
\Cref{prop:ch10-balancing-subdivision}.
\end{proof}

\begin{remark}[Two kinds of weight]
\label{rem:ch18-two-weights}
Gubler and K\"unnemann permit the weight $m_\sigma$ to be a function smooth
up to the faces of $\sigma$, rather than a constant, subject to the balancing
condition
$\sum_{\sigma\supset\tau}m_\sigma(x)u_{\sigma/\tau}\in N_{\tau,\mathbf R}$ for every
$x\in\RelInt\tau$.  The ambient cycles used below have constant weights
unless smooth weights are explicitly mentioned.  The corner locus of a
piecewise smooth function, however, generally has smooth weights, so its
balancing and differential calculus use the smooth-weight theory of
\cite[\S\S1--3]{GK17}.  In particular, with smooth weights the differential
$\dprime\delta_C$ acquires an extra term from $\rd m_\sigma$, and
$\delta_C$ need not be closed.  Constant weights and balancing give
$\dprime\delta_C=\ddprime\delta_C=0$ by \Cref{thm:ch11-current-closed}.
\end{remark}

\subsection{Delta-preforms}
\label{sec:ch18-delta-preforms}

\begin{definition}[Delta-preform]
A \emph{delta-preform} of bidegree \((p,q)\) and codimension \(l\) on
\(\Omega\) is a finite sum
\begin{equation}
                       \sum_i \alpha_i\wedge\delta_{C_i},
\label{eq:delta-preform}
\end{equation}
where \(C_i\) is a tropical cycle with smooth weights of codimension \(l\) and
\(\alpha_i\in A^{p-l,q-l}(\Omega)\) is an ambient smooth superform.
Finite sums over different codimensions are called
\emph{generalized delta-preforms}.
\end{definition}

The codimension shift makes \(\delta_C\) itself have bidegree \((l,l)\).
Ordinary smooth superforms embed as the codimension-zero terms
\(\alpha\wedge\delta_{N_{\mathbf R}}\) with unit constant weight.  They are
not all codimension-zero delta-preforms: allowing smooth weights on a
full-dimensional polyhedral complex produces the piecewise-smooth
superforms of \cite[Example~2.10]{GK17}.

Two elementary conventions must be fixed at once, since they are the source
of most confusion in reading \cite{GK17}.  First, only the restriction of
$\alpha_i$ to $|C_i|\cap\Omega$ affects
\eqref{eq:delta-preform}.  Thus changing an ambient coefficient away from
the cycle changes no current; conversely, a coefficient initially given on
a neighbourhood of the cycle can be handled after restriction and a smooth
partition-of-unity extension.  Second, the expression
$\alpha_i\wedge\delta_{C_i}$ already denotes the action of a smooth form on
the integration current of $C_i$.  Thus a delta-preform is the resulting
polyhedral supercurrent, not a raw formal presentation.  What still requires
proof is that stable intersection gives a product depending only on that
current and not on the chosen presentation; this is
\Cref{prop:ch18-product-descends,thm:delta-preform-algebra}.

\begin{definition}[The associated current]
\label{def:ch18-associated-current}
The \emph{current associated} to the delta-preform
$\beta=\sum_i\alpha_i\wedge\delta_{C_i}$ of bidegree $(p,q)$ is
\begin{equation}
\label{eq:ch18-associated-current}
 \langle[\beta],\omega\rangle
 :=\sum_i\sum_{\sigma\in(C_i)_{d_i}}
    \int_\sigma m_{i,\sigma}\,\alpha_i\wedge\omega ,
\end{equation}
$\omega$ running over compactly supported superforms of bidegree
$(n-p,n-q)$, where $n=\dim N_{\mathbf R}$ and $d_i=n-l_i$.  The integrals
are taken with the lattice measures of \Cref{def:ch11-lattice-measure}, and
the restriction of $\alpha_i\wedge\omega$ to $\sigma$ is taken in the sense
of \Cref{def:ch10-forms-on-support}.
\end{definition}

The assignment $\beta\mapsto[\beta]$ is linear and, by
\Cref{lem:ch18-subdivision-invariance}, unchanged under subdivision of the
$C_i$.  A displayed sum is only a presentation of this current; the next
proposition gives the exact equality criterion after a common subdivision.

\begin{proposition}
\label{prop:ch18-current-faithful}
Let $\beta=\sum_i\alpha_i\wedge\delta_{C_i}$ be a generalized delta-preform
of bidegree $(p,q)$ on $\Omega$, and choose a polyhedral complex
$\mathscr C$ in $\Omega$ subdividing every $C_i$
(\Cref{prop:tropical-common-refinement}).  For a face $\sigma\in\mathscr C$
of dimension $d=n-l$ write
\[
 A_\sigma:=\sum_{i:\,l_i=l}m_{i,\sigma}\,\alpha_i\big|_\sigma
 \ \in\ A^{p-l,q-l}(\sigma),
\]
the sum over those $i$ of codimension $l$ for which $\sigma$ is contained in
a maximal face of $C_i$, with $m_{i,\sigma}$ the corresponding weight.  Then
\[
 [\beta]=0
 \quad\Longleftrightarrow\quad
 A_\sigma=0\ \text{ for every }\sigma\in\mathscr C .
\]
\end{proposition}

\begin{proof}
\emph{Step 1: reduction to a single codimension.}  Suppose $[\beta]=0$ and
let $l_{\max}$ be the largest codimension occurring in $\beta$, so that
$d_{\min}=n-l_{\max}$ is the smallest dimension of a supporting cycle.  Let
$\sigma$ be a $d_{\min}$-dimensional face of $\mathscr C$ and let $x$ be a
point of $\RelInt\sigma$ which lies on no other face of $\mathscr C$ of
dimension $\leq d_{\min}$; such $x$ exists because the union of the other
faces of dimension $\leq d_{\min}$ meets $\sigma$ in a closed subset of
empty interior in $\sigma$.  Choose a compactly supported
$\omega$ of complementary bidegree with support in a ball $B$ around $x$ so
small that $B$ meets only those faces of $\mathscr C$ that contain $\sigma$.
By \eqref{eq:ch18-associated-current},
\[
 0=\langle[\beta],\omega\rangle
 =\int_\sigma A_\sigma\wedge\omega
 +\sum_{\dim\tau>d_{\min}}\int_{\tau\cap B}A_\tau\wedge\omega .
\]
Shrink $B$ to $x$ after normalizing $\omega$ so that the first integral is
fixed: the terms with $\dim\tau>d_{\min}$ are integrals over sets of
dimension $>d_{\min}$ and are $O(r^{\dim\tau-d_{\min}})$ relative to the
first one, where $r$ is the radius of $B$; more precisely, take
$\omega=\chi_r\eta$ with $\eta$ a fixed form and $\chi_r$ a cutoff equal to
$1$ on $B_{r/2}$ and supported in $B_r$, and divide by
$\mu_\sigma(\sigma\cap B_r)\asymp r^{d_{\min}}$.  The quotient of the first
term tends to the value of $A_\sigma\wedge\eta$ at $x$, and the quotients of
the remaining terms tend to $0$.  Hence $A_\sigma\wedge\eta$ vanishes at
$x$, for every $\eta$; by the perfect pairing of complementary superform
degrees on the $d_{\min}$-dimensional space $N_{\sigma,\mathbf R}$
(\Cref{lem:ch09-complement-pairing}), $A_\sigma(x)=0$.  Since the admissible
$x$ are dense in $\sigma$ and $A_\sigma$ is continuous, $A_\sigma=0$ on
$\sigma$.

\emph{Step 2: descending induction.}  Remove from $\beta$ the codimension-%
$l_{\max}$ terms; by Step 1 they contribute nothing to $[\beta]$, so the
remaining preform still induces the zero current, and the argument of Step 1
applies to the next largest codimension.  After finitely many steps every
$A_\sigma$ vanishes.

\emph{Step 3: the converse.}  If every $A_\sigma$ vanishes then every
integrand in \eqref{eq:ch18-associated-current} vanishes after subdivision
to $\mathscr C$, and $[\beta]=0$ by
\Cref{lem:ch18-subdivision-invariance}.
\end{proof}

\begin{definition}[Delta-preforms on an open subset of $N_{\mathbf R}$]
\label{def:ch18-delta-form-euclidean}
We write $P^{p,q}(\Omega)$ for the vector space of generalized
delta-preforms of bidegree $(p,q)$ on $\Omega$, regarded as polyhedral
supercurrents.  Equivalently, one may start with presentations
\eqref{eq:delta-preform} and divide by those whose associated current
vanishes.  The letter $P$ is deliberate: the letter $B$ will be reserved
for the smaller differential algebra of global delta-forms constructed in
\Cref{sec:ch18-delta-analytic}.
\end{definition}

\begin{remark}[Which normalization is in force]
\label{rem:ch18-normalization-choice}
On a fixed open subset of \(N_{\mathbf R}\), a delta-preform is the
polyhedral supercurrent represented by an expression
\eqref{eq:delta-preform}, not the raw expression itself.  Equality of two
such preforms is therefore equality of their associated polyhedral
supercurrents; after a common subdivision,
\Cref{prop:ch18-current-faithful} identifies this with equality of every
restricted coefficient.  This is the fixed-ambient convention of
\cite[Definition~2.9 and Remark~2.11]{GK17}, and
\Cref{def:ch18-delta-form-euclidean} records it explicitly.  In particular,
the fixed-ambient space $P(\Omega)$ is an algebra, but it is not in general
preserved by the current differentials $\dprime$ and $\ddprime$; see
\Cref{rem:ch18-differentials-descend}.

The statement changes on an analytic variety.  Equality of the single
ordinary current on the base is not the functorial equivalence used to
transport through every change of tropical chart.  The global definition in
\Cref{sec:ch18-delta-analytic} therefore uses compatible tropical charts
after every algebraic morphism, exactly as in \cite[\S4]{GK17}; the corresponding
current-detection criterion is the all-algebraic-pullback criterion of
\cite[Proposition~6.8]{GK17}.  Fixed-chart faithfulness must not be used as
a substitute for that global criterion.
\end{remark}

\begin{proposition}
\label{prop:ch18-product-descends}
Let $\beta=\sum_i\alpha_i\wedge\delta_{C_i}$ and
$\gamma=\sum_j\beta_j\wedge\delta_{D_j}$ be delta-preforms and suppose
$[\beta]=0$.  Then the delta-preform
$\sum_{i,j}(\alpha_i\wedge\beta_j)\wedge\delta_{C_i\cdot D_j}$ produced by
\Cref{thm:delta-preform-algebra} also induces the zero current.
Consequently the product of \Cref{thm:delta-preform-algebra} is well defined
on $P^{\bullet,\bullet}(\Omega)$.
\end{proposition}

\begin{proof}
Decompose \(\beta\) and \(\gamma\) first by codimension.  It is enough by
bilinearity to treat pieces of fixed codimensions \(l\) and \(l'\).  Choose
a common polyhedral complex \(\mathscr C\) for every supporting cycle.
After collecting all terms supported on the same cell, write
\[
 \beta=\sum_{\sigma\in\mathscr C_{n-l}}
             A_\sigma\wedge\delta_\sigma,
 \qquad
 \gamma=\sum_{\sigma'\in\mathscr C_{n-l'}}
             B_{\sigma'}\wedge\delta_{\sigma'} .
\]
Here \(A_\sigma\) and \(B_{\sigma'}\) are precisely the weighted sums of
the restricted coefficient forms.  Since \([\beta]=0\),
\Cref{prop:ch18-current-faithful} gives
\(A_\sigma=0\) for every \(\sigma\).

Choose one sufficiently small generic displacement used simultaneously to
compute all stable intersections and refine the resulting intersections.
For every resulting cell \(\tau\), the coefficient of \(\delta_\tau\) in
\(\beta\wedge\gamma\) is
\[
 \sum_{\substack{\sigma,\sigma'\\
                  \tau=\sigma\cap\sigma'\\
                  \sigma\cap(\sigma'+\varepsilon v)\ne\varnothing}}
 [N:N_\sigma+N_{\sigma'}]\,
 \bigl(A_\sigma\wedge B_{\sigma'}\bigr)\big|_\tau .
\]
This is the coefficient formula in the proof of
\cite[Proposition~2.12(ii)]{GK17}; it is also
\eqref{eq:ch18-stable-weight} with the smooth coefficients retained.
Every summand vanishes because \(A_\sigma=0\).  Hence every cell coefficient
of the product is zero, and
\Cref{prop:ch18-current-faithful} gives
\([\beta\wedge\gamma]=0\).

Bilinearity of the product then makes it well defined on
$P^{\bullet,\bullet}(\Omega)$: if $\beta\equiv\beta'$ and
$\gamma\equiv\gamma'$, then
$\beta\wedge\gamma-\beta'\wedge\gamma'
=(\beta-\beta')\wedge\gamma+\beta'\wedge(\gamma-\gamma')$ induces the zero
current by the case just proved and its mirror image.
\end{proof}

\begin{warning}{The fixed-ambient algebra is not a differential algebra}
\label{rem:ch18-differentials-descend}
The currents $\dprime[\beta]$ and $\ddprime[\beta]$ are defined for every
$\beta\in P(\Omega)$, but they need not belong to $P(\Omega)$.  In
particular, differentiating the underlying current does \emph{not} define
operators $P\to P$.  This is the obstruction isolated in
\cite[Remark~2.17 and Proposition~2.18]{GK17}.  The global differential algebra
$B_{\mathrm{GK}}$ is obtained only after imposing functorial closedness and
then adjoining smooth coefficients; see
\Cref{prop:ch18-global-pb-structure}.  Confusing $P$ and $B_{\mathrm{GK}}$
would incorrectly allow arbitrary piecewise-smooth degree-zero functions to
act on delta-currents by the differential Leibniz rule.
\end{warning}

\subsection{Stable intersection as the polyhedral multiplication}
\label{sec:ch18-stable-intersection}

Tropical cycles possess a stable intersection product \(C\cdot D\).  After a
common subdivision and a small generic translation, it is computed by
intersecting faces and multiplying their weights by lattice indices.  The
balancing condition makes the limit independent of the translation.

We recall the formula, since it is the one place where a lattice index that
is not $1$ can appear and it is the source of every multiplicity in the
chapter.  If $C$ and $D$ are balanced cycles in $N_{\mathbf R}$ of
codimensions $l$ and $l'$, then by
\Cref{thm:tropical-stable-intersection} the stable intersection $C\cdot D$ is
the balanced cycle of codimension $l+l'$ whose weight on a maximal face
$\tau$ is
\begin{equation}
\label{eq:ch18-stable-weight}
 m_\tau(C\cdot D)(x)
 =\sum_{\substack{\sigma\in C,\ \sigma'\in D\\
                   \tau=\sigma\cap\sigma'\\
                   \sigma\cap(\sigma'+\varepsilon v)\ne\varnothing}}
 m_\sigma(x)m_{\sigma'}(x)\,
 \bigl[N:N_\sigma+N_{\sigma'}\bigr],
\end{equation}
where the cycles have first been put on a common complex,
$x\in\RelInt\tau$, $v$ is generic, and $\varepsilon>0$ is sufficiently
small.  Thus $\tau$ is the \emph{unshifted} intersection; the translated
intersection occurs only in the nonemptiness condition.  This is
\cite[(1.4.1)]{GK17}, in the notation of
\eqref{eq:tropical-stable-weight}.  Three
consequences are used repeatedly:
\begin{enumerate}
\item\label{itm:ch18-si:support}
 $|C\cdot D|\subseteq|C|\cap|D|$, so the coefficient forms of a product
 delta-preform need only be defined near $|C|\cap|D|$;
\item\label{itm:ch18-si:star}
 the weight \eqref{eq:ch18-stable-weight} depends only on the stars of $C$
 and $D$ along $\tau$ (\Cref{cor:ch10-stable-star}), so the product is
 local;
\item\label{itm:ch18-si:unimodular}
 if $N_\sigma+N_{\sigma'}=N$ for all contributing pairs, all indices are
 $1$ and \eqref{eq:ch18-stable-weight} is a plain count with multiplicities.
\end{enumerate}

\begin{example}[Two tropical lines, with the numbers]
\label{ex:ch18-two-lines-numbers}
Let $L$ be the standard tropical line of \eqref{eq:ch18-gauge-rays} and let
$L'=L+w$ be a generic translate.  Each of the three rays of $L$ meets each
of the three rays of $L'$ in at most one point, and for generic $w$ exactly
one pair meets.  The two tangent lattices are $\mathbf Zv_i$ and
$\mathbf Zv_j$ with $i\neq j$, and
\[
 [\mathbf Z^2:\mathbf Zv_1+\mathbf Zv_2]
 =\bigl|\det\begin{pmatrix}1&0\\0&1\end{pmatrix}\bigr|=1,
 \quad
 [\mathbf Z^2:\mathbf Zv_1+\mathbf Zv_3]
 =\bigl|\det\begin{pmatrix}1&-1\\0&-1\end{pmatrix}\bigr|=1,
\]
and similarly for the third pair.  So $L\cdot L'$ is one point of weight
$1$, and by the translation invariance in
\Cref{thm:tropical-stable-intersection}, $L\cdot L=1\cdot\{0\}$.  This is
the number used in \Cref{ex:ch18-corner-max}, and the general statement is
the tropical B\'ezout theorem.  Had we taken the doubled line $2L$, the
weight would have been $4$; had we taken the fan
$L_{(2)}$ with rays $(1,0),(0,1),(-1,-1)$ but weights $2,2,2$, the same
computation gives $L_{(2)}\cdot L_{(2)}=4\cdot\{0\}$.  Weights are not
normalizations; they are part of the cycle.
\end{example}

\subsection{The algebra of delta-preforms}
\label{sec:ch18-delta-algebra}

\begin{theorem}\label{thm:delta-preform-algebra}
There is a unique $C^\infty(\Omega)$-bilinear, local,
graded-commutative product extending wedge of superforms and stable tropical
intersection such that
\[
 (\alpha\wedge\delta_C)\wedge(\beta\wedge\delta_D)
       =(\alpha\wedge\beta)\wedge\delta_{C\cdot D}
\]
whenever $C\cdot D$ is computed by fan displacement.  A common refinement
followed by a generic displacement defines the product without a
proper-intersection hypothesis.  Integral affine pullback of delta-preforms
is defined and respects this product.
\end{theorem}
\begin{proof}
For two displayed presentations, define the product by bilinearity and by
the displayed rule.  Choose a common complex for all coefficients and
cycles.  The fan-displacement formula then gives, on every cell $\tau$, the
coefficient written in the proof of
\Cref{prop:ch18-product-descends}.  That formula depends only on the
collected coefficients of the two polyhedral currents, so
\Cref{prop:ch18-product-descends} proves independence of both
presentations.

Locality is \Cref{lem:ch18-product-local}.  Graded commutativity follows
from the Koszul rule for superforms and commutativity of stable
intersection.  For three factors, take one common subdivision and two
generic displacement vectors.  Associativity of stable intersection
(\Cref{thm:tropical-stable-intersection}) identifies the lattice-indexed
coefficient sums for the two bracketings, while associativity of the
coefficient wedge identifies their summands.  The constant superform $1$
on $N_{\mathbf R}$, paired with its unit-weight cycle, is the unit.  These
properties also prove uniqueness, since the elementary terms
$\alpha\wedge\delta_C$ span $P(\Omega)$.

If $F:N'_{\mathbf R}\to N_{\mathbf R}$ is integral affine, define
$F^*(\alpha\wedge\delta_C)=F^*\alpha\wedge\delta_{F^*C}$, where the cycle
pullback is defined by stable intersection with the graph of $F$, with the
lattice indices of \cite[\S1.4(v)]{GK17}.  Its common-refinement
coefficient formula depends only on the source current, so it descends to
$P$; the projection formula and associativity of stable intersection give
$F^*(\beta\wedge\gamma)=F^*\beta\wedge F^*\gamma$.  This is exactly the
construction of \cite[Proposition~2.12]{GK17}.  It is a pullback on the
polyhedral algebra $P$, not a pullback of arbitrary currents.
\end{proof}

Three points in the preceding proof deserve to be isolated, because they
are what one actually verifies when computing.

\begin{lemma}
\label{lem:ch18-differential-delta}
Let $C$ be a constant-weight balanced pure $d$-dimensional tropical cycle
in $\Omega$ and let $\alpha$ be a smooth $(p,q)$-form near $|C|$.  Then, as
currents,
\begin{equation}
\label{eq:ch18-leibniz-delta}
 \dprime\bigl[\alpha\wedge\delta_C\bigr]
 =\bigl[\dprime\alpha\wedge\delta_C\bigr],
 \qquad
 \ddprime\bigl[\alpha\wedge\delta_C\bigr]
 =\bigl[\ddprime\alpha\wedge\delta_C\bigr].
\end{equation}
\end{lemma}

\begin{proof}
The integration current $\delta_C$ is $\dprime$- and
$\ddprime$-closed by \Cref{thm:ch11-current-closed}.  Apply the current
Leibniz rule of \Cref{prop:current-differential-calculus} to the action of
the smooth form $\alpha$ on $\delta_C$.  The terms involving a differential
of $\delta_C$ vanish and give the two displayed identities.
\end{proof}

\begin{warning}{Smooth weights require the full polyhedral derivative}
The hypothesis ``constant-weight balanced'' in
\Cref{lem:ch18-differential-delta} cannot be dropped.  For smooth weights,
the balancing condition still cancels the codimension-one boundary terms,
but the facewise differentials of the weight functions remain.  Equivalently,
if $C=(\mathscr C,m)$, then \cite[Proposition~2.16]{GK17} gives
\[
 \dprime\delta_C=\dprime m\wedge\delta_{\mathscr C}
   :=\sum_{\sigma\in\mathscr C_d}
        \dprime m_\sigma\wedge\delta_\sigma,
 \qquad
 \ddprime\delta_C=\ddprime m\wedge\delta_{\mathscr C},
\]
and the current Leibniz rule supplies the corresponding formula for
$\alpha\wedge\delta_C$.  These polyhedral derivatives need not lie in
$P(\Omega)$; this is precisely why
\Cref{rem:ch18-differentials-descend} is necessary.
\end{warning}

\begin{lemma}
\label{lem:ch18-product-local}
Let $\beta,\gamma$ be delta-preforms on $\Omega$ and $\Omega'\subseteq\Omega$
open.  Then $(\beta\wedge\gamma)|_{\Omega'}=\beta|_{\Omega'}\wedge
\gamma|_{\Omega'}$.  Moreover, if $\beta$ vanishes on $\Omega'$ as a
delta-preform, so does $\beta\wedge\gamma$.
\end{lemma}

\begin{proof}
Restriction of a delta-preform means restriction of every coefficient form
and intersection of every supporting cycle with $\Omega'$; this is
compatible with subdivision.  By
\cref{itm:ch18-si:star} the stable-intersection weight
\eqref{eq:ch18-stable-weight} at a face $\tau$ depends only on the stars of
the two cycles along $\tau$, hence only on their germs near any point of
$\RelInt\tau$.  Both assertions follow, the second from
\Cref{prop:ch18-product-descends} applied on $\Omega'$.
\end{proof}

\begin{remark}[Why uniqueness needs the two normalizations]
\label{rem:ch18-uniqueness-normalizations}
The word ``unique'' in \Cref{thm:delta-preform-algebra} refers to a product
that (i) restricts to the wedge product on codimension-zero terms, and (ii)
restricts to stable intersection on the supporting cycles when the
coefficients are $1$.  Without (ii), another symmetric unital factor
$\lambda(l,l')$ can preserve associativity when it satisfies
$\lambda(l,l')\lambda(l+l',l'')=
 \lambda(l',l'')\lambda(l,l'+l'')$; for example,
$\lambda(l,l')=c^{ll'}$.  Requirement (ii) rules out every such twist and is a genuine
normalization, and it is the one that makes
\Cref{thm:gk-delta-cl-measure} come out with the intersection numbers of
\Cref{thm:cl-model-properties} rather than with a multiple of them.  The
reader who suspects a factor should test on \Cref{ex:ch18-two-lines-numbers},
where the answer must be $1$, and on the doubled line, where it must be $4$.
\end{remark}

\subsection{The corner-locus formula}
\label{sec:ch18-corner-formula}

The following proposition is the technical heart of the whole delta-calculus.
Everything later---Poincar\'e--Lelong, Chern delta-forms, the positivity
criterion, and the comparison with model functions---is an application of it
in a particular tropical chart.

\begin{proposition}\label{prop:tropical-delta-pl}
Let \(C\) be a tropical cycle and \(\varphi\) a piecewise smooth function on
\(|C|\).  Then
\begin{equation}
 \dprimeddprime[\varphi]
    =[\dprime_P \ddprime_P\varphi]+\delta_{\varphi\cdot C}.
\label{eq:tropical-delta-pl}
\end{equation}
The first term is the facewise Hessian and the second is the weighted corner
locus.  For piecewise affine \(\varphi\), only the second term remains.
\end{proposition}
\begin{proof}
Choose a subdivision on which $\varphi$ is smooth on every facet and test
the current on a compactly supported form of complementary bidegree.  Apply
integration by parts twice on each facet.  The interior term is the
facewise Hessian
$[\dprime_P\ddprime_P\varphi]$.  Along an internal codimension-one face
$\tau$, the two boundary orientations are opposite, while the normal
derivatives of $\varphi$ need not agree.  Their weighted difference is
\[
       \sum_{\sigma\supset\tau}m_\sigma
          \partial_{u_{\sigma/\tau}}\varphi.
\]
This is exactly the weight of the corner cycle $\varphi\cdot C$ at $\tau$.
The tangential terms cancel because $\varphi$ is continuous, and the terms
coming from an affine change of local representative cancel by balancing of
$C$.  Summing over all faces proves \eqref{eq:tropical-delta-pl}.

If $\varphi$ is piecewise affine, its facewise Hessian is zero.  In dimension
one, the remaining weights are the outgoing-slope sums in
\eqref{eq:delta-graph-laplacian}.  This is the cellwise Stokes proof of
\cite[Theorem~3.16 and Corollary~3.19]{GK17}.
\end{proof}

We write out the definition of the corner cycle that the proof produces,
since it is used as a formula from here on.

\begin{definition}[Corner cycle]
\label{def:ch18-corner-cycle}
\index{corner locus}%
Let $C$ be a balanced weighted pure $d$-dimensional cycle and let $\varphi$
be continuous on $|C|$ and smooth on each maximal cell of some subdivision
$\mathscr C$ of $C$.  Fix a codimension-one face $\tau$ and a point
$x\in\RelInt\tau$.  For each maximal $\sigma\supset\tau$ let
$u_{\sigma/\tau}\in N_\sigma$ be a lift of the primitive normal of
\Cref{def:ch10-star}, pointing from $\tau$ into $\sigma$, and let
$\varphi_\sigma$ be the smooth extension of $\varphi|_\sigma$ near $\tau$.
Balancing of $C$ gives
\begin{equation}
\label{eq:ch18-corner-defect}
 w_\tau:=\sum_{\sigma\in\mathscr C_d,\ \sigma\supset\tau}
   m_\sigma\,u_{\sigma/\tau}\ \in\ N_\tau .
\end{equation}
The \emph{corner cycle} $\varphi\cdot C$ is the weighted
$(d-1)$-dimensional cycle whose support is contained in the codimension-one
skeleton of $\mathscr C$ and whose weight at $\tau$ is
\begin{equation}
\label{eq:ch18-corner-weight}
 m_\tau(\varphi\cdot C)(x)
 =\sum_{\sigma\in\mathscr C_d,\ \sigma\supset\tau}
   m_\sigma\,\partial_{u_{\sigma/\tau}}\varphi_\sigma(x)
   \ -\ \partial_{w_\tau}\bigl(\varphi|_\tau\bigr)(x) .
\end{equation}
\end{definition}

\begin{remark}[The correction term]
\label{rem:ch18-corner-correction}
The second term of \eqref{eq:ch18-corner-weight} is not decoration.  The
lift $u_{\sigma/\tau}$ is only canonical modulo $N_\tau$, so the first sum
alone depends on choices; the correction cancels exactly that dependence,
as \Cref{lem:ch18-corner-welldefined}\,\cref{itm:ch18-cc:affine} shows.  In
the case used most often in this book---$C$ one-dimensional and $\tau$ a
vertex---one has $N_\tau=0$, hence $w_\tau=0$ and no correction, and
\eqref{eq:ch18-corner-weight} reduces to the atom \eqref{eq:ch18-atom}.  The
correction also vanishes whenever $\varphi$ is affine along $\tau$ in the
direction $w_\tau$, in particular whenever $\varphi|_\tau$ is constant.
\end{remark}

Three checks make \eqref{eq:ch18-corner-weight} usable.

\begin{lemma}
\label{lem:ch18-corner-welldefined}
In the situation of \Cref{def:ch18-corner-cycle}:
\begin{enumerate}
\item\label{itm:ch18-cc:affine}
 $m_\tau(\varphi\cdot C)$ is independent of the choice of the lifts
 $u_{\sigma/\tau}$, and is unchanged if $\varphi$ is replaced by
 $\varphi+\ell$ with $\ell$ the restriction of an affine function on
 $N_{\mathbf R}$;
\item\label{itm:ch18-cc:subdiv}
 $\varphi\cdot C$ does not depend on the subdivision $\mathscr C$;
\item\label{itm:ch18-cc:balanced}
 $\varphi\cdot C$ is balanced.
\end{enumerate}
\end{lemma}

\begin{proof}
\cref{itm:ch18-cc:affine}  \emph{Independence of the lifts.}  Replace
$u_{\sigma/\tau}$ by $u_{\sigma/\tau}+z_\sigma$ with $z_\sigma\in N_\tau$.
Since $\varphi$ is continuous, the restriction of $\varphi_\sigma$ to $\tau$
is $\varphi|_\tau$ for every $\sigma$, so
$\partial_{z_\sigma}\varphi_\sigma(x)=\partial_{z_\sigma}(\varphi|_\tau)(x)$.
Hence the first sum in \eqref{eq:ch18-corner-weight} changes by
$\sum_\sigma m_\sigma\partial_{z_\sigma}(\varphi|_\tau)(x)
 =\partial_{z}(\varphi|_\tau)(x)$ with
$z=\sum_\sigma m_\sigma z_\sigma\in N_\tau$, while $w_\tau$ changes to
$w_\tau+z$, so the second term changes by exactly the same amount.  The
difference is unchanged.

\emph{Affine invariance.}  Adding $\ell$ with linear part $\rd\ell$ changes
the first sum by
$\sum_\sigma m_\sigma\langle \rd\ell,u_{\sigma/\tau}\rangle
 =\langle \rd\ell,w_\tau\rangle$, by \eqref{eq:ch18-corner-defect}, and
changes the second term by $\partial_{w_\tau}\ell=\langle \rd\ell,w_\tau
\rangle$ as well, because $\ell|_\tau$ has the same linear part along
$N_\tau$.  The two changes cancel.  Balancing
(\Cref{def:ch10-weighted-balanced}) is used exactly once, to know that
$w_\tau\in N_\tau$, so that $\partial_{w_\tau}(\varphi|_\tau)$ makes sense.

\cref{itm:ch18-cc:subdiv}  A finer subdivision introduces new
codimension-one faces $\tau'$ lying in the relative interior of some
maximal $\sigma$.  There exactly two cells of the finer complex meet along
$\tau'$, with the same weight $m_\sigma$ and opposite normals $\pm u$; hence
$w_{\tau'}=m_\sigma(u-u)=0$ and the correction term vanishes, while
$\varphi$ is smooth across $\tau'$, so the two normal derivatives are
negatives of each other and the first sum in
\eqref{eq:ch18-corner-weight} is $0$.  At a face $\tau$ already present, the
maximal cells through $\tau$ are subdivided without changing weights or
normals, so both terms of \eqref{eq:ch18-corner-weight} are unchanged.

\cref{itm:ch18-cc:balanced}  This is the corner-locus balancing theorem
\cite[Proposition~1.12]{GK17}, applied to the constant-weight tropical cycle
$C$ and the piecewise smooth function $\varphi$.  Its conclusion is a
tropical cycle with smooth weights.  Notice that the current identity in
\Cref{prop:tropical-delta-pl} does not by itself give a shortcut to this
conclusion: the facewise Hessian has smooth coefficients, and its
integration current need not be separately $\dprime$-closed because the
derivatives of those coefficients contribute as in
\Cref{rem:ch18-two-weights}.
\end{proof}

\subsection{The corner-locus formula in examples}
\label{sec:ch18-corner-examples}

\begin{example}[Dimension one, all cases]
\label{ex:ch18-corner-dim-one}
Let $C\subset\mathbf R$ be the line with weight $1$ and let
$\varphi=\max\{0,t\}$.  The subdivision at $t=0$ has two maximal cells with
$\varphi$ affine on each, so $\dprime_P\ddprime_P\varphi=0$, and
\eqref{eq:ch18-corner-weight} at $\tau=\{0\}$ gives
$1\cdot0+1\cdot1=1$ (the outgoing slope is $0$ to the left, $1$ to the
right).  Thus $\varphi\cdot C=\{0\}$ with weight $1$ and
$\dprimeddprime[\varphi]=\delta_0$, as computed already in
\Cref{sec:ch16-obstruction}.

For $\varphi=\min\{0,t\}=-\max\{0,-t\}$ the same computation gives weight
$-1$, so $\varphi\cdot C=-\{0\}$ and $\dprimeddprime[\varphi]=-\delta_0$.
The corner cycle is therefore \emph{not} effective in general, and its
effectivity is one of the two positivity conditions isolated in
\Cref{sec:delta-positivity}.

For $\varphi(t)=t^2$, smooth, the corner cycle is $0$ and the facewise
Hessian is $2\,\dprime t\wedge\ddprime t$, so
$\dprimeddprime[\varphi]=2\,\dprime t\wedge\ddprime t$, an absolutely
continuous form; consistently, $\varphi$ has no atom.
\end{example}

\begin{example}[The gauge on the tropical line, as a corner cycle]
\label{ex:ch18-gauge-corner}
Return to \Cref{ex:ch18-gauge-atom}: $C=L$ is the standard tropical line
with unit weights and $\varphi=g|_L$ has slope $1$ on each ray.  The only
codimension-one face is $\tau=\{0\}$, and \eqref{eq:ch18-corner-weight}
gives
\[
 m_0(\varphi\cdot L)=1\cdot1+1\cdot1+1\cdot1=3 ,
\]
so $\varphi\cdot L=3\{0\}$ and $\dprimeddprime[\varphi]=3\delta_0$,
reproving \eqref{eq:ch18-gauge-atom}.  Note that $\varphi\cdot L$ is
effective, as it must be, since $\varphi$ is psh; and that
$m_0(\varphi\cdot L)$ is unchanged if $\varphi$ is modified by an affine
function, in accordance with
\Cref{lem:ch18-corner-welldefined}\,\cref{itm:ch18-cc:affine}: adding
$\ell(z)=z_1$ changes the three slopes to $1+1=2$, $1+0=1$, $1-1=0$, whose
sum is again $3$.
\end{example}

\begin{example}[A two-dimensional corner cycle and a product]
\label{ex:ch18-two-dim-product}
On $C=\mathbf R^2$ with weight $1$, let $\varphi_1=\max\{0,x_1\}$ and
$\varphi_2=\max\{0,x_2\}$.  Then $\varphi_1\cdot C$ is the line
$\{x_1=0\}$ with weight $1$, and $\varphi_2\cdot C$ is $\{x_2=0\}$ with
weight $1$; these meet transversally at the origin with lattice index
$[\mathbf Z^2:\mathbf Ze_2+\mathbf Ze_1]=1$, so by
\eqref{eq:ch18-stable-weight}
\[
 (\varphi_1\cdot C)\cdot(\varphi_2\cdot C)=1\cdot\{0\},
\]
and the delta-form product gives
$\dprimeddprime\varphi_1\wedge\dprimeddprime\varphi_2=\delta_0$.  This is
the second row of the backward dictionary of \Cref{sec:ch16-dict-backward},
computed there by the Bedford--Taylor recursion; the two answers agree, both
being the measure $\delta_0$ of mass $1$.

For $\varphi=\max\{0,x_1,x_2\}$ of \Cref{ex:ch18-corner-max} one has
$\varphi\cdot C=C_\varphi$, the standard tropical line, and
$(\varphi\cdot C)^{\cdot2}=1\cdot\{0\}$ by
\Cref{ex:ch18-two-lines-numbers}; so
$(\dprimeddprime\varphi)^{\wedge2}=\delta_0$, of mass $1$, matching the toric
computation \eqref{eq:ch16-fan-p2-mass} for $\mathbf P^2$.
\end{example}

\begin{example}[A nonunimodular index]
\label{ex:ch18-nonunimodular}
On $\mathbf R^2$ let $\varphi_1=\max\{0,2x_1\}$ and
$\varphi_2=\max\{0,x_2\}$.  Then $\varphi_1\cdot C$ is $\{x_1=0\}$ with
weight $2$, because the jump of the normal derivative is $2$, and
$\varphi_2\cdot C$ is $\{x_2=0\}$ with weight $1$.  The product is
$2\cdot\{0\}$.  If instead $\varphi_1=\max\{0,x_1+x_2\}$ and
$\varphi_2=\max\{0,x_1-x_2\}$, the two corner lines have primitive tangent
vectors $(1,-1)$ and $(1,1)$, weights $1$, and
\[
 [\mathbf Z^2:\mathbf Z(1,-1)+\mathbf Z(1,1)]
 =\bigl|\det\begin{pmatrix}1&1\\-1&1\end{pmatrix}\bigr|=2 ,
\]
so the product is $2\cdot\{0\}$ again, although both corner cycles have
weight $1$.  This is the phenomenon of
\Cref{ex:ch10-lattice-index-test}: a lattice index of $2$ appears where a
na\"\i ve count would give $1$.  Any calculation of a Chambert--Loir measure
that omits such an index will produce a wrong total mass.
\end{example}

\section{Delta-forms on analytifications}
\label{sec:ch18-delta-analytic}

\subsection{Very affine charts}
\label{sec:ch18-very-affine}

Let \(X\) be an algebraic variety over an algebraically closed, complete,
nontrivially valued field \(K\).  If \(U\subset X\) is very affine, its
canonical embedding in a split torus gives a tropicalization
\[
                      \operatorname{trop}_U:U^{\mathrm{an}}\to N_{U,\mathbf R}.
\]
A tropical chart is an inverse image \(V=\operatorname{trop}_U^{-1}(\Omega)\).

We spell out the three properties of this chart system which are used below;
each is a result of \Cref{chap:analytic-tropicalization} transcribed into the
very affine setting.

\begin{lemma}
\label{lem:ch18-very-affine-charts}
Let $X$ be an algebraic variety over $K$ as above.
\begin{enumerate}
\item\label{itm:ch18-va:cover}
 The very affine open subsets of $X$ form a basis of the Zariski topology,
 and $U\mapsto\Trop(U):=\operatorname{trop}_U(U^{\An})$ is a finite
 $\Gamma$-rational polyhedral subset of $N_{U,\mathbf R}$ of pure dimension
 $\dim U$, canonically calibrated by
 \Cref{def:trop-canonical-calibration} and balanced by
 \Cref{thm:analytic-tropical-balancing}.
\item\label{itm:ch18-va:refine}
 If $g:X'\to X$ is an algebraic morphism and
 $(V',\varphi_{U'})$, $(V,\varphi_U)$ are compatible tropical charts, the
 induced integral affine map
 $F:N_{U',\mathbf R}\to N_{U,\mathbf R}$ satisfies
 $\operatorname{trop}_U\circ g^{\An}=F\circ\operatorname{trop}_{U'}$ and
 \[
       F_*\Trop(U')=\deg(g|_{U'})\,\Trop(g(U')).
 \]
 Here the degree is the function-field degree onto the schematic image,
 and is zero when that extension is not finite.  In particular, a dense
 very-affine open immersion into an irreducible $U$ has degree one and
 schematic image $U$.
\item\label{itm:ch18-va:filtered}
 Any two very affine opens containing a given point admit a common very
 affine refinement, and the resulting category of tropical charts at a point
 is filtered, by \Cref{lem:ch13-charts-filtered}.
\end{enumerate}
\end{lemma}

\begin{proof}
\cref{itm:ch18-va:cover} A very affine variety is by definition a closed
subvariety of a split torus $T=\Spec K[M]$; such opens form a basis because
any affine open is covered by principal opens of a closed embedding into
affine space, and a principal open of an affine variety embedded in
$\mathbf A^N$ embeds in $\mathbf G_m^{N+1}$.  Polyhedrality and pure
dimensionality are \Cref{thm:analytic-tropical-polyhedrality} and
\Cref{prop:analytic-tropical-dimension}; the calibration is
\Cref{def:trop-canonical-calibration}; balancing is
\Cref{thm:analytic-tropical-balancing}, whose hypotheses hold because
$U^{\An}$ has no boundary in the torus.

\cref{itm:ch18-va:refine} Compatibility gives the homomorphism of tori and
the integral affine map $F$; functoriality of tropicalization is
\Cref{lem:trop-moment-basic}\,\cref{itm:trop-moment:funct}.  The equality of
weighted tropical cycles is the Sturmfels--Tevelev multiplicity formula
\cite[(4.3.1)]{GK17}, equivalently
\eqref{eq:analytic-tropical-pushforward}.  For a dense open immersion the
two function fields agree, so the last assertion follows.

\cref{itm:ch18-va:filtered} The intersection of two very affine opens is very
affine, and \Cref{lem:ch13-charts-filtered} gives the filteredness statement
in the form used later.
\end{proof}

\begin{definition}[Generalized delta-forms and delta-forms]
\label{def:ch18-global-pb}
Fix a tropical chart $(V,\varphi_U)$, put
$\Omega=\operatorname{trop}_U(V)$, and choose an open
$\widetilde\Omega\subset N_{U,\mathbf R}$ with
$\Omega=\widetilde\Omega\cap\Trop(U)$.  If
$f:X'\to X$ is an algebraic morphism and $(V',\varphi_{U'})$ is a tropical
chart compatible with $(V,\varphi_U)$, let
$F:N_{U',\mathbf R}\to N_{U,\mathbf R}$ be the induced integral affine map
and put $\Omega'=\operatorname{trop}_{U'}(V')$.  For
$\widetilde\beta\in P(\widetilde\Omega)$ set
\begin{equation}
\label{eq:ch18-all-pullback-restriction}
 \mathcal R_{f,V'}(\widetilde\beta)
 :=\bigl(F^*\widetilde\beta\bigr)\wedge
      \delta_{\Trop(U')}\big|_{\Omega'}\in D(\Omega').
\end{equation}
Here one first restricts $F^*\widetilde\beta$ to a sufficiently small open
neighbourhood of $\Omega'$.

Let $\mathcal N(V,\varphi_U)$ be the ideal of all
$\widetilde\beta$ for which
$\mathcal R_{f,V'}(\widetilde\beta)=0$ for \emph{every} algebraic $f$ and
every compatible chart.  Define
\[
 P(V,\varphi_U):=P(\widetilde\Omega)/\mathcal N(V,\varphi_U).
\]
An element of this algebra is a chartwise generalized delta-form.  Let
$Z(V,\varphi_U)$ be the subalgebra of those classes $\omega$ for which every
current $\mathcal R_{f,V'}(\omega)$ is both $\dprime$- and
$\ddprime$-closed.  Finally let $AZ(V,\varphi_U)$ be the subalgebra generated
by the smooth forms and $Z(V,\varphi_U)$.  On a presentation
\begin{equation}
\label{eq:ch18-az-presentation}
                 \beta=\sum_i\alpha_i\wedge\omega_i,
 \qquad \alpha_i\in A(\Omega),\quad
        \omega_i\in Z(V,\varphi_U),
\end{equation}
put
\[
 \dprime\beta=\sum_i\dprime\alpha_i\wedge\omega_i,
 \qquad
 \ddprime\beta=\sum_i\ddprime\alpha_i\wedge\omega_i.
\]

Taking the directed colimits over tropical charts, followed by
sheafification, gives the sheaf $P$ of \emph{generalized delta-forms} and
its subsheaf $B_{\mathrm{GK}}$ of \emph{delta-forms}.  We suppress the
subscript when no confusion with another $B$ is possible.
\end{definition}

The universal quantifier in
\eqref{eq:ch18-all-pullback-restriction} is essential.  Testing only
refinements on $X$ does not recover the functorial quotient of
\cite[Definition~4.4]{GK17}; one must test compatible tropical charts after
every algebraic pullback.

\begin{proposition}
\label{prop:ch18-global-pb-structure}
The sheaves $P$ and $B_{\mathrm{GK}}$ are sheaves of bigraded,
graded-commutative algebras, and
\[
                         A\lhook\joinrel\longrightarrow
                         B_{\mathrm{GK}}\lhook\joinrel\longrightarrow P.
\]
The operators $\dprime$ and $\ddprime$ in
\Cref{def:ch18-global-pb} are well defined on $B_{\mathrm{GK}}$, satisfy the
two differential-algebra identities, and commute with algebraic pullback.
No differential on all of $P$ is asserted.
\end{proposition}

\begin{proof}
Integral affine pullback and the product on $P(\widetilde\Omega)$ commute by
\Cref{thm:delta-preform-algebra}.  Hence if every current
$\mathcal R_{f,V'}(\widetilde\beta)$ vanishes, the same is true after
multiplying $\widetilde\beta$ by any delta-preform; thus
$\mathcal N(V,\varphi_U)$ is an ideal and the product descends.  The same
formula defines pullback between compatible charts and shows its
functoriality.

Suppose $\omega$ and $\eta$ lie in $Z(V,\varphi_U)$.  After any algebraic
pullback and on any compatible chart, take a common polyhedral complex for
their restricted currents.  The fan-displacement formula expresses the
coefficient of their product on a cell as a lattice-indexed sum of wedges of
the two coefficient forms.  Applying $\dprime$ splits that sum into the two
Leibniz terms.  The closedness relations for $\omega$ and $\eta$, reindexed
using associativity of stable intersection, make both sums vanish.  The
same argument applies to $\ddprime$.  This is the cellwise proof of
\cite[Lemma~3.7 and Proposition~4.7]{GK17}, and proves that $Z$ is a
subalgebra.

Now let a sum \eqref{eq:ch18-az-presentation} represent zero.  After every
algebraic pullback its associated current is zero.  Differentiate that
current.  Since each pulled-back $\omega_i$ is closed, the current Leibniz
rule gives exactly the pullback of
$\sum_i\dprime\alpha_i\wedge\omega_i$, which must therefore vanish.  The
all-pullback quotient shows that this sum is zero.  This proves that
$\dprime$ is independent of the presentation; the proof for $\ddprime$ is
identical.  The relations $\dprime^2=\ddprime^2=0$,
$\dprime\ddprime+\ddprime\dprime=0$, and both Leibniz rules reduce to the
corresponding identities for the smooth coefficients.  Pullback
compatibility follows from the same formulas.

Finally, compatible charts form a directed system.  Restriction gives the
presheaf colimits, and sheafification gives $P$ and $B_{\mathrm{GK}}$.
Chartwise partitions of unity glue finite compatible families.  The first
inclusion is injective by fixed-chart faithfulness of smooth superforms; the
second is injective by construction.  This proves all assertions and is
the construction of \cite[Definitions~4.9, 4.13 and
Proposition~4.15]{GK17}.
\end{proof}

\begin{remark}[Reconciling the two quotients]
\label{rem:ch18-two-quotients-analytic}
On a fixed open subset of the ambient vector space,
\Cref{prop:ch18-current-faithful} says that the polyhedral current determines
its collected cell coefficients.  On an analytic chart, however, one first
restricts to $\Trop(U)$, and a class in $P(V,\varphi_U)$ is required to have
well-defined behaviour after every algebraic pullback.  Equality of the two
restricted currents on $X^{\An}$ tests only the identity morphism and is not
the definition of equality in $P$.  The exact faithful replacement is the
all-pullback criterion in
\Cref{prop:ch18-delta-vs-cld-currents}.  This statement concerns generalized
delta-forms; it neither asserts nor requires injectivity of the restriction
map from delta-currents to CLD currents.
\end{remark}

\subsection{Integration, Stokes, and delta-currents}
\label{sec:ch18-delta-currents}

Integration of compactly supported top-degree generalized delta-forms is
defined on one very-affine chart and is independent of that choice.  It
satisfies Stokes' formula on the differential subalgebra $B_{\mathrm{GK}}$.
Dualizing compactly supported delta-forms gives delta-currents, and every
generalized delta-form determines a delta-current.

We state the three assertions of the previous paragraph precisely, since the
whole comparison with \Cref{chap:cld-currents} rests on them.

\begin{proposition}
\label{prop:ch18-delta-integration}
Let $X$ be as above, of pure dimension $n$, and let $\beta\in P^{n,n}_c(X^{\An})$.
\begin{enumerate}
\item\label{itm:ch18-di:defined}
 A very affine chart of integration exists: this means a nonempty very
 affine $U\subset X$ for which
 $\beta|_{U^{\An}}=\operatorname{trop}_U^*\beta_U$ with
 $\beta_U\in P_c^{n,n}(U^{\An},\varphi_U)$.  The number
 \[
       \int_{X^{\An}}\beta:=\int_{|\Trop(U)|}\beta_U
 \]
 is independent of the chart of integration and of the polyhedral
 representative.
\item\label{itm:ch18-di:measure}
 $\beta$ determines a unique signed Radon measure $\mu_\beta$ characterized
 by
 $\int u\,\mathrm d\mu_\beta=\int u\beta$ for smooth compactly supported
 $u$.  The left side then defines the integral for arbitrary continuous
 compactly supported $u$.  If $\beta$
 is a positive generalized delta-form in the sense of
 \Cref{sec:delta-positivity}, then this measure is positive.  Positivity of
 this single measure is not asserted to imply functorial positivity of
 $\beta$.
\item\label{itm:ch18-di:stokes}
 If $X^{\An}$ is boundaryless and $\gamma$ is a compactly supported
 delta-form of bidegree $(n-1,n)$, then
 $\int_{X^{\An}}\dprime\gamma=0$, and likewise for $\ddprime$.
\end{enumerate}
\end{proposition}

\begin{proof}
\cref{itm:ch18-di:defined} Cover $\Supp\beta$ by finitely many tropical
charts on which $\beta$ is represented.  The intersection of their
very-affine algebraic opens contains a nonempty very-affine open $U$.
The gluing argument for the all-pullback quotient gives a unique
$\beta_U\in P^{n,n}(U^{\An},\varphi_U)$ representing $\beta$ on $U^{\An}$.
Because the bidegree is top degree, the support-removal statement
\cite[Corollary~5.6]{GK17} shows that $\Supp\beta\subset U^{\An}$,
and $\beta_U$ has compact support in $\Trop(U)$.  This is the construction
of a chart of integration in \cite[Proposition~5.7 and
Definition~5.8]{GK17}.

Choose a representative $\sum_i\alpha_i\wedge\delta_{C_i}$ of $\beta_U$.
The number
$\sum_i\sum_\sigma\int_\sigma m_{i,\sigma}\alpha_i$ is unchanged under
subdivision by \Cref{lem:ch18-subdivision-invariance}, and unchanged by a
change of representative inducing the same current, by
\Cref{prop:ch18-current-faithful}.  For a nonempty very-affine
$U'\subseteq U$, the representative on $\Trop(U')$ is the pullback along
the induced surjection $F$ of
\Cref{lem:ch18-very-affine-charts}\,\cref{itm:ch18-va:refine}, and
$\int_{F^{-1}(\sigma)\cap\Trop(U')}F^*\alpha$ equals
$\int_\sigma\alpha$ multiplied by the pushforward weight, by
\Cref{cor:ch11-projection-formula}; the degree-one multiplicity formula for
the open immersion makes the sums equal.  Two charts of integration are
compared after passing to a nonempty very-affine open in their intersection.
This proves independence and is \cite[Proposition~5.9]{GK17}.

\cref{itm:ch18-di:measure} Initially take $u\in A_c^{0,0}(X^{\An})$.
By linearity it is enough to treat one fixed-codimension-$l$ summand of
$\beta$; the constants for the finitely many codimensions are then added.
A common chart of integration for $\beta$ and $u\beta$ exists by the first
part.  Fix the preceding chart $U$ and a polyhedral-current presentation
\[
                    \widetilde\beta_U=\sum_\sigma
                       \alpha_\sigma\wedge\delta_\sigma .
\]
After refining so that $\Trop(U)=(\mathscr C_{\leq n},m)$, fan displacement
gives a finite constant
\[
 C_\beta=
 \sum_{\substack{\Delta\in\mathscr C_n,\ \sigma\in\mathscr C,\
                  \operatorname{codim}_{N_{U,\mathbf R}}\sigma=l\\
                  L_\Delta+L_\sigma=N_{U,\mathbf R},\
                  \tau:=\Delta\cap\sigma,\ \dim\tau=n-l}}
 [N_U:N_\Delta+N_\sigma]m_\Delta
 \int_\tau|\alpha_\sigma|_\tau.
\]
On a common tropical refinement on which $u$ is a smooth coefficient, the
projection formula for lattices and the Sturmfels--Tevelev multiplicity
formula cancel exactly the extra indices and sheets.  Termwise estimation
therefore gives
\begin{equation}
\label{eq:ch18-measure-sup-bound}
               \left|\int_{X^{\An}}u\beta\right|
                    \leq C_\beta\|u\|_\infty .
\end{equation}
The lattice cancellation is the identity proved in
\cite[Lemma~6.14]{GK17}; substituting it in the fan-displacement sum gives
\cite[Proposition~6.13]{GK17}.  Smooth compactly supported functions are
sup-norm dense in continuous compactly supported functions, so
\eqref{eq:ch18-measure-sup-bound} extends the functional uniquely to
$C_c(X^{\An})$.  The Riesz representation theorem now gives the unique
signed Radon measure $\mu_\beta$.  Only after this extension do we define
$\int u\beta:=\int u\,\mathrm d\mu_\beta$ for nonsmooth continuous $u$.

A positive generalized
delta-form has positive restriction for the identity morphism, so its
top-degree cell coefficients and hence its measure are positive.  The
all-pullback requirement is stronger than this one necessary test, as
\Cref{lem:ch18-positivity-consistent}\,\cref{itm:ch18-pos:global} explains.

\cref{itm:ch18-di:stokes} By
\Cref{prop:ch18-global-pb-structure}, the current associated to
$\dprime\gamma$ is the current derivative of the current associated to
$\gamma$.  Choose a compactly supported smooth function $\chi$ equal to $1$
on a neighbourhood of $\Supp\gamma$.  With the current sign convention of
\eqref{eq:ch14-dprime-sign},
\[
 \int_{X^{\An}}\dprime\gamma
 =\langle\dprime[\gamma]_D,\chi\rangle
 =\pm\langle[\gamma]_D,\dprime\chi\rangle=0,
\]
because $\dprime\chi$ vanishes near $\Supp\gamma$.  The same argument gives
the $\ddprime$ identity.  This proof also shows why the assertion is made
for $B_{\mathrm{GK}}$, not for all of $P$.
\end{proof}

We first make the continuity condition precise.  Fix a tropical chart
$(V,\varphi_U)$ with $\Omega=\operatorname{trop}_U(V)$ and a finite tuple
$\omega_J=(\omega_j)_{j\in J}$ in $Z(V,\varphi_U)$.  Let
$AZ(V,\varphi_U;\omega_J)$ be the subspace of elements admitting a
presentation $\sum_{j\in J}\alpha_j\wedge\omega_j$ with
$\alpha_j\in A(\Omega)$.  For a compact $Q\subset\Omega$ and $s\geq0$, let
$p_{Q,s}(\alpha)$ be the maximum on $Q$ of the coefficient derivatives of
$\alpha$ of order at most $s$, and set
\begin{equation}
\label{eq:ch18-delta-test-seminorm}
 p_{Q,s,\omega_J}(\beta)
 :=\inf_{\beta=\sum_j\alpha_j\wedge\omega_j}
       \max_{j\in J}p_{Q,s}(\alpha_j).
\end{equation}
These quotient seminorms give $AZ(V,\varphi_U;\omega_J)$ a locally convex
topology.

Now fix a tropical-chart cover $(V_i,\varphi_{U_i})_{i\in I}$ of an open
$W\subseteq X^{\An}$ and a finite tuple $\omega_{J_i}$ in
$Z(V_i,\varphi_{U_i})$ for each $i$.  Write
\[
 B(W;V_i,\varphi_{U_i},\omega_{J_i}:i\in I)
\]
for the delta-forms whose restriction to $V_i$ belongs to
$AZ(V_i,\varphi_{U_i};\omega_{J_i})$ for every $i$.  Give this space the
coarsest locally convex topology for which all restriction maps to the
chart spaces are continuous.  If $C\subset W$ is compact, append a subscript
$C$ to denote the subspace of forms supported in $C$, with its induced
topology.

\begin{definition}[Delta-current]
\label{def:ch18-delta-current}
\index{delta-current}%
A \emph{delta-current} on $W$ is a real linear functional on $B_c(W)$ whose
restriction to
\[
 B_C(W;V_i,\varphi_{U_i},\omega_{J_i}:i\in I)
\]
is continuous for every compact $C$, every tropical-chart cover, and every
choice of the finite tuples $\omega_{J_i}$.  It has bidegree $(p,q)$ when it
vanishes outside test bidegree $(n-p,n-q)$.  Its differentials
$\dprime,\ddprime$ and its support are defined by the formulas of
\Cref{chap:cld-currents}.
A symmetric delta-current is weakly positive, positive, or strongly
positive when it is nonnegative, respectively, on compactly supported
strongly positive, positive, or weakly positive delta-forms of
complementary bidegree.  These are the same dual-cone conventions as in
\Cref{sec:ch14-positivity}.  This is precisely
\cite[Definitions~6.1--6.4]{GK17} and
\cite[Definition~2.10]{GK19}.
\end{definition}

\begin{proposition}
\label{prop:ch18-delta-vs-cld-currents}
Let $X$ be as above.
\begin{enumerate}
\item\label{itm:ch18-dc:restrict}
 Restriction along the inclusion of smooth forms into delta-forms is a
 linear map
 $\iota^\ast$ from delta-currents to CLD currents, commuting with
 $\dprime$, $\ddprime$ and with multiplication by smooth functions, and
 not increasing supports.  It sends each of the three positivity cones of
 delta-currents into the corresponding cone of CLD currents.
\item\label{itm:ch18-dc:form}
 Every generalized delta-form $\beta\in P^{p,q}(X^{\An})$ determines a
 delta-current $[\beta]_E$ by $\gamma\mapsto\int\beta\wedge\gamma$, and
 $\iota^\ast[\beta]_E$ is the CLD current
 \eqref{eq:ch18-associated-current}.
\item\label{itm:ch18-dc:criterion}
 For generalized delta-forms \(\beta,\beta'\in P^{p,q}(W)\) on an open
 $W\subseteq X^{\An}$, the faithful current criterion is functorial:
 \(\beta=\beta'\) if and only if
 \[
  [f^*\beta|_{W'}]_{D}=[f^*\beta'|_{W'}]_{D}
 \]
 for every morphism \(f:X'\to X\) of algebraic \(K\)-varieties and every
 open \(W'\subseteq(X')^{\An}\) with $f(W')\subseteq W$.  No injectivity or
 noninjectivity assertion about $\iota^*:E\to D$ is part of this criterion.
\end{enumerate}
\end{proposition}

\begin{proof}
\cref{itm:ch18-dc:restrict} A smooth form is a delta-form of codimension
zero.  In the continuity test of \Cref{def:ch18-delta-current}, choose on
each chart the one-element tuple $(1)$.  Then
$AZ(V_i,\varphi_{U_i};1)=A(\Omega_i)$, and the seminorms
\eqref{eq:ch18-delta-test-seminorm} are the usual coefficient-derivative
seminorms defining the CLD test-form topology of
\Cref{sec:ch13-topology}.  Thus restriction of a delta-current is a
continuous CLD current.

On every fixed test space, $\dprime$ and $\ddprime$ differentiate only the
smooth coefficients because the $\omega_{J_i}$ are closed; they send an
order-$(s+1)$ seminorm continuously to an order-$s$ seminorm.
Multiplication by a fixed smooth function is continuous by the ordinary
product estimates for these seminorms.  The dual definitions therefore give
the asserted compatibilities.  Supports do not increase because the smooth
test forms are fewer.  Finally, a smooth positive test form is a positive
delta-form, so the three dual-cone implications follow directly from
\Cref{def:ch18-delta-current}; this is \cite[Remark~2.12]{GK19}.

\cref{itm:ch18-dc:form} The pairing is well defined by
\Cref{prop:ch18-delta-integration}\,\cref{itm:ch18-di:defined} applied to
$\beta\wedge\gamma$, which is a compactly supported generalized delta-form
by \Cref{prop:ch18-global-pb-structure}.  We verify the nontrivial
continuity assertion.  Fix a compact $C$, a chart cover, and tuples
$\omega_{J_i}$ as in \Cref{def:ch18-delta-current}.  Refine the cover so
that $\beta$ has a representative on every chart, and choose a finite
subcover $(V_i)_{i\in I_0}$ of $C$.  The chart-of-integration construction in
the proof of \cref{itm:ch18-di:defined} puts every product
$\beta\wedge\gamma$ on one common very-affine chart.  Its tropical support
is contained in a fixed compact polyhedral set $Q$, independent of
$\gamma$.

Write on $V_i$
$\gamma_i=\sum_{j\in J_i}\alpha_{ij}\wedge\omega_{ij}$.  On a common
subdivision, the coefficient of $\beta\wedge\gamma$ is a finite sum of
fixed smooth coefficients of $\beta$ and $\omega_{ij}$ multiplied by the
variable coefficients $\alpha_{ij}$ and fixed nonnegative lattice indices.
Integration over $Q$ therefore gives a constant $K$, depending only on the
fixed data, such that
\[
 \left|\int_{X^{\An}}\beta\wedge\gamma\right|
 \leq K\max_{i\in I_0,\,j\in J_i}p_{Q_i,0}(\alpha_{ij})
\]
for suitable compact $Q_i\subset\Omega_i$.  Taking the infimum over all
presentations of each $\gamma_i$ turns the right-hand side into finitely many
seminorms \eqref{eq:ch18-delta-test-seminorm}.  Hence the pairing is
continuous on every required test space.  This is the estimate of
\cite[Proposition~6.6]{GK17}.

If $\gamma$ is smooth then
$\beta\wedge\gamma$ is the preform with coefficients $\alpha_i\wedge\gamma$
and the same cycles, whose integral is \eqref{eq:ch18-associated-current}.

\cref{itm:ch18-dc:criterion} The forward implication follows from
functoriality.  Conversely, cover $W$ by tropical charts on which $\beta$
and $\beta'$ are represented by ambient delta-preforms.  Fix one such chart
and a compatible chart after an arbitrary algebraic morphism $f$.  The
hypothesis says that the associated pulled-back currents agree there.  By
the fixed-ambient equality criterion
\Cref{prop:ch18-current-faithful}, their difference has zero restriction
\eqref{eq:ch18-all-pullback-restriction}.  Since this holds for every $f$ and
every compatible chart, the difference belongs to the ideal
$\mathcal N$ of \Cref{def:ch18-global-pb}.  Hence the two chart classes
agree; sheaf gluing gives $\beta=\beta'$.  This is
\cite[Proposition~6.8]{GK17}.
\end{proof}

\begin{warning}{Delta-currents are not currents}
The test space defining a delta-current is $B_{\mathrm{GK},c}$, whereas a
CLD current is tested only on $A_c$.  Restriction therefore gives a
canonical map $E\to D$, but the cited theory supplies no general assertion
that this map is injective or noninjective.  Equality of delta-currents
always implies equality of their restrictions, so
\Cref{thm:gk-poincare-lelong} implies the CLD Poincar\'e--Lelong identity on
their common domain.  The separate question whether two generalized
delta-forms agree is settled by the all-pullback criterion
\cref{itm:ch18-dc:criterion}, not by a claim about $E\to D$.
\end{warning}

\subsection{The exact module boundary}
\label{sec:ch18-delta-current-module}

The dual definition of a delta-current does not make it a module over every
coefficient object occurring in the enlarged polyhedral calculi.  In the
original notation of \cite{GK17}, \(B_{\mathrm{GK}}\) is the differential
algebra of delta-forms, \(P\) is the larger space of generalized
delta-forms, and \(P\!S\!P\) also contains the piecewise-smooth algebra.
The canonical action on delta-currents is the action of
\(B_{\mathrm{GK}}\), not an action of all of \(P\) or \(P\!S\!P\).  The
following calculation records the complete Leibniz rule, including the
signs and the precise current-level detection hypothesis.

% Exact statement for production Rethlas and source audit, 26 August 2026.
% This proposition isolates the original Gubler--Kunnemann module category;
% it does not assert the unresolved continuous analytic Bedford--Taylor case.

\begin{proposition}
\label{prop:p21p3-delta-module}
Let \(K\) be an algebraically closed complete non-Archimedean field with
nontrivial absolute value, let \(X\) be an \(n\)-dimensional algebraic
\(K\)-variety, and let \(W\subset X^{\mathrm{an}}\) be open.  Write
\(B_{\mathrm{GK}}^{a,b}(W)\) for the Gubler--K\"unnemann delta-forms and
\(E_{\mathrm{GK}}^{p,q}(W)\) for their delta-currents.  Let
\[
 \alpha\in B_{\mathrm{GK}}^{a,b}(W),
 \qquad
 T\in E_{\mathrm{GK}}^{p,q}(W),
\]
where all bidegrees occurring below are in the range \(0,\ldots,n\).  Then
\begin{align}
 \mathrm d'(\alpha\wedge T)
 &=\mathrm d'\alpha\wedge T
   +(-1)^{a+b}\alpha\wedge\mathrm d'T,
 \label{eq:p21p3-dprime-module}\\
 \mathrm d''(\alpha\wedge T)
 &=\mathrm d''\alpha\wedge T
   +(-1)^{a+b}\alpha\wedge\mathrm d''T,
 \label{eq:p21p3-ddprime-module}
\end{align}
and consequently
\begin{align}
 \mathrm d'\mathrm d''(\alpha\wedge T)
 ={}&\mathrm d'\mathrm d''\alpha\wedge T
   +(-1)^{a+b+1}\mathrm d''\alpha\wedge\mathrm d'T
 \notag\\
 &+(-1)^{a+b}\mathrm d'\alpha\wedge\mathrm d''T
   +\alpha\wedge\mathrm d'\mathrm d''T.
 \label{eq:p21p3-dd-module}
\end{align}
In particular, if \(\mathrm d'T=\mathrm d''T=0\), then
\begin{equation}
 \mathrm d'\mathrm d''(\alpha\wedge T)
 =\mathrm d'\mathrm d''\alpha\wedge T.
 \label{eq:p21p3-closed-module}
\end{equation}
Thus form-level vanishing
\(\mathrm d'\mathrm d''\alpha=0\) in
\(B_{\mathrm{GK}}^{a+1,b+1}(W)\) forces
\(\mathrm d'\mathrm d''(\alpha\wedge T)=0\) as a delta-current.

Assume again that \(\mathrm d'T=\mathrm d''T=0\).  The same vanishing
conclusion follows from the functorial current-level hypothesis
\begin{equation}
 \bigl[f^*(\mathrm d'\mathrm d''\alpha)|_{W'}\bigr]_{D}=0
 \label{eq:p21p3-all-pullbacks}
\end{equation}
for every morphism \(f:X'\to X\) of algebraic \(K\)-varieties and every
open \(W'\subset (X')^{\mathrm{an}}\) with
\(f^{\mathrm{an}}(W')\subset W\), where the
left-hand side is the ordinary current on \(W'\) associated to the pulled
back generalized delta-form.

This statement supplies no multiplication of an arbitrary continuous or
piecewise-smooth function by an arbitrary delta-current.  Indeed,
\(B_{\mathrm{GK}}^{0,0}=A^{0,0}\), so a degree-zero multiplier in the
displayed module is smooth.  In this smooth degree-zero case there is a
useful sharper conclusion: if \(a=b=0\),
\(\mathrm d'T=\mathrm d''T=0\), and the single ordinary current
\begin{equation}
 \bigl[\mathrm d'\mathrm d''\alpha\bigr]_{D}=0
 \quad\text{on }W,
 \label{eq:p21p3-smooth-single-current}
\end{equation}
then \(\mathrm d'\mathrm d''(\alpha\wedge T)=0\).  By contrast, for a
general delta-form of general bidegree, vanishing of only its associated
ordinary current on \(W\) need not provide the form-level premise via the
original Gubler--K\"unnemann detection theorem; that theorem requires all
algebraic pullbacks.  Moreover, although the generalized piecewise-smooth
space \(P\!S\!P\) is an algebra and its elements define delta-currents, the
construction does not furnish a \(P\!S\!P\)-module structure on arbitrary
delta-currents extending the displayed \(B_{\mathrm{GK}}\)-module.  Hence
the proposition does not settle the continuous analytic Bedford--Taylor
question of
\cite[\S8.6.1 and Definition~8.6.2]{CLD25}.
No positivity, symmetry, or order hypothesis on \(T\) was used.
\end{proposition}


% Exact proof for production Rethlas and source audit, 26 August 2026.

\begin{proof}
Put \(A=a+b\) and \(P=p+q\).  The delta-form module structure on
delta-currents is defined in \cite[\S6.4, equation~(6.4.1)]{GK17} by
\begin{equation}
 \langle\alpha\wedge T,\beta\rangle
 =(-1)^{AP}\langle T,\alpha\wedge\beta\rangle
 \label{eq:p21p3-module-pairing}
\end{equation}
for every compactly supported delta-form \(\beta\) of complementary
bidegree.  The current differentials are defined in \cite[\S6.9]{GK17} by
\begin{equation}
 \langle\mathrm d'T,\gamma\rangle
 =(-1)^{P+1}\langle T,\mathrm d'\gamma\rangle,
 \qquad
 \langle\mathrm d''T,\gamma\rangle
 =(-1)^{P+1}\langle T,\mathrm d''\gamma\rangle.
 \label{eq:p21p3-current-differential}
\end{equation}
The delta-forms themselves are a bigraded differential algebra by
\cite[\S4.13 and Proposition~4.15(iii)]{GK17}.

Let \(\beta\) have the bidegree complementary to
\(\mathrm d'(\alpha\wedge T)\).  From
\eqref{eq:p21p3-module-pairing} and
\eqref{eq:p21p3-current-differential},
\begin{equation}
 \langle\mathrm d'(\alpha\wedge T),\beta\rangle
 =(-1)^{A+P+1+AP}
   \langle T,\alpha\wedge\mathrm d'\beta\rangle.
 \label{eq:p21p3-left-sign}
\end{equation}
On the other hand,
\begin{align*}
 &\left\langle
   \mathrm d'\alpha\wedge T+(-1)^A\alpha\wedge\mathrm d'T,
   \beta\right\rangle\\
 &\quad=(-1)^{(A+1)P}
   \langle T,\mathrm d'\alpha\wedge\beta\rangle
   +(-1)^{A+A(P+1)+P+1}
   \langle T,\mathrm d'(\alpha\wedge\beta)\rangle.
\end{align*}
Expanding
\(
 \mathrm d'(\alpha\wedge\beta)
 =\mathrm d'\alpha\wedge\beta+(-1)^A
   \alpha\wedge\mathrm d'\beta
\), the two terms containing
\(\mathrm d'\alpha\wedge\beta\) cancel, while the remaining sign is
\((-1)^{A+P+1+AP}\).  Comparison with
\eqref{eq:p21p3-left-sign} proves
\eqref{eq:p21p3-dprime-module}.  Replacing \(\mathrm d'\) everywhere by
\(\mathrm d''\) gives \eqref{eq:p21p3-ddprime-module}; the total-degree
signs are identical.

Apply \eqref{eq:p21p3-dprime-module} to
\eqref{eq:p21p3-ddprime-module}.  Since \(\mathrm d''\alpha\) has total
degree \(A+1\), one obtains
\begin{align*}
 \mathrm d'\mathrm d''(\alpha\wedge T)
 ={}&\mathrm d'\mathrm d''\alpha\wedge T
   +(-1)^{A+1}\mathrm d''\alpha\wedge\mathrm d'T\\
 &+(-1)^A\mathrm d'\alpha\wedge\mathrm d''T
   +\alpha\wedge\mathrm d'\mathrm d''T,
\end{align*}
which is \eqref{eq:p21p3-dd-module}.  If
\(\mathrm d'T=\mathrm d''T=0\), then also
\(\mathrm d'\mathrm d''T=0\), and
\eqref{eq:p21p3-closed-module} follows.  Its form-level consequence is
immediate.

For the functorial premise, apply
\cite[Proposition~6.8]{GK17} to the generalized delta-form
\(\eta=\mathrm d'\mathrm d''\alpha\) and to zero.  That proposition says
precisely that equality of generalized delta-forms is detected by equality
of the associated ordinary currents after every algebraic pullback and on
every open subset lying over \(W\).  Hypothesis
\eqref{eq:p21p3-all-pullbacks} therefore gives
\(\mathrm d'\mathrm d''\alpha=0\) as a delta-form, so the already proved
case applies.

Finally, \cite[equation~(4.19.1)]{GK17} gives
\(B_{\mathrm{GK}}^{0,0}=A^{0,0}\).  Thus the module used above has only
smooth degree-zero multipliers.  Suppose now that \(a=b=0\) and
\eqref{eq:p21p3-smooth-single-current} holds.  Then
\(\mathrm d'\mathrm d''\alpha\) is a smooth form.  The local tropical
pairing sends it to the ordinary current in
\eqref{eq:p21p3-smooth-single-current}.  A boundaryless pure-dimensional
Berkovich space is very rich in the sense of Chambert--Loir--Ducros, and
the map from forms with tropically continuous coefficients to currents is
injective on such a space by \cite[Corollary~7.3.9]{CLD25}.  Applied on
the open set \(W\), this shows that
\eqref{eq:p21p3-smooth-single-current} implies
\(\mathrm d'\mathrm d''\alpha=0\) as a smooth form and hence as a
delta-form.  The closed-current case already proved gives the asserted
smooth degree-zero conclusion.

For general bidegree, Remark~6.7 and Proposition~6.8 of \cite{GK17} give
the all-pullback criterion used above, not a single-current detection
statement.  Piecewise-smooth functions instead lie in
\(P\!S^{0,0}=P^{0,0}\) by \cite[equation~(9.3.1)]{GK17}.  Section~9.5 of
\cite{GK17} constructs the algebra \(P\!S\!P\) and associates a
delta-current to each of its elements, but it does not extend the module
pairing of equation~(6.4.1) to a pairing of \(P\!S\!P\) with arbitrary
delta-currents.  In addition, the polyhedral differentials of an associated
piecewise-smooth delta-current have residue terms and are not compatible
with current differentiation in general \cite[Remark~9.6]{GK17}.  These
observations prove the scope assertions and keep the analytic question of
\cite[\S8.6.1]{CLD25} outside the conclusion.
\end{proof}


The degree-zero clause is especially restrictive:
\(B_{\mathrm{GK}}^{0,0}=A^{0,0}\).  Thus a smooth pluriharmonic difference
cannot create a defect against a closed delta-current, but the proposition
does not permit the same inference for an arbitrary continuous or real-PL
difference.  For such a function, multiplication by an order-zero CLD
current is defined in the analytic current category; multiplication by an
arbitrary delta-current requires a separately constructed action and cannot
be inferred from the \(B_{\mathrm{GK}}\)-module.  This distinction is used
when formulating the residual Bedford--Taylor problem in
\Cref{prob:frontier-cross-fan-upstairs}.

\section{Poincar\'e--Lelong in the enlarged calculus}
\label{sec:ch18-delta-pl}

\begin{theorem}\label{thm:gk-poincare-lelong}
For a nonzero rational function \(f\) on \(X\),
\[
                   \dprimeddprime[\log|f|]=\delta_{\operatorname{cyc}(f)}
\]
as delta-currents.
\end{theorem}
\begin{proof}
The assertion is local on $X$.  By writing a rational function as a quotient
of two regular functions, using additivity of both sides, and treating the
irreducible components separately, we may assume that $X=\Spec A$ is
integral and that $f\in A$ is nonconstant.  Put
$U=X\setminus\Supp(\divis(f))$.  The function $\log|f|$ is smooth on
$U^{\An}$, and the support-removal theorem for delta-currents identifies
$E^{0,0}(U^{\An})$ with $E^{0,0}(X^{\An})$; this defines
$[\log|f|]$ across the divisor.

Both currents in the asserted equality are symmetric.  It is therefore
enough to pair them with a symmetric
$\alpha\in B_c^{n-1,n-1}(X^{\An})$.  Choose finitely many tropical charts
$(V_i,\varphi_{U_i})$ covering $\Supp\alpha$ and compact polytopal affinoid
domains $W_i\Subset V_i$ whose interiors still cover that support.  Since
$X^{\An}$ is separated, every nonempty intersection
$W_E=\bigcap_{i\in E}W_i$ is affinoid.  Shrink the polytopes slightly if
necessary and put
\[
       U_0:=U_1\cap\cdots\cap U_m\setminus\Supp(\divis(f)).
\]
This is very affine.  On the common chart, let
$F:N_{U_0,\mathbf R}\to\mathbf R$ be the integral affine map satisfying
\[
 F\circ\operatorname{trop}_{U_0}=-\log|f|,
\]
and let $\alpha_{U_0}$ represent $\alpha$.  The top-degree support-removal
statement implies that $\Supp(\ddprime\alpha)$ is contained in $U_0^{\An}$.
Writing $W=\bigcup_iW_i$, the polytopes may be chosen so that
$\Supp\alpha_{U_0}$ lies in the relative interior of
$\operatorname{trop}_{U_0}(W\cap U_0^{\An})$.
The compact set $\Supp(\ddprime\alpha)$ is contained in $U_0^{\An}$, so
there is $s>0$ such that $|f|>s$ on it.  Set
\[
 \begin{aligned}
 W_i(s)&:=\{x\in W_i:|f(x)|\geq s\},&
 W(s)&:=\bigcup_iW_i(s),\\
 P_s&:=\operatorname{trop}_{U_0}(W(s)\cap U_0^{\An}).
 \end{aligned}
\]
We give $P_s$ its canonical weighted polyhedral structure.

By the definition of the current differential and the sign convention
$F\circ\operatorname{trop}_{U_0}=-\log|f|$,
\begin{equation}
\label{eq:ch18-pl-pairing}
 \left\langle\dprimeddprime[\log|f|],\alpha\right\rangle
       =-\int_{P_s}F^*(x_0)\dprimeddprime\alpha_{U_0}.
\end{equation}
Green's formula for delta-preforms, applied cell by cell, gives
\[
 \int_{P_s}F^*(x_0)\dprimeddprime\alpha_{U_0}
 =\int_{\partial P_s}
     \bigl(F^*(x_0)\ddprime\alpha_{U_0}
       -\ddprime F^*(x_0)\wedge\alpha_{U_0}\bigr),
\]
because $\dprimeddprime F^*(x_0)=0$.  The first boundary term is zero:
$\Supp(\ddprime\alpha_{U_0})$ lies in the relative interior of $P_s$.
Consequently \eqref{eq:ch18-pl-pairing} becomes
\begin{equation}
\label{eq:ch18-pl-boundary-pairing}
 \left\langle\dprimeddprime[\log|f|],\alpha\right\rangle
   =\int_{\partial P_s}\ddprime F^*(x_0)\wedge\alpha_{U_0}.
\end{equation}
The boundary--fibre calculation in
\Cref{lem:ch18-boundary-divisor} identifies the last member with
$\langle\delta_{\operatorname{cyc}(f)},\alpha\rangle$.  This proves the
equality on every symmetric test form and hence as delta-currents.  Notice
that no torus chart on which $f$ is invertible can by itself see the divisor:
there $\log|f|$ is integral affine and has no corner.  The boundary limit is
the indispensable step.
\end{proof}

\subsection{The analytic fibre and boundary calculation}
\label{sec:ch18-pl-ingredients}

\begin{lemma}
\label{lem:ch18-fibre-cylinder}
Let $W$ be an affinoid domain in a pure $n$-dimensional analytic space,
let $g:W\to T^{\An}$ be an analytic moment map, and let $f$ be a nonzero
analytic function whose zero locus is an effective Cartier divisor on $W$.
There are $r>0$ and a weighted integral polyhedral complex $(\mathscr C,m)$
of pure dimension $n-1$ such that:
\begin{enumerate}
\item for every analytic point $t$ of the closed disc $|t|\leq r$, the
 $(n-1)$-skeleton of $g_{\trop}(W_t)$, with its canonical calibration, is
 $(\mathscr C,m)$;
\item for every closed interval $I\subset(0,r]$, the $n$-skeleton of
 $(-\log|f|,g_{\trop})(W_I)$ is the weighted product
 $(-\log I,1)\times(\mathscr C,m)$;
\item at $t=0$, the same $(\mathscr C,m)$ is the calibrated
 tropicalization of $\divis(f|_W)$.
\end{enumerate}
\end{lemma}

\begin{proof}
Apply the analytic tropical fibre theorem
\cite[Proposition~4.6.6]{CLD25} to the analytic map
$(f,g):W\to\mathbf A^{1,\An}\times T^{\An}$.  That theorem first gives
polyhedrality of the images of the fibres and of the annular domains, and
then says that their canonical calibrations are locally constant as the
parameter approaches zero.  After decreasing $r$, the finitely many cells
and their calibrations are therefore independent of $t$.  On an annulus the
coordinate $-\log|f|$ has primitive slope one; the calibration product rule
then gives $(-\log I,1)\times(\mathscr C,m)$.  The fibre over zero is the
closed analytic subspace defined by the Cartier divisor, and the
fundamental-cycle calibration includes its algebraic multiplicities.  This
gives the third assertion.  These are precisely the three conclusions of
the cited proposition, translated from calibrations to weighted complexes
using \Cref{def:trop-canonical-calibration}.
\end{proof}

\begin{lemma}
\label{lem:ch18-boundary-divisor}
With the notation used in the proof of
\Cref{thm:gk-poincare-lelong}, and after choosing $s>0$ small enough for
\Cref{lem:ch18-fibre-cylinder} on every nonempty $W_E$, one has
\begin{equation}
\label{eq:ch18-boundary-divisor}
 \int_{\partial P_s}\ddprime F^*(x_0)\wedge\alpha_{U_0}
       =\left\langle\delta_{\operatorname{cyc}(f)},\alpha\right\rangle .
\end{equation}
\end{lemma}

\begin{proof}
Put $U_E=\bigcap_{i\in E}U_i$ and
$W_E(s)=\bigcap_{i\in E}W_i(s)$.  Choose smooth functions
$\phi_\ell:\mathbf R\to[0,1]$ which equal $1$ on
$(-\infty,-\log s-1/\ell]$ and $0$ on
$[-\log s-1/(2\ell),\infty)$, and put $\psi_\ell=1-\phi_\ell$.  The form
$\phi_\ell(F^*x_0)\ddprime F^*x_0\wedge\alpha_{U_0}$ vanishes on the
truncation face $x_0=-\log s$; on every other boundary face it vanishes
because $\Supp\alpha_{U_0}$ lies in the relative interior of the
untruncated tropicalization.  Its boundary integral is therefore zero.
Hence the left side of
\eqref{eq:ch18-boundary-divisor} equals the limit, as $\ell\to\infty$, of
\[
 \int_{\partial P_s}\psi_\ell(F^*x_0)\ddprime F^*x_0
                         \wedge\alpha_{U_0}.
\]

Use finite inclusion--exclusion for the covering $W(s)=\bigcup_iW_i(s)$.
For every nonempty $E\subseteq\{1,\dots,m\}$ put
$V_E=\bigcap_{i\in E}V_i$, let $\alpha_E$ be the chart class determined by
$\alpha|_{V_E}=\operatorname{trop}_{U_E}^*\alpha_E$, and choose an ambient
delta-preform $\widetilde\alpha_E$ representing it.  Let
$g_E=\varphi_{U_E}$, $h_E=(f,g_E)$, and let
$\Sigma_E(s)=h_{E,\trop}(W_E(s))$.  The Sturmfels--Tevelev formula and the
projection formula carry the corresponding summand to
\[
 \int_{\partial\Sigma_E(s)}
       \psi_\ell(x_0)\ddprime x_0\wedge P_E^*\widetilde\alpha_E,
\]
where $P_E:\mathbf R\times N_{U_E,\mathbf R}\to N_{U_E,\mathbf R}$.
More explicitly, if $G_E:N_{U_0,\mathbf R}\to N_{U_E,\mathbf R}$ is the
compatible affine map and $H_E=(F,G_E)$, then
\[
 h_{E,\trop}(U_E\setminus\Supp(\divis(f)))
       =(H_E)_*\Trop(U_0).
\]
Applying the projection formula to this equality gives the displayed
boundary integral.  Passing to the limit
localizes this integral on the face $x_0=-\log s$.  By
\Cref{lem:ch18-fibre-cylinder}, that face, with its full lattice and
residue-degree calibration, is
\(
 \{-\log s\}\times
 \operatorname{trop}_{U_E}(\divis(f)\cap W_E)
\).
The limiting summand is therefore
\[
        \int_{\operatorname{trop}_{U_E}
          (\divis(f)\cap W_E\cap U_E^{\An})}\alpha_E .
\]
Finally sum with the inclusion--exclusion signs
$(-1)^{|E|+1}$.  On each irreducible component of $\divis(f)$ these terms
are exactly the inclusion--exclusion formula for integration over the part
covered by the $W_i$; those interiors cover $\Supp\alpha$.  Linearity in the
components, with their orders in $\operatorname{cyc}(f)$, gives
$\langle\delta_{\operatorname{cyc}(f)},\alpha\rangle$.  This proves
\eqref{eq:ch18-boundary-divisor}.  The calculation is the one in
\cite[Theorem~7.2 and Lemma~7.3]{GK17}; every multiplicity enters through
the calibrated fibre in \Cref{lem:ch18-fibre-cylinder}, rather than through
a nonexistent corner on the unit locus.
\end{proof}

\subsection{The constants, checked}
\label{sec:ch18-pl-constants}

The delta-form formula must reduce to the CLD formula and to the normalization
sheet.  We verify this rather than assert it, because the two calculi were
set up independently and there is no a priori reason for the constants to
match.

\begin{corollary}
\label{cor:ch18-pl-restricts}
Let $X$ be an algebraic variety over $K$ as above and $f\in K(X)^\times$.
\begin{enumerate}
\item\label{itm:ch18-plc:cld}
 Applying $\iota^\ast$ of
 \Cref{prop:ch18-delta-vs-cld-currents}\,\cref{itm:ch18-dc:restrict} to
 \Cref{thm:gk-poincare-lelong} gives
 $\dprimeddprime[\log|f|]=\delta_{\divis(f)}$, which is
 \Cref{thm:cld-poincare-lelong} and \eqref{eq:appJ-cld-pl}.
\item\label{itm:ch18-plc:val}
 Equivalently, in the valuation convention
 $v(f)=-\log|f|$ of \eqref{eq:appJ-valuation-sign},
 $\dprimeddprime[v(f)]=-\delta_{\divis(f)}$, which is
 \eqref{eq:appJ-valuation-pl}.
\item\label{itm:ch18-plc:section}
 For a delta-metrized line bundle $\overline L$ and a nonzero rational
 section $s$,
 \[
  c_1^\delta(\overline L)
  =\dprimeddprime[-\log\|s\|]+\delta_{\divis(s)},
 \]
 which is \eqref{eq:delta-c1-section}, and its image under $\iota^\ast$ is
 \eqref{eq:appJ-cld-section-pl}, that is, \Cref{cor:ch15-pl-section}.
\item\label{itm:ch18-plc:mass}
 If $X$ is a proper curve then $\int_{X^{\An}}\delta_{\divis(f)}=0$, in
 agreement with \Cref{cor:ch15-degree-zero}.
\end{enumerate}
No factor of $2$, of $2\pi$, or of $\log$-base occurs anywhere in these
identities.
\end{corollary}

\begin{proof}
\cref{itm:ch18-plc:cld} By
\cite[Example~6.11 and Remark~6.7]{GK17}, the restriction of the algebraic
integration delta-current $\delta_{\operatorname{cyc}(f)}$ to smooth test
forms is the CLD cycle current of
\Cref{prop:ch14-delta-Y-current}, namely $\delta_{\divis(f)}$.
Restriction commutes with $\dprimeddprime$ by
\Cref{prop:ch18-delta-vs-cld-currents}\,\cref{itm:ch18-dc:restrict}, and
$\iota^\ast[\log|f|]_\delta=[\log|f|]$ because both are the
$L^1_{\loc}$-current of the same function.  Hence the identity restricts as
claimed.  Comparing with \eqref{eq:appJ-cld-pl} shows that the two
normalizations agree.

\cref{itm:ch18-plc:val} Multiply \cref{itm:ch18-plc:cld} by $-1$ and use
$v(f)=-\log|f|$; this is \Cref{cor:ch15-appJ-consistency}.

\cref{itm:ch18-plc:section} Work locally on an open where $L$ has a frame
$e$; write $s=fe$ and $\lambda_e=-\log\|e\|$.  Then
$-\log\|s\|=\lambda_e-\log|f|$ on the locus where $f$ is invertible, and
$c_1^\delta(\overline L)=\dprimeddprime[\lambda_e]$ by the definition of the
Chern delta-form in \Cref{sec:ch18-chern-delta}.  Subtracting
\Cref{thm:gk-poincare-lelong} gives the displayed formula; the passage to
$\iota^\ast$ is \cref{itm:ch18-plc:cld} and produces
\eqref{eq:appJ-cld-section-pl}.

\cref{itm:ch18-plc:mass} Evaluate the current identity on the constant test
function $1$, which is admissible because $X^{\An}$ is compact.  By the
definition of the current differential,
\[
 \int_{X^{\An}}\delta_{\divis(f)}
 =\langle\dprimeddprime[\log|f|],1\rangle
 =\pm\langle[\log|f|],\dprimeddprime1\rangle=0.
\]
This is also \Cref{cor:ch15-degree-zero} after restriction to CLD currents.

Finally, no constant intervenes because every step is an equality of
polyhedral weights: the corner weight \eqref{eq:ch18-corner-weight} is a sum
of integer slopes against primitive normals, and the tropical multiplicity
\eqref{eq:trop-multiplicity-formula} is a sum of lengths times residue
degrees.  Both are the numbers appearing in
\Cref{sec:appJ-cld} without further normalization.
\end{proof}

\subsection{Worked verifications}
\label{sec:ch18-pl-examples}

\begin{example}[A character on the torus]
\label{ex:ch18-pl-character}
Let $X=\mathbf G_{m,K}^r$ and $f=\chi^u$ a character, $u\in M$.  Then
$\Trop(X)=\mathbf R^r$ with weight $1$, and
$\log|f|=-\langle u,\cdot\rangle$ is a linear form.  Its corner cycle is
$0$ by \Cref{lem:ch18-corner-welldefined}\,\cref{itm:ch18-cc:affine}, and its
facewise Hessian is $0$; so $\dprimeddprime[\log|f|]=0$.  This is correct:
$\operatorname{cyc}(\chi^u)=0$ on the torus, since a character is
invertible.  The example is the first sanity check on the sign convention,
and it fails for any convention in which $\log|f|$ is replaced by
$|f|$ or by $\log|f|^2$ without a compensating change on the right.
\end{example}

\begin{example}[A rational function on $\mathbf P^1$]
\label{ex:ch18-pl-p1}
Let $X=\mathbf P^1_K$ and $f=(T-a)/(T-b)$ with $|a|,|b|\leq1$ and
$a\neq b$.  Then $\operatorname{cyc}(f)=[a]-[b]$ and
\Cref{thm:gk-poincare-lelong} asserts
$\dprimeddprime[\log|f|]=\delta_a-\delta_b$ as delta-currents.  Restricting
to the standard apartment and writing $r=|a-b|$, the function $\log|f|$ is
$0$ on the part of the tree outside the disc of radius $r$ about $a$, and on
the two branches from the Gauss point of that disc it has slope $-1$
towards the zero $a$ and slope $+1$ towards the pole $b$.  Both atoms have absolute value $1$
and the total mass is $0$, as \Cref{cor:ch15-degree-zero} and
\Cref{cor:ch18-pl-restricts}\,\cref{itm:ch18-plc:mass} require.  This is the
computation of \Cref{ex:ch15-rational-p1}, now read in the delta-current
test category.  Its right-hand side is the algebraic integration
delta-current $\delta_{[a]-[b]}$, so it may be paired with compactly
supported delta-forms; prescribed Green-current intersections are also
available when the relevant algebraic cycles meet properly.  It is not being
asserted that this algebraic current is a generalized delta-form, and the bare
measure $\delta_a-\delta_b$ has no product with an arbitrary current.  No
injectivity assertion for $E\to D$ is being used.
\end{example}

\begin{example}[The tropical line, once more]
\label{ex:ch18-pl-line}
Let $U=\{1+x+y=0\}\cap\mathbf G_{m,K}^2$, so that $\Trop(U)=L$ is the
standard tropical line \eqref{eq:ch18-gauge-rays} with unit weights.  Take
$f=x|_U$, invertible on $U$.  Then $\log|f|=-x_1$ in the tropical
coordinates, an affine function, so
$\dprimeddprime[\log|f|]=0$: correct, since $x$ is invertible on $U$.  Take
instead $g=1+x=-y$ on $U$, again invertible; $\log|g|=-x_2$, again affine,
again giving $0$.  Now pass to the compactification
$\overline U\subset\mathbf P^2$ and take the restriction of the coordinate
$Z_0/Z_2$: the corresponding function is no longer invertible on
$\overline U$.  Its divisor is
\[
 [\overline U\cap\{Z_0=0\}]-[\overline U\cap\{Z_2=0\}],
\]
two boundary points with weights $+1$ and $-1$ and total degree zero.  The
example shows that on the unit locus all the
Poincar\'e--Lelong content lies at the boundary of the chart.  It is the
calibrated boundary--fibre limit of
\Cref{lem:ch18-boundary-divisor}, not the addition of another invertible
torus coordinate, that recovers the divisor.
\end{example}

\begin{example}[The gauge is not a Poincar\'e--Lelong term]
\label{ex:ch18-gauge-not-pl}
The function $u=g|_L$ of \Cref{ex:ch18-gauge-atom} has
$\dprimeddprime[u]=3\delta_0$, which by
\Cref{thm:gk-poincare-lelong} is \emph{not} of the form
$\delta_{\operatorname{cyc}(f)}$ for a rational function $f$ on $U$.
Indeed, an algebraic divisor current on the curve is supported at rigid
type-$1$ points, whereas the CLD measure associated to the polyhedral atom
$\delta_0$ is supported at the type-$2$ Shilov point of the tropical fibre
over the vertex.  The supports are disjoint.  This is the difference between
a curvature current and an algebraic divisor current, and it is the reason the Chern delta-form of
\Cref{sec:ch18-chern-delta} is needed: the mass $3$ is a positivity datum of
a metric, not the divisor of a function.
\end{example}

\section{Chern delta-forms and formal metrics}
\label{sec:ch18-chern-delta}

\subsection{Delta-metrics}
\label{sec:ch18-delta-metrics}

For a piecewise smoothly metrized line bundle \(\overline L\), a local frame
has a piecewise smooth weight \(\phi=-\log\|e\|\).  Its current derivative
splits, as in \eqref{eq:tropical-delta-pl}, into a facewise smooth part and a
corner-locus part.  The former always globalizes in the piecewise smooth
sheaf, and the latter globalizes canonically as a residual generalized
delta-form of codimension one.  The extra requirement for an
\emph{ordinary delta-metric} is that the facewise Hessian be the restriction
of an ambient superform with piecewise smooth coefficients.  Only under that
condition does their sum define the first Chern delta-form
\(c_1^\delta(\overline L)\).

Keeping these three assertions, and their different codimension conventions,
separate is the delicate point.

\begin{lemma}
\label{lem:ch18-delta-metric-gluing}
Let \(K\) be an algebraically closed, complete, nontrivially valued
non-Archimedean field, let \(X\) be a pure \(n\)-dimensional algebraic
variety over \(K\), let \(W\subset X^{\An}\) be open, and let
\(\overline L=(L,\|\cdot\|)\) be a piecewise smoothly metrized line bundle
on \(W\).  The following statements hold.
\begin{enumerate}
\item\label{itm:ch18-dm:frame}
 If \((V,\operatorname{trop}_U)\) is a tropical chart in \(W\), \(e\) is a
 nowhere-vanishing frame of \(L\) over \(U\), and
 \[
       -\log\|e\|=\phi_e\circ\operatorname{trop}_U
       \qquad\hbox{on }V,
 \]
 then the facewise Hessians
 \(\mathrm d'_{\mathrm P}\mathrm d''_{\mathrm P}\phi_e\) define a global
 piecewise smooth form
 \[
       c_1(L,\|\cdot\|)_{\mathrm{ps}}\in
       \operatorname{PS}^{1,1}(W).
 \]
 This class is independent of the frame and of the tropical chart after
 passage to common refinements.  It is not thereby asserted to be a
 generalized delta-form of codimension zero.

\item\label{itm:ch18-dm:current}
 There is a canonical generalized delta-form
 \[
       c_1(L,\|\cdot\|)_{\mathrm{res}}
       \in P^{1,1}(W)
 \]
 of codimension one.  It is uniquely characterized by the following
 universal pullback property: for every morphism \(f:X'\to X\) from an
 algebraic variety over \(K\) and every open subset
 \(W'\subseteq f^{-1}(W)\), the delta-current associated to
 \[
        f^*c_1(L,\|\cdot\|)_{\mathrm{res}}\big|_{W'}
 \]
 is the residual part of the first Chern current of the pulled-back
 metrized line bundle \(f^*\overline L\) on \(W'\).  For the identity
 pullback one obtains
 \begin{equation}
 \label{eq:ch18-dm-current-decomposition}
 [c_1(L,\|\cdot\|)]
   =[c_1(L,\|\cdot\|)_{\mathrm{ps}}]
    +[c_1(L,\|\cdot\|)_{\mathrm{res}}].
 \end{equation}

\item\label{itm:ch18-dm:delta}
 The metric is an ordinary delta-metric if and only if every point of \(W\)
 has data as in \cref{itm:ch18-dm:frame} for which there is an ambient
 superform \(\gamma\) of bidegree \((1,1)\), with piecewise smooth
 coefficients on \(N_{U,\mathbf R}\), such that
 \begin{equation}
 \label{eq:ch18-dm-hessian-lift}
    \gamma|_{\operatorname{Trop}(U)}
      =\mathrm d'_{\mathrm P}\mathrm d''_{\mathrm P}\phi_e
      \quad\hbox{on }\operatorname{trop}_U(V).
 \end{equation}
 Equivalently, \(c_1(L,\|\cdot\|)_{\mathrm{ps}}\) is induced by a
 generalized delta-form of codimension zero.  In that case the sum of this
 codimension-zero form and
 \(c_1(L,\|\cdot\|)_{\mathrm{res}}\) is the uniquely determined first Chern
 delta-form \(c_1^\delta(\overline L)\).

\item\label{itm:ch18-dm:degree}
 The ambient and relative bidegrees in the preceding assertions are related
 as follows.  In a chart satisfying \eqref{eq:ch18-dm-hessian-lift}, let
 \(\widetilde\Omega\subset N_{U,\mathbf R}\) be open and chosen so that,
 with
 \[
       C:=\operatorname{Trop}(U)\cap\widetilde\Omega,
       \qquad c:=\operatorname{codim}_{N_{U,\mathbf R}} C,
 \]
 one has \(C\subseteq\operatorname{trop}_U(V)\).  Let
 \(\widetilde\phi_e\) be a piecewise smooth ambient extension of
 \(\phi_e\).  Whenever \eqref{eq:ch18-dm-hessian-lift} holds, the local
 current on \(C\) is represented, after a common subdivision, by
 \begin{equation}
 \label{eq:ch18-dm-ambient-representative}
 \begin{aligned}
   \widetilde\Theta_{U,e}
    &:=(\gamma\wedge\delta_C)
       +\delta_{(\widetilde\phi_e\cdot N_{U,\mathbf R})\cdot C}\\
     &=(\mathrm d'_{\mathrm P}\mathrm d''_{\mathrm P}\phi_e)
          \wedge\delta_C
       +\delta_{\phi_e\cdot C}.
 \end{aligned}
 \end{equation}
 Both sides have ambient bidegree \((c+1,c+1)\) and relative bidegree
 \((1,1)\) on \(C\).  The first summand is supported in ambient codimension
 \(c\), while the second is supported in ambient codimension \(c+1\).
 Thus the first summand has ambient support-codimension zero only when
 \(c=0\); it is the coefficient form \(\gamma\), before intersection with
 \(C\), that is a codimension-zero generalized delta-preform in every
 ambient dimension.
\end{enumerate}
\end{lemma}

\paragraph{Imported inputs.}
The proof uses the following results from Gubler--K\"unnemann
\cite{GK17}, with their original hypotheses.
\begin{enumerate}
\item Their \S9.7 constructs the global piecewise smooth part of the first
 Chern current and the current decomposition (9.7.1).
\item Their Proposition~9.8 constructs the residual generalized delta-form,
 proves that it has codimension one, and proves its uniqueness and
 functoriality.
\item Their Definition~9.9 gives the ambient-Hessian criterion for a
 delta-metric.  Their Remark~9.10 identifies its ambient superform with a
 generalized delta-preform of codimension zero.  Their Proposition~9.11 and
 \S9.13 identify the resulting sum with the unique first Chern delta-form.
\item Their Proposition~1.8 extends a piecewise smooth function from a
 tropical cycle to the ambient vector space.  Their Proposition~1.14 and
 Corollary~1.15 identify the corner locus after restriction with stable
 intersection of the ambient corner locus.  Their Corollary~3.19 is the
 tropical Poincar\'e--Lelong decomposition into facewise Hessian and corner
 locus.
\end{enumerate}
The bidegree convention is the one in \cite[Definition~2.9 and
Example~2.10]{GK17}: if \(D\) has ambient codimension \(l\), then
\(\alpha\wedge\delta_D\), with \(\alpha\) of bidegree \((p-l,q-l)\), has
ambient bidegree \((p,q)\).

\begin{proof}
\cref{itm:ch18-dm:frame}
The construction is \cite[\S9.7]{GK17}.  We spell out its invariance
because it prevents a misleading fixed-chart assertion.  If \(e'=ge\),
then
\[
  -\log\|e'\|=-\log\|e\|-\log|g|.
\]
After adjoining the invertible regular function \(g\) to a common very
affine refinement, \(-\log|g|\) is the restriction of an integral affine
function on the new ambient tropical vector space.  Its facewise Hessian
vanishes.  Hence the two local facewise Hessians agree on that common
refinement.  The same argument for two tropical charts, using their common
very-affine refinement, gives a global element of
\(\operatorname{PS}^{1,1}(W)\).  This construction takes place in the
piecewise smooth sheaf.  It supplies no ambient bidegree-\((1,1)\) lift, so
it does not by itself put the class in the codimension-zero part of the
generalized delta-form sheaf.

\cref{itm:ch18-dm:current}
Equation \eqref{eq:ch18-dm-current-decomposition} is
\cite[(9.7.1)]{GK17}.  Proposition~9.8 there constructs a generalized
delta-form of codimension one and proves exactly the displayed universal
pullback property.  It also proves uniqueness among generalized delta-forms
satisfying that property.  Taking \(f=\operatorname{id}_X\) gives the
residual summand in \eqref{eq:ch18-dm-current-decomposition}.  No injectivity
of the map from generalized delta-forms to associated currents on \(W\)
alone is used or asserted.

\cref{itm:ch18-dm:delta}
Condition \eqref{eq:ch18-dm-hessian-lift} is exactly
\cite[Definition~9.9(ii)]{GK17}; condition (i) of that definition is the
piecewise smooth local presentation already recorded in
\cref{itm:ch18-dm:frame}.  A superform with piecewise smooth coefficients on
the ambient vector space is a generalized delta-preform of codimension zero
by \cite[Remark~9.10 and Example~2.10]{GK17}.  That remark therefore proves
the stated equivalence.  Proposition~9.11 and \S9.13 of the same paper show
that the sum of this codimension-zero piece and the residual codimension-one
piece represents the full first Chern delta-current and is a delta-form.
Its uniqueness follows from \cite[Proposition~6.8 and
Definition~9.12]{GK17}.

\cref{itm:ch18-dm:degree}
By \cite[Proposition~1.8]{GK17}, after shrinking if necessary,
\(\phi_e\) has a piecewise smooth ambient extension
\(\widetilde\phi_e\).  The ambient corner locus
\(\widetilde\phi_e\cdot N_{U,\mathbf R}\) has codimension one.  Its stable
intersection with \(C\) is the corner locus of the restricted function:
\[
   (\widetilde\phi_e\cdot N_{U,\mathbf R})\cdot C
      =\phi_e\cdot C,
\]
by \cite[Proposition~1.14 and Corollary~1.15]{GK17}.  The facewise
restriction in \eqref{eq:ch18-dm-hessian-lift} gives the next identity on
all of \(C\), because \(C\subseteq\operatorname{trop}_U(V)\):
\[
   \gamma\wedge\delta_C
    =(\mathrm d'_{\mathrm P}\mathrm d''_{\mathrm P}\phi_e)
       \wedge\delta_C.
\]
Together with the tropical Poincar\'e--Lelong formula
\cite[Corollary~3.19]{GK17}, these two identities prove
\eqref{eq:ch18-dm-ambient-representative}.

It remains only to keep the two degree conventions separate.  The cycle
\(C\) has ambient codimension \(c\), so a bidegree-\((1,1)\) coefficient
form wedged with \(\delta_C\) has ambient bidegree
\((c+1,c+1)\).  The cycle
\((\widetilde\phi_e\cdot N_{U,\mathbf R})\cdot C\) has ambient codimension
\(c+1\), so its integration preform has the same ambient bidegree.  After
subtracting the ambient codimension \(c\) of \(C\), both have relative
bidegree \((1,1)\).  Their support codimensions are respectively \(c\) and
\(c+1\).  On the other hand \(\gamma\), before wedging with
\(\delta_C\), is by definition a codimension-zero generalized
delta-preform.  This proves every assertion in
\cref{itm:ch18-dm:degree}.
\end{proof}


\begin{warning}{Not every continuous psh metric is a delta-metric}
Being a delta-metric is a regularity hypothesis of polyhedral type, not a
positivity hypothesis, and it is genuinely restrictive.  A uniform limit of
delta-metrics need not be one, because the corner cycles need not converge
as cycles even when the curvature currents converge weakly; the
non-Archimedean Calabi--Yau solutions of \Cref{chap:bfj-ma} are the standard
source of continuous psh metrics for which no polyhedral presentation is
known.  Every theorem in this section is stated for delta-metrics and none
of them is asserted for arbitrary continuous psh metrics.
\end{warning}

For a meromorphic section \(s\),
\begin{equation}
 c_1^\delta(\overline L)
   =\dprimeddprime[-\log\|s\|]+\delta_{\operatorname{div}(s)}.
\label{eq:delta-c1-section}
\end{equation}
This is the same sign as
\eqref{eq:cld-pl-section-sign}.  Formal metrics are delta-metrics in the
algebraic setting of \cite{GK17}; see its Proposition~8.11 and
Section~9.

\subsection{Mixed Chern delta-forms and Chambert--Loir measures}
\label{sec:ch18-delta-cl}

If \(X\) is proper of dimension \(n\) and
\(\overline L_1,\ldots,\overline L_n\) are delta-metrized, then
\[
                 c_1^\delta(\overline L_1)\wedge\cdots\wedge
                 c_1^\delta(\overline L_n)
\]
is a top-degree delta-form and hence a signed Radon measure.

Two remarks on this sentence.  The product exists because of
\Cref{thm:delta-preform-algebra} and \Cref{prop:ch18-product-descends}; no
Bedford--Taylor approximation or recursion is used, and no regularity beyond
the piecewise-smooth, hence continuous, delta-metric hypothesis is required.
The passage to a measure is
\Cref{prop:ch18-delta-integration}\,\cref{itm:ch18-di:measure}.  The measure
is signed because effectivity of the corner cycles has not been assumed;
positivity is exactly the content of \Cref{sec:delta-positivity}.

\begin{theorem}\label{thm:gk-delta-cl-measure}
For formally metrized line bundles on a proper \(n\)-dimensional algebraic
variety, the preceding mixed delta-form measure equals the corresponding
Chambert--Loir measure.
\end{theorem}
\begin{proof}
Choose one proper formal model carrying all the line bundles that define the
metrics.  On a formal chart, each metric weight is piecewise affine, so its
Chern delta-form is its weighted corner cycle by
\eqref{eq:tropical-delta-pl}.  Take a common subdivision.  The wedge of the
$n$ Chern delta-forms is the stable intersection of these corner cycles.

At a point corresponding to a generic component $Y$ of the special fibre,
the stable-intersection weight is the mixed degree
\[
       \deg_Y\bigl(c_1(\mathcal L_1|_Y)\cdots
                         c_1(\mathcal L_n|_Y)\cap[Y]\bigr),
\]
multiplied by the component length and the residue-field degree.  This is
the toric intersection formula in the graded reduction of the formal chart.
No other point contributes to a top-degree stable intersection.  Hence the
delta-form measure is the sum of these coefficients at the divisorial
points, which is precisely the Chambert--Loir measure.

The tropical/residual mixed-degree comparison is
\cite[Theorem~10.5]{GK17}.  It is invariant under a dominating formal model
by the projection formula, so the calculation is independent of the chosen
common model.  This also proves directly that the result agrees with
\Cref{thm:cld-recovers-cl}.
\end{proof}

\begin{corollary}
\label{cor:ch18-delta-total-mass}
In the situation of \Cref{thm:gk-delta-cl-measure},
\[
 \int_{X^{\An}}c_1^\delta(\overline L_1)\wedge\cdots\wedge
 c_1^\delta(\overline L_n)=(L_1\cdots L_n)_X ,
\]
the intersection number on the generic fibre.  In particular the mass does
not change when the formal model is changed.
\end{corollary}

\begin{proof}
By \Cref{thm:gk-delta-cl-measure} the measure is the Chambert--Loir measure
\eqref{eq:appJ-cl-measure}, whose total mass is
\eqref{eq:appJ-cl-total-mass}; that identity is
\Cref{thm:cl-model-properties}.  Independence of the model is the projection
formula, quoted in the proof of \Cref{thm:gk-delta-cl-measure}.
Alternatively, one can argue entirely within the delta-calculus: the mass is
$\int\dprimeddprime[\phi_1]\wedge\cdots$, and successive applications of
\Cref{prop:ch18-delta-integration}\,\cref{itm:ch18-di:stokes} reduce a change
of metric $\phi_1\rightsquigarrow\phi_1'$ to a term
$\int\dprimeddprime[\phi_1-\phi_1']\wedge(\cdots)$ which vanishes; hence the
mass depends only on the isomorphism classes of the $L_i$, and evaluating on
a model metric gives the intersection number.
\end{proof}

\subsection{The formal case computed}
\label{sec:ch18-formal-computed}

\begin{example}[The chain of three components, as delta-forms]
\label{ex:ch18-chain-delta}
Continue \Cref{ex:ch18-chain-numbers}, with the curve $C$ of genus zero and
$L=\cO(1)|_C$ of degree $2$, and let $\overline L$ be the formal metric
attached to $\mathcal L=\mathcal L_0\otimes\cO(D)$ with
$D=a_0C_0+a_1E+a_2C_2$.  The skeleton is $\Gamma=[0,2]$ with unit weights.
By \Cref{lem:ch18-delta-metric-gluing} and \eqref{eq:tropical-delta-pl}, the
Chern delta-form is the corner cycle of the metric weight, a $0$-cycle on
$\Gamma$; by \eqref{eq:ch18-corner-weight} its coefficients are the atoms.
Adding the untwisted part $\mu_0=\delta_{v_0}+\delta_{v_2}$ of
\Cref{ex:metrics-cl-curve-chain} to the atoms
\eqref{eq:ch18-chain-atoms} gives
\[
 c_1^\delta(\overline L)
 =(1+a_1-a_0)\{v_0\}+(a_0+a_2-2a_1)\{v_1\}+(1+a_1-a_2)\{v_2\},
\]
which is exactly \eqref{eq:metrics-cl-chain-measure}.  The total weight is
$2=\deg_CL$ for every $a$, as \Cref{cor:ch18-delta-total-mass} requires.

Two numerical instances.  For $a=(0,0,0)$ one gets
$\{v_0\}+\{v_2\}$: the metric is the restriction of the standard metric of
$\mathbf P^2$ and the middle component, being contracted, carries no mass.
For $a=(0,1,0)$ one gets $2\{v_0\}-2\{v_1\}+2\{v_2\}$, of total weight $2$
but not effective; consistently, the corresponding model line bundle is not
nef, since $\deg(\mathcal L|_E)=0-2+0=-2<0$.  This is the first appearance,
in numbers, of the positivity criterion of
\Cref{thm:gk-formal-positivity-comparison}: effectivity of the Chern
delta-form is equivalent to nefness of the residual bundles.
\end{example}

\begin{example}[A toric surface]
\label{ex:ch18-toric-delta}
Let $X$ be the toric surface of \Cref{ex:metrics-cl-toric-surface} with the
model metric whose weight on the skeleton $\mathbf R^2$ is the piecewise
affine function \eqref{eq:metrics-toric-weight}.  The Chern delta-form is
the corner cycle of that weight, a one-dimensional balanced cycle in
$\mathbf R^2$ whose weights are the lattice lengths of the edges of the
associated polytope, and the top self-intersection is computed by
\eqref{eq:ch18-stable-weight}.  The resulting mass is
\eqref{eq:metrics-toric-measure}, that is, twice the Euclidean area of the
polytope in the normalization of \Cref{sec:appJ-toric}.  The lattice indices
of \Cref{ex:ch18-nonunimodular} are precisely what makes the answer come out
as the normalized volume rather than as a naive edge count.
\end{example}

\begin{example}[Where the mass sits: the blowup of $\mathbf P^1$]
\label{ex:ch18-p1-blowup-delta}
Continue \Cref{ex:ch18-p1-blowup-defect}.  For
$\mathcal L'=\rho^*\cO(1)\otimes\cO(-E)$ the weight in the frame $Z_0$ was
computed in \Cref{ex:metrics-p1-blowup} to be $\max\{-1,\log|T|\}$, whose
corner cycle on the apartment is the single point $x_E$ with weight $1$;
the degrees of $\mathcal L'$ on the two components are
$\deg(\mathcal L'|_{\widetilde C})=1-1=0$ and
$\deg(\mathcal L'|_E)=0+1=1$.  Thus
$c_1^\delta(\overline{L'})=\{x_E\}$, of total weight $1=\deg\cO(1)$,
effective, and supported at a single divisorial point.  The corresponding
Chambert--Loir measure is $\delta_{x_E}$, in agreement with
\Cref{thm:gk-delta-cl-measure} and with \Cref{thm:cl-model-properties}.
\end{example}

\subsection{Green delta-currents and local heights}
\label{sec:ch18-green}

The delta-calculus also permits Green delta-currents and their
\(*\)-products, but only in a prescribed one-sided framework.  If $Z$ is
an algebraic cycle of codimension $p$, a delta-current
$g_Z\in E^{p-1,p-1}$ is a Green delta-current when
\[
          \omega(Z,g_Z):=\dprimeddprime g_Z+\delta_Z
\]
is the current of a delta-form $\omega_Z\in B^{p,p}$.

Let $\overline L=(L,\|\cdot\|)$ be a delta-metrized line bundle, let $s$ be
a nonzero meromorphic section, and put
\[
 D=\divis(s),\qquad Y=\operatorname{cyc}(s),\qquad
 g_Y=[-\log\|s\|],\qquad \omega_Y=c_1^\delta(\overline L).
\]
Suppose that $D$ meets $Z$ properly.  For a prime component $Z$ with
inclusion $i_Z:Z^{\An}\to X^{\An}$, define
\begin{equation}
\label{eq:ch18-green-restriction}
       g_Y\wedge\delta_Z
       :=(i_Z)_*[-\log\|s|_Z\|],
\end{equation}
and extend linearly in $Z$.  This is a restriction followed by pushforward,
not a product of arbitrary currents.  If $g_Z$ is a Green delta-current, set
\begin{equation}
\label{eq:ch18-green-star}
             g_Y*g_Z:=g_Y\wedge\delta_Z+\omega_Y\wedge g_Z
             \in E^{p,p}(X^{\An}).
\end{equation}
The second term is the canonical action of the differential algebra $B$ on
delta-currents.  The first is the special operation
\eqref{eq:ch18-green-restriction}.  Proper intersection alone does not
define $\delta_Z\wedge g_W$ for an arbitrary Green current $g_W$, and the
polyhedral algebra of \Cref{thm:delta-preform-algebra} does not supply such
a product.

We verify that \eqref{eq:ch18-green-star} has the required curvature.
Poincar\'e--Lelong on every prime component of $Z$, followed by proper
pushforward, gives
\begin{equation}
\label{eq:ch18-green-restricted-pl}
 \dprimeddprime(g_Y\wedge\delta_Z)+\delta_{D\cdot Z}
            =\omega_Y\wedge\delta_Z .
\end{equation}
Since $\omega_Y$ is closed and acts on delta-currents, the Leibniz rule and
\eqref{eq:ch18-green-restricted-pl} yield
\begin{align*}
 \omega(D\cdot Z,g_Y*g_Z)
 &=\dprimeddprime(g_Y\wedge\delta_Z)
    +\omega_Y\wedge\dprimeddprime g_Z+\delta_{D\cdot Z}\\
 &=\omega_Y\wedge(\delta_Z+\dprimeddprime g_Z)
  =\omega_Y\wedge\omega(Z,g_Z).
\end{align*}
Thus $g_Y*g_Z$ is a Green delta-current for $D\cdot Z$.  This is
\cite[Lemma~11.3 and Proposition~11.4]{GK17} with all operations displayed.
The operation is not literally a commutative product on arbitrary Green
currents; even for two divisor factors, interchange is controlled only
modulo the explicit $\dprime$- and $\ddprime$-exact correction.  If
$\eta_{Y_i}=-\log\|s_i\|$, then
\begin{equation}
\label{eq:ch18-green-commutator}
 g_{Y_1}*g_{Y_2}-g_{Y_2}*g_{Y_1}
 =\dprime[\ddprime_P\eta_{Y_1}\wedge\eta_{Y_2}]
  +\ddprime[\eta_{Y_1}\wedge\dprime_P\eta_{Y_2}],
\end{equation}
by \cite[Proposition~11.5]{GK17}.

For $n+1$ properly intersecting delta-metrized Cartier divisors on a proper
$n$-fold, one may iterate this divisor-first construction and evaluate the
right-associated top current
\[
       \bigl(g_{Y_0}*(g_{Y_1}*(\cdots*g_{Y_n}))\bigr)(1).
\]
The value is symmetric and agrees with the
model-theoretic local height for algebraic metrics by
\cite[Theorem~0.4 and \S12]{GK17}.  This is the precise connection with
non-Archimedean Arakelov intersection theory.  We use the construction
nowhere below, but record it because it is a principal application of the
delta-current category.

\section{Positivity of delta-forms}
\label{sec:delta-positivity}

\subsection{The definition, and what it must satisfy}
\label{sec:ch18-positivity-def}

Positivity for delta-preforms is defined by testing the intrinsic collected
coefficient on every cell of a common subdivision.  A presentation may use
signed cycles and signed weights; no sign is imposed term by term.  Weak,
ordinary, and strong positivity are distinguished in intermediate
bidegrees, exactly as for superforms.  On the analytic space the test is
repeated after every compatible algebraic pullback; the word
\emph{functorial} records that universal requirement.

Written out, the definition is the following.

\begin{definition}[Positive delta-preform]
\label{def:ch18-positive-delta}
\index{positivity!of delta-forms}%
A delta-preform $\beta=\sum_i\alpha_i\wedge\delta_{C_i}$ of bidegree $(p,p)$
on $\Omega$ is \emph{positive} (respectively \emph{strongly positive},
\emph{weakly positive}) if, after passing to a common subdivision
$\mathscr C$ as in \Cref{prop:ch18-current-faithful}, for every relevant face
$\sigma\in\mathscr C$ the collected coefficient
\[
 A_\sigma=\sum_{i:\,l_i=\operatorname{codim}\sigma}m_{i,\sigma}\alpha_i|_\sigma
\]
is a positive (respectively strongly, weakly positive) superform on the
vector space $N_{\sigma,\mathbf R}$ in the sense of
\Cref{def:ch09-positivity}.  There is no condition on the signs of the
individual weights in a chosen presentation: only the collected coefficient,
which is intrinsic to the polyhedral current, is tested.

A chart class in $P^{p,p}(V,\varphi_U)$ is positive in the corresponding
sense when every restriction \eqref{eq:ch18-all-pullback-restriction} after
every algebraic morphism is a positive delta-preform.  A generalized
delta-form on an analytic open set is positive when this holds locally on
tropical charts.  This is the functorial definition of
\cite[Definitions~2.2 and 2.6]{GK19}.
\end{definition}

\begin{lemma}
\label{lem:ch18-positivity-consistent}
\begin{enumerate}
\item\label{itm:ch18-pos:welldef}
 \Cref{def:ch18-positive-delta} is independent of the subdivision and
 of the presentation, and hence is well defined on $P(\Omega)$.
\item\label{itm:ch18-pos:current}
 On a fixed ambient open set, a delta-preform is weakly positive, positive,
 or strongly positive if and only if its associated polyhedral current has
 the corresponding positivity property.
\item\label{itm:ch18-pos:global}
 A generalized delta-form $\beta\in P^{p,p}(W)$ is weakly positive,
 positive, or strongly positive if and only if $[f^*\beta]_D$ has the
 corresponding positivity property on $f^{-1}(W)$ for every algebraic
 morphism $f:X'\to X$.
\item\label{itm:ch18-pos:top}
 On a fixed $n$-dimensional ambient tropical space, the three notions in
 bidegree $(n,n)$ coincide and are equivalent to positivity of the
 associated measure.  For a generalized delta-form on an analytic variety,
 positivity implies positivity of the base measure; the converse still
 requires the all-pullback test in \cref{itm:ch18-pos:global}.
\end{enumerate}
\end{lemma}

\begin{proof}
\cref{itm:ch18-pos:welldef} A refinement of $\mathscr C$ replaces
$A_\sigma$ by its restriction to subcells, so the condition is unchanged.
The collected coefficients depend only on the polyhedral current by
\Cref{prop:ch18-current-faithful}; in particular, cancellation between
positive and negative weights has already occurred before positivity is
tested.

\cref{itm:ch18-pos:current} If all $A_\sigma$ are positive, test
\eqref{eq:ch18-associated-current} against a strongly positive $\omega$.
Each $A_\sigma\wedge\omega|_\sigma$ is a positive top-degree superform by
\Cref{prop:ch09-positivity-products}, hence has nonnegative integral.  The
weak and strong variants use the dual positivity cones in the same way.

Conversely, order the cells by increasing dimension.  Near a point of the
relative interior of a cell of minimal remaining dimension, use nonnegative
cutoffs supported in a shrinking transverse tube, normalized by the lattice
volume of that cell.  Contributions from higher-dimensional cells tend to
zero, exactly as in Step~1 of \Cref{prop:ch18-current-faithful}.  Positivity
of the current therefore says that $A_\sigma$ pairs nonnegatively with every
test form in the dual cone.  The finite-dimensional duality of
\Cref{prop:ch09-positivity-dictionary} puts $A_\sigma$ in the required cone.
Remove those cells and continue by induction.  This proves the converse for
each of the three cones in every bidegree; compare
\cite[Remark~1.11]{GK19}.

\cref{itm:ch18-pos:global} If $\beta$ is positive, every algebraic pullback
is positive by definition, and \cref{itm:ch18-pos:current} makes its
associated current positive.  Conversely, assume all the currents
$[f^*\beta]_D$ are positive.  On any compatible tropical chart, their
ambient polyhedral representatives are positive by
\cref{itm:ch18-pos:current}.  Since this is true after every algebraic
pullback, \Cref{def:ch18-positive-delta} makes the original chart class
positive.  Locality gives positivity on $W$.  The weak and strong cases are
identical.  This is \cite[Proposition~2.13]{GK19}.

\cref{itm:ch18-pos:top} A codimension-$l$ term of total bidegree $(n,n)$ has
a coefficient of top bidegree $(n-l,n-l)$ on its $(n-l)$-dimensional cell.
Its three positivity cones are therefore the same half-line, by
\Cref{prop:ch09-positivity-dictionary}.  The fixed-ambient form--current
equivalence in \cref{itm:ch18-pos:current} identifies this condition with
positivity of the associated measure.  The analytic assertions follow from
\Cref{prop:ch18-delta-integration}\,\cref{itm:ch18-di:measure} and the
all-pullback criterion \cref{itm:ch18-pos:global}.

\end{proof}

\begin{warning}{Three positivity words, three theorems}
Positivity of a generalized delta-form, positivity of one current on the
base, and semipositivity of a metric must not be conflated.  On a fixed
ambient polyhedral space the form--current criterion is an equivalence in
every bidegree.  On an analytic variety, the correct converse requires the
currents of \emph{all} algebraic pullbacks, as in
\Cref{lem:ch18-positivity-consistent}\,\cref{itm:ch18-pos:global}; positivity
of the single base current is only a necessary condition.  The relation
with semipositivity of a metric is the content of
\Cref{thm:gk-tropical-frame-positivity} and
\Cref{thm:gk-formal-positivity-comparison}, which require the metric to be
piecewise smooth, respectively piecewise linear.  Nothing here applies to an
arbitrary continuous psh metric.
\end{warning}

\subsection{The two conditions}
\label{sec:ch18-two-conditions}

For a piecewise smooth local weight \(\varphi\) on a tropical chart, formula
\eqref{eq:tropical-delta-pl} separates two necessary conditions:
\begin{enumerate}
\item the restriction of \(\varphi\) to every maximal face is convex;
\item the corner locus \(\varphi\cdot\operatorname{Trop}(U)\) is effective.
\end{enumerate}

Neither condition implies the other, and both are needed.  The following
four examples, all on one-dimensional tropical spaces where every number can
be written down, separate the possibilities.

\begin{example}[The four cases]
\label{ex:ch18-positivity-four-cases}
Work on the standard tropical line $L$ of \eqref{eq:ch18-gauge-rays}, with
unit weights, and write a function by its three outgoing slope data at the
origin and its behaviour on the rays.
\begin{enumerate}
\item\label{itm:ch18-4c:both}
 $\varphi=g|_L$, the gauge of \Cref{ex:ch18-gauge-atom}: affine on each ray,
 hence facewise convex, and corner weight $3>0$.  Both conditions hold, and
 $\varphi$ is psh, as verified there by explicit smooth approximation.
\item\label{itm:ch18-4c:cornerfails}
 $\varphi=-g|_L$: affine on each ray, so facewise convex, but corner weight
 $-3<0$.  The second condition fails; and indeed
 $\dprimeddprime[\varphi]=-3\delta_0$ is not positive.
\item\label{itm:ch18-4c:facefails}
 $\varphi(z)=-|z|^2$ in the ambient Euclidean norm, restricted to $L$: the
 corner weight at the origin is $0$, since the outgoing slopes all vanish,
 so the second condition holds; but $\varphi$ is concave, not convex, on
 each ray, so the first fails.  Here
 $\dprimeddprime[\varphi]=-2\sum_i\rd s_i$ is absolutely continuous and
 negative.
\item\label{itm:ch18-4c:mixed}
 $\varphi=g|_L-\epsilon|z|^2$ for small $\epsilon>0$: the corner weight is
 still $3>0$, but the facewise Hessian is $-2\epsilon<0$ on each ray, so
 $\dprimeddprime[\varphi]=3\delta_0-2\epsilon\sum_i\rd s_i$ is a signed
 measure which is neither positive nor negative.  Positivity of the total
 mass on a large ball, namely $3-2\epsilon\cdot3R$, is not enough:
 positivity of a delta-form is a cellwise condition, and this example shows
 why it must be.
\end{enumerate}
\end{example}

\begin{example}[A nonsmooth pluriharmonic function]
\label{ex:ch18-cross-positivity}
On the cross-fan $\Sigma$ of \eqref{eq:ch16-cross-fan} with unit weights,
the function $h=|x_1|-|x_2|$ of \eqref{eq:ch16-cross-functions} has slopes
$(1,1,-1,-1)$ and corner weight $1+1-1-1=0$.  Thus $h$ satisfies the second
condition with equality, and it is affine on each ray, so it satisfies the
first.  Therefore \cite[Theorem~8.4]{GK19} says that both $h$ and $-h$ are
psh, and indeed $\dprimeddprime[h]=0$.  The example is not a failure of the
psh criterion: it is a nonsmooth pluriharmonic function.  It reappears in
\Cref{sec:harmonic-trop} as the basic example showing that pluriharmonic need
not mean smooth on a tropical space.
\end{example}

\subsection{The tropical-frame criterion}
\label{sec:ch18-frame-criterion}

\begin{theorem}\label{thm:gk-tropical-frame-positivity}
Let \(L\) be a line bundle on an algebraic variety and let its metric be
piecewise smooth over an analytic open set.  The metric is psh in the
current sense if and only if both preceding conditions hold in every
tropical frame.
\end{theorem}
\begin{proof}
In a fixed tropical frame, \eqref{eq:tropical-delta-pl} decomposes the first
Chern delta-current as the sum of the facewise Hessian and the corner-cycle
current.  If the weight is convex on each facet, its Hessian is a positive
$(1,1)$-form there.  If the corner cycle is effective, its integration
current is positive.  Their sum is therefore positive.

Conversely, suppose that the Chern current is positive.  Test it by forms
supported in the relative interior of one facet.  The corner term vanishes,
so positivity of the current implies positivity of the Hessian; hence the
weight is convex on that facet.  Next choose a nonnegative test form supported
in a thin tubular neighborhood of a generic point of a codimension-one wall
and constant in the normal direction.  After shrinking the tube, the
facewise Hessian contribution tends to zero, while the current pairing tends
to the wall weight times a positive tangential volume.  Positivity forces
that wall weight to be nonnegative.  Repeating this on every wall proves
effectivity of the corner locus.  Common refinements show that the two
conditions are independent of the tropical frame.  This is
\cite[Theorem~8.4]{GK19}.
\end{proof}

The tube argument in the converse deserves to be recorded separately,
because it is the standard device for reading off a wall weight from a
current and it will be used again in \Cref{sec:harmonic-trop}.

\begin{lemma}
\label{lem:ch18-tube}
Let $T$ be a positive $(1,1)$-current on an open subset $\Omega$ of a
tropical space of pure dimension $d$, let $\tau$ be a codimension-one face
of a polyhedral structure on $\Omega$, and let $x\in\RelInt\tau$.  Suppose
that on a neighbourhood of $x$ one has $T=[\Theta]+\delta_W$ with $\Theta$ a
smooth $(1,1)$-form and $W$ a weighted cycle supported on $\tau$ with weight
$w$.  Then $w\geq0$.
\end{lemma}

\begin{proof}
Choose a chart in which $\tau$ is $\{x_d=0\}$ near $x$ and take a strongly
positive test form
$\omega_\epsilon=\chi(x')\psi_\epsilon(x_d)\,\Theta_0$, where $\Theta_0$ is
a fixed strongly positive $(d-1,d-1)$-form of complementary bidegree,
$\chi\geq0$ is a fixed bump function on $\tau$ with
$\int_\tau\chi\,\Theta_0=1$, and $\psi_\epsilon\geq0$ is supported in
$[-\epsilon,\epsilon]$ with $\psi_\epsilon\leq1$ and
$\psi_\epsilon(0)=1$.  Then $\langle T,\omega_\epsilon\rangle\geq0$ by
positivity.  On the other hand
\[
 \langle[\Theta],\omega_\epsilon\rangle
 =O\bigl(\|\Theta\|_\infty\,\epsilon\bigr)\to0,
 \qquad
 \langle\delta_W,\omega_\epsilon\rangle=w\int_\tau\chi\,\Theta_0=w ,
\]
the first because the integrand is supported in a slab of thickness
$2\epsilon$.  Letting $\epsilon\to0$ gives $w\geq0$.
\end{proof}

\begin{corollary}
\label{cor:ch18-positivity-onedim}
Let $C$ be a one-dimensional balanced tropical space with positive weights
and let $\varphi$ be continuous on $|C|$ and $C^2$ on each edge.  Then
$\dprimeddprime[\varphi]\geq0$ if and only if $\varphi$ is convex on each
edge and $\operatorname{at}_v(\varphi)\geq0$ at every vertex $v$.  If $C$
has a boundary, the condition at a boundary vertex is instead that the
single outgoing slope be nonnegative.
\end{corollary}

\begin{proof}
This is \Cref{lem:ch15-vertex-atom} together with
\Cref{lem:ch18-tube} in dimension one: the edgewise part of
$\dprimeddprime[\varphi]$ is $\varphi''\rd s$ on each edge and its atomic
part at $v$ is $\operatorname{at}_v(\varphi)\delta_v$.  A measure which is
a sum of an absolutely continuous part and atoms is nonnegative if and only
if both parts are.  The boundary statement is the same computation with only
one term in \eqref{eq:ch18-atom}.
\end{proof}

\begin{example}[The criterion on the running examples]
\label{ex:ch18-frame-criterion-numbers}
On the tropical line $L$: $g|_L$ passes both tests
(\Cref{ex:ch18-positivity-four-cases}\,\cref{itm:ch18-4c:both}) and $-g|_L$
fails the second.  On the chain $\Gamma=[0,2]$ of
\Cref{ex:ch18-chain-delta}, the delta-form
$c_1^\delta(\overline L)$ is effective if and only if
\[
 1+a_1-a_0\geq0,\qquad a_0+a_2-2a_1\geq0,\qquad 1+a_1-a_2\geq0 ,
\]
which is exactly the condition that the three degrees
$\deg(\mathcal L|_{C_0})$, $\deg(\mathcal L|_E)$, $\deg(\mathcal L|_{C_2})$
be nonnegative, that is, that $\mathcal L$ be nef on the special fibre.  For
$a=(0,1,0)$ the middle inequality fails, giving $-2$; for $a=(0,0,0)$ all
three hold, with values $1,0,1$; for $a=(0,-1,0)$ they hold with values
$0,2,0$.  The middle case $a=(0,0,0)$ shows that a zero weight is allowed:
effectivity, not strict positivity, is the criterion.
\end{example}

\subsection{Piecewise linear metrics: the functorial comparison}
\label{sec:ch18-pl-comparison}

For piecewise linear metrics the comparison is global and functorial.

\begin{theorem}\label{thm:gk-formal-positivity-comparison}
Let \(K\) be algebraically closed, complete, and nontrivially valued.  Let
\(X\) be an algebraic variety, \(L\) a line bundle, \(W\subset X^{\mathrm{an}}\)
open, and let \(\|\cdot\|\) be a piecewise linear metric on \(L|_W\).  The
following are equivalent:
\begin{enumerate}
\item the metric is semipositive over \(W\), meaning that its residual line
      bundles are nef there;
\item the metric is functorially delta-psh;
\item the metric is functorially psh;
\item \(c_1^\delta(L|_W,\|\cdot\|)\) is a positive delta-form;
\item for every closed algebraic curve \(C\subset X\), the metric restricted
      to \(W\cap C^{\mathrm{an}}\) is psh.
\end{enumerate}
\end{theorem}
\begin{proof}
We give the argument of \cite[Theorem~6.10]{GK19}, including the two
reductions that are needed for the positivity implication.  Nagata
compactification first permits us to assume that $X$ is proper.  All five
conditions are local for the analytic topology, so, after shrinking around
an arbitrary point of $W$, the PL metric extends to the formal metric defined
by a line bundle $\mathfrak L$ on a formal model $\mathfrak X$ of $X$.  It is
therefore enough to work with this formal model.

Assume (1).  Semipositivity of formal metrics is preserved by algebraic
pullback.  To prove (2), it consequently suffices to prove that the current
associated to $c_1^\delta(L|_W,\|\cdot\|)$ is positive.  The projection formula
allows us to test this after a generically finite projective covering; by
Chow's lemma we may thus suppose that $X$ is projective.

Fix $x\in W$.  There is a compact strictly analytic neighbourhood
$V\Subset W$ such that $(V,L^{\An}|_V)$ has a formal model
$(\mathfrak V,\mathfrak M)$ induced by $(\mathfrak X,\mathfrak L)$.  Write
$Z$ for the closure in $\mathfrak V_s$ of the reduction of $x$.  Because
$x$ is an inner point of $V$, the variety $Z$ is proper over the residue
field.  Semipositivity at $x$ says precisely that $\mathfrak M|_Z$ is nef.
After passing to a dominating formal model, Raynaud's theorem and formal
GAGA let us arrange that the model is induced from a projective flat model
and that $\mathfrak L$ is algebraic.  Choose a very ample model line bundle
$\mathfrak H$, write $H$ for its generic fibre, and let
$\|\cdot\|_{\mathfrak H}$ be its semipositive algebraic metric.  For every
rational $\varepsilon>0$, set
\[
 \mathfrak L_\varepsilon
   :=\mathfrak L\otimes\mathfrak H^{\varepsilon},
 \qquad
 \|\cdot\|_\varepsilon
   :=\|\cdot\|\otimes\|\cdot\|_{\mathfrak H}^{\varepsilon}.
\]
The restriction $\mathfrak L_\varepsilon|_Z$ is ample.  The local
representability criterion \cite[Lemma~6.9]{GK19} therefore shows that
$c_1^\delta(L|_W,\|\cdot\|_\varepsilon)$ is positively representable on a
neighbourhood of $x$.  As $x$ was arbitrary, its associated delta-current is
positive on $W$.  If $n=\dim X$ and
$\beta\in B_c^{n-1,n-1}(W)$ is positive, then
\begin{align*}
 &\big\langle[c_1^\delta(L|_W,\|\cdot\|)],\beta\big\rangle
 +\varepsilon
  \big\langle[c_1^\delta(H|_W,\|\cdot\|_{\mathfrak H})],\beta\big\rangle \\
 &\hspace{7em}=
 \big\langle[c_1^\delta(L|_W,\|\cdot\|_\varepsilon)],\beta\big\rangle
 \geq0 .
\end{align*}
Letting $\varepsilon\downarrow0$ proves positivity of the first current.
Applying the same argument after every algebraic pullback gives (2).

A positive delta-current is a positive ordinary current
\cite[Remark~2.12]{GK19}, so (2) implies (3).  The functorial positivity
criterion \cite[Proposition~2.13]{GK19} gives (3)$\Leftrightarrow$(4), and
(3) immediately implies (5) by pullback to closed algebraic curves.

It remains to prove (5)$\Rightarrow$(1).  By the algebraization of PL
metrics, we may take $(\mathfrak X,\mathfrak L)$ to be an algebraic model.
Let $Y$ be any closed curve in $\mathfrak X_s$ whose generic point lies in
the reduction of $W$.  Since $X$ is proper, $Y$ is proper over the residue
field.  The curve-lifting theorem \cite[Theorem~4.1]{GK19} supplies a closed
algebraic curve $C\subset X$ such that $Y$ is an irreducible component of the
special fibre of the closure of $C$.  By (5), the Chern current of the
restricted metric on $W\cap C^{\An}$ is positive.  Its extension is the
discrete Chambert--Loir measure of the formal metric on $C$.  The curve
criterion \cite[Lemma~6.8]{GK19} now yields
\(
  \deg_{\mathfrak L}(Y)\geq0.
\)
Thus $\mathfrak L$ is vertically nef over $W$, which is (1).

The argument uses the hypotheses on $K$ precisely in the formal-model and
curve-lifting steps.  It concerns PL metrics only; the approximation result
for formal metrics is the separate \cite[Proposition~7.2]{GK19}.
\end{proof}

\begin{remark}[Where each hypothesis is used]
\label{rem:ch18-gk-positivity-hypotheses}
The five equivalent conditions of
\Cref{thm:gk-formal-positivity-comparison} are not equally robust.
\begin{itemize}
\item Conditions (1)--(5) concern PL, hence formal, data.  For a merely
 piecewise-smooth metric there is no residual line bundle in general; the
 valid replacement is \Cref{thm:gk-tropical-frame-positivity}, namely
 facewise convexity plus effectivity of the corner locus.
\item The implication (1)$\Rightarrow$(2),(3) uses functoriality of residual
 line bundles under algebraic morphisms of varieties.  No functoriality for
 arbitrary analytic morphisms is asserted here.
\item The equivalence with (5) uses the curve-lifting theorem, and this is
 where $K$ algebraically closed and nontrivially valued is essential:
 residual curves are lifted to closed algebraic curves after refinement, and
 both steps can fail over a general base.
\item The restriction to PL metrics enters twice, once to have a formal
 model at all and once to know that the residual data are finite.
\end{itemize}
The corresponding statement for a merely continuous semipositive metric is
\emph{not} proved here and is not true in the same form: continuity of a
metric gives no residual line bundle.  What is available in that generality
is the approximation statement \Cref{thm:cl-continuous-extension} and the
comparison of \Cref{chap:compare-ma}.
\end{remark}

\begin{corollary}
\label{cor:ch18-zhang}
Let $X$ be a proper algebraic variety over $K$ as above and let $L$ carry a
PL metric.  The metric is semipositive in the sense of Zhang, that is, it is
induced by a nef model line bundle after passing to a positive tensor power,
if and only if $c_1^\delta(L,\|\cdot\|)$ is a positive delta-form on all of
$X^{\An}$.
\end{corollary}

\begin{proof}
By \Cref{thm:gm-metric-comparison} a PL metric on a proper variety is a
model metric after passing to a power, so the two sides of the assertion are
statements about a model line bundle $\mathcal L$.  Nefness of
$\mathcal L$ on the special fibre is condition (1) of
\Cref{thm:gk-formal-positivity-comparison} with $W=X^{\An}$, and positivity
of $c_1^\delta$ is condition (4); the theorem gives the equivalence.  The
equivalence with the definition of semipositivity used in
\Cref{sec:metrics-semipositive} is \Cref{thm:model-semipositive-closed}, and
the compatibility with \cite{GM19} is discussed there.
\end{proof}

\begin{example}[The criterion tested on the chain]
\label{ex:ch18-zhang-chain}
For $\overline L$ as in \Cref{ex:ch18-chain-delta}, the delta-form
$c_1^\delta(\overline L)$ is positive exactly when
\Cref{ex:ch18-frame-criterion-numbers} says, and by
\Cref{cor:ch18-zhang} this is exactly semipositivity of the metric.  The
case $a=(0,1,0)$, with a residual degree $-2$, therefore gives a PL metric
which is \emph{not} semipositive although its curvature measure
$2\delta_{v_0}-2\delta_{v_1}+2\delta_{v_2}$ has positive total mass $2$.
Total mass is never a positivity criterion.
\end{example}

\section{Harmonic tropicalizations}
\label{sec:harmonic-trop}

\subsection{The defect that delta-forms do not repair}
\label{sec:ch18-harmonic-motivation}

The next extension addresses a different defect.  On a Berkovich curve,
there are harmonic functions which are not smooth in the original CLD
sheaf.  As a result, the CLD Dolbeault cohomology of a smooth proper curve
can be too large and Poincar\'e duality can fail.  Enlarging coefficients by
delta-currents does not by itself make those harmonic functions into
coordinate functions.

The phenomenon is already visible on a fan, and it is worth seeing it there
before any analytic geometry is introduced.

\begin{example}[A harmonic function which is not smooth]
\label{ex:ch18-harmonic-not-smooth}
Let $\Sigma\subset\mathbf R^2$ be the cross-fan \eqref{eq:ch16-cross-fan},
that is, the four rays $\mathbf R_{\geq0}(\pm e_1)$ and
$\mathbf R_{\geq0}(\pm e_2)$, each with weight $1$; it is balanced because
opposite rays cancel.  Let
\[
 h=|x_1|-|x_2|
\]
be the function of \eqref{eq:ch16-cross-functions}, restricted to $\Sigma$.
Its outgoing slopes at the origin are $1,1,-1,-1$ along
$e_1,-e_1,e_2,-e_2$ respectively, and it is affine on each open ray.  Hence,
by \eqref{eq:delta-graph-laplacian},
\[
 \operatorname{at}_0(h)=1+1-1-1=0,\qquad
 \dprimeddprime[h]=0 :
\]
$h$ is harmonic.  Nevertheless $h$ is not smooth at the origin.  Indeed, by
the computation in \Cref{prop:ch18-smooth-zero-atom}, a smooth function has
slopes $\langle \rd F(0),u\rangle$; along $e_1$ and $-e_1$ these are
negatives of each other, whereas $h$ has slope $1$ along both.  Thus no
ambient smooth $F$ restricts to $h$.

Two lessons.  First, harmonicity is strictly weaker than smoothness plus
$\dprimeddprime=0$: the sheaf of harmonic functions is larger than the sheaf
of \emph{smooth} harmonic functions, which on a fan consists of the affine
functions.  Second, the deficiency cannot be repaired by enlarging
coefficients, because $h$ is a function, that is, an object of bidegree
$(0,0)$, where there are no coefficients to enlarge.  The only remedy is to
allow $h$ itself to be used as a coordinate; that is the definition of a
harmonic tropicalization.
\end{example}

\begin{remark}[Consequences for cohomology, and what we use]
\label{rem:ch18-cohomology}
The cohomological statement in the opening paragraph is quoted, not proved,
in this book.  What \Cref{ex:ch18-harmonic-not-smooth} does prove is the
sheaf-theoretic input: on the local model the inclusion
\[
 \{\text{smooth functions }u\text{ with }\dprimeddprime u=0\}
 \subsetneq
 \{\text{harmonic functions}\},
\]
is strict, by \Cref{prop:ch18-harmonic-basic}\,\cref{itm:ch18-hb:smooth}
applied to $h=|x_1|-|x_2|$, which is not the restriction of an affine
function.  The passage from that strictness to a statement about Dolbeault
cohomology and Poincar\'e duality is carried out in
\cite[\S\S10 and 13]{GJR26}, with the local theory of the smooth complex
supplied by Jell's Poincar\'e lemma \cite{Jell16}.  Nothing in
\Crefrange{chap:cld-forms}{chap:cld-metrized-bundles} depends on this;
Dolbeault cohomology is used nowhere in the pluripotential theory of this
book.  We record the defect because it is the stated motivation of the
construction and because it explains why the enlargement is by charts rather
than by coefficients.
\end{remark}

\subsection{Harmonic PL functions}
\label{sec:ch18-harmonic-pl}

Let \(X\) be a good strict analytic space over a complete nontrivially valued
field.  An \(\mathbf R\)-PL function \(h\) has at each point a residual real
line bundle \(L_h(x)\) on the reduction of the germ.

\begin{definition}[Harmonic PL function]
An integral PL function \(h\) is \emph{harmonic} if locally it is induced by
a formal model of the trivial line bundle whose restriction to the special
fibre is numerically trivial.  An \(\mathbf R\)-harmonic function is a local
real linear combination of integral harmonic functions.
\end{definition}

Equivalently, both \(L_h(x)\) and its dual are nef.  On a rig-smooth analytic
curve this is exactly Thuillier harmonicity
\cite[Proposition~15.7]{GJR26}.

The description ``$L_h(x)$ and its dual are both nef'' is the one to use in
proofs, because it makes every stability property a formal consequence of
the corresponding property of nefness.  We record the ones the book needs.

\begin{proposition}
\label{prop:ch18-harmonic-basic}
Let $X$ be a good strict $K$-analytic space.
\begin{enumerate}
\item\label{itm:ch18-hb:sheaf}
 The $\mathbf R$-harmonic functions form a sheaf $\cH_X$ of
 $\mathbf R$-vector spaces, containing the constants.
\item\label{itm:ch18-hb:units}
 If $f\in\cO(U)^\times$ then $-\log|f|$ is harmonic on $U$; more generally
 $-\log|f|$ is harmonic wherever $f$ is invertible.
\item\label{itm:ch18-hb:smooth}
 A smooth function $u$ on a tropical chart satisfies $\dprimeddprime[u]=0$
 if and only if it is locally the restriction of an ambient affine function,
 that is, if and only if locally
 $u=c+\sum_i\lambda_i\bigl(-\log|f_i|\bigr)$ with $c,\lambda_i\in\mathbf R$
 and $f_i$ the invertible functions defining the chart.  Such a $u$ is
 harmonic.
\item\label{itm:ch18-hb:curve}
 On a rig-smooth $K$-analytic curve, an $\mathbf R$-PL function is harmonic
 if and only if the sum of its outgoing slopes vanishes at every point of
 the skeleton, that is, if and only if $\dprimeddprime[h]=0$.
\end{enumerate}
\end{proposition}

\begin{proof}
\cref{itm:ch18-hb:sheaf} Harmonicity is a local condition by definition, so
$\cH_X$ is a subsheaf of the sheaf of $\mathbf R$-PL functions.  It is a
vector subspace because the residual line bundle of $h+h'$ is
$L_h\otimes L_{h'}$ and of $\lambda h$ is $L_h^{\otimes\lambda}$ in
$\Pic\otimes\mathbf R$, and because nefness of both a class and its dual is
preserved under sums and real scalars: numerical triviality is a linear
condition.  A constant $c$ has $L_c$ trivial, hence numerically trivial.

\cref{itm:ch18-hb:units} If $f$ is invertible on $U$ then the trivial line
bundle with the frame $f$ has norm one, so the corresponding formal model is
the trivial one, and the residual bundle is trivial, hence numerically
trivial in both directions.  Equivalently, one may quote
\Cref{lem:ch15-log-smooth}: $-\log|f|$ is the pullback of an affine
coordinate, so $\dprimeddprime$ of it vanishes, and the model is trivial.

\cref{itm:ch18-hb:smooth} If $u=c+\sum_i\lambda_i(-\log|f_i|)$ then
$\dprimeddprime[u]=0$ by \Cref{lem:ch15-log-smooth} and linearity, and $u$ is
harmonic by \cref{itm:ch18-hb:sheaf} and \cref{itm:ch18-hb:units}.
Conversely, let $u=F|_{|C|}$ with $F$ smooth on an ambient open set and
$\dprimeddprime[u]=0$.  Fix $x\in|C|$ and work on the star of $x$, so that
every cell of the local polyhedral structure contains $x$.  Testing
$\dprimeddprime[u]=0$ against forms supported in the relative interior of a
maximal cell $\sigma$ gives, by
\Cref{prop:tropical-delta-pl}, that the facewise Hessian vanishes, that is,
$u|_\sigma$ is affine and $\langle \rd F(y),v\rangle$ is independent of
$y\in\sigma$ for every $v\in N_{\sigma,\mathbf R}$.  Put
$\ell(y)=F(x)+\langle \rd F(x),y-x\rangle$.  For $y$ in a maximal cell
$\sigma\ni x$ we get $u(y)=u(x)+\langle \rd F(x),y-x\rangle=\ell(y)$, and
these cells cover the star.  Hence $u=\ell|_{|C|}$ near $x$, and writing
$\ell$ in the moment coordinates of the chart gives the displayed form.
Note that \Cref{prop:ch18-smooth-zero-atom} makes the atoms of $u$ vanish
automatically, so the hypothesis $\dprimeddprime[u]=0$ carries no
information beyond the facewise one; compare
\Cref{prop:ch16-cell-rigidity}.

\cref{itm:ch18-hb:curve} On a curve, the reduction of a germ at a point of
the skeleton is a curve over the residue field, and the residual line bundle
of $h$ has degree equal to the sum of the outgoing slopes of $h$ at that
point, with the multiplicities of the residual components; this is
\cite[Proposition~15.7]{GJR26} and, in the model language,
\Cref{prop:skeletal-curvature-curve}.  Numerical triviality of a line bundle
on a curve is vanishing of its degree on every component.  The identification
of that sum with $\dprimeddprime[h]$ is \eqref{eq:delta-graph-laplacian}.
\end{proof}

The property that makes the theory usable is functoriality.  It is what
allows a harmonic tropicalization to be pulled back along a morphism and
hence to define a sheaf of forms.

\begin{proposition}
\label{prop:ch18-harmonic-functorial}
Let $\varphi:Y\to X$ be a morphism of good strict $K$-analytic spaces and
let $h$ be an $\mathbf R$-harmonic function on an open $U\subseteq X$.
\begin{enumerate}
\item\label{itm:ch18-hf:pullback}
 $\varphi^*h=h\circ\varphi$ is $\mathbf R$-harmonic on $\varphi^{-1}(U)$.
\item\label{itm:ch18-hf:base}
 If $K'/K$ is a complete valued extension, then the pullback of $h$ to
 $U\widehat\otimes_KK'$ is $\mathbf R$-harmonic.
\item\label{itm:ch18-hf:domain}
 The restriction of $h$ to a strict analytic domain of $U$ is
 $\mathbf R$-harmonic.
\item\label{itm:ch18-hf:compose}
 If $h_1,\ldots,h_r$ are integral harmonic on $U$ and
 $\lambda\in\mathbf R^r$, then $\sum\lambda_ih_i$ is $\mathbf R$-harmonic;
 consequently, if $h=(h_1,\ldots,h_r):U\to\mathbf R^r$ is a harmonic
 tropicalization and $A:\mathbf R^r\to\mathbf R^{r'}$ is an integral affine
 map, then $A\circ h$ is a harmonic tropicalization.
\end{enumerate}
\end{proposition}

\begin{proof}
\cref{itm:ch18-hf:pullback} Harmonicity is local, so we may assume that $h$
is integral and is induced by a formal model $\cL$ of $\cO_U$ whose
restriction to the special fibre is numerically trivial.  Pulling back the
formal model along a formal model of $\varphi$---which exists locally after
an admissible formal blowup, by the flattening and domination results of
\Cref{chap:models}---produces a formal model of $\cO_{\varphi^{-1}(U)}$
inducing $\varphi^*h$.  Its restriction to the special fibre is the pullback
of a numerically trivial class, hence numerically trivial, because degrees on
residual curves are computed by the projection formula: for a residual curve
$C'$ mapping to a residual curve $C$ with degree $e$, one has
$\deg(\varphi^*\cL|_{C'})=e\deg(\cL|_C)=0$, and for a contracted $C'$ the
degree is $0$.  Real coefficients are handled by linearity.

\cref{itm:ch18-hf:base} Scalar extension of a formal model is a formal model,
and numerical triviality is preserved because the reduction of the extension
is the base change of the reduction along a field extension of the residue
field, under which degrees of line bundles on curves are multiplied by the
residue-degree factor of \Cref{cor:trop-finite-field-extension} and in
particular remain zero.

\cref{itm:ch18-hf:domain} A special case of
\cref{itm:ch18-hf:pullback} for the inclusion of a domain.

\cref{itm:ch18-hf:compose} The first assertion is
\Cref{prop:ch18-harmonic-basic}\,\cref{itm:ch18-hb:sheaf}.  For the second,
the coordinates of $A\circ h$ are integral affine combinations of the $h_i$,
which are integral harmonic by \cref{itm:ch18-hb:sheaf} applied with integer
coefficients, the constant term being harmonic as well.
\end{proof}

The integral metric criterion of Chambert--Loir--Ducros and the real residue
formalism of Gubler--Jell--Rabinoff fit together without pulling currents
back to curves.  The following result uses an ample residual class as an
order unit for the finitely many integral PL summands of a real PL function.

\begin{theorem}
\label{thm:p01-real-pl-positive-current-candidate}
Let $K$ be a complete nontrivially valued non-Archimedean field, let $X$
be a good, separated, boundaryless, pure-dimensional, strictly
$K$-analytic space, and let
$h:X\to\mathbf R$ be an $\mathbf R$-PL function.  The following conditions
are equivalent:
\begin{enumerate}
\item the Chambert--Loir--Ducros current
      $\dprimeddprime[h]$ is positive;
\item the residue class
      \[
      L_h(x)\in\operatorname{Pic}(\operatorname{red}(X,x))
                    \otimes_{\mathbf Z}\mathbf R
      \]
      is nef for every $x\in X$;
\item $h$ is semipositive in the sense of Gubler--Jell--Rabinoff.
\end{enumerate}
Consequently, if the same $X$ is allowed to have boundary, positivity of
$\dprimeddprime[h]$ on $\operatorname{Int}(X)$ is equivalent to
semipositivity of $h|_{\operatorname{Int}(X)}$.
\end{theorem}

\begin{proof}
The equivalence of (2) and (3) is
\cite[Proposition~7.1]{GJR26}, and (3) implies (1) by
\cite[Theorem~7.14]{GJR26}.  We prove that (1) implies (2).

Fix $x\in X$.  By \cite[Proposition~5.8 and Lemma~5.10]{GJR26}, after
shrinking to an analytic neighbourhood of $x$, write
\begin{equation}
\label{eq:p01-real-pl-decomposition}
                         h=\sum_{i=1}^{s}\lambda_i h_i,
\end{equation}
where the $h_i$ are integral PL functions and
$\lambda_i\in\mathbf R$.

We first construct one integral PL function that dominates the finitely
many $h_i$ in both current directions.  Choose a compact strictly analytic
neighbourhood $Y$ of $x$, with $x\in\operatorname{Int}(Y)$, on which all
$h_i$ are defined.  By \cite[Corollary~5.9 and \S6.16]{GJR26}, after passing
to a formal model dominating the finitely many models involved, the closure
$V_0$ of the reduction of $x$ in the special fibre is a model of
$\operatorname{red}(X,x)$, and line bundles $M_{i,0}$ on $V_0$ represent
$L_{h_i}(x)$.  Since $x$ is inner, $V_0$ is proper; this is the inner-point
input cited in the proof of \cite[Proposition~7.10]{GJR26}.  As in that
proof, Chow's lemma gives a proper birational morphism
$\pi:V\to V_0$ with $V$ projective.  By \cite[\S6.2]{GJR26}, properness of
$\pi$ shows that $V$ is again a model of $\operatorname{red}(X,x)$.  Put
$M_i=\pi^*M_{i,0}$ and choose an ample line bundle $A$ on $V$.  The
realization step \cite[\S6.6.5]{CLD12}, equivalently the strict lifting
statement \cite[Theorem~17.5.7(b)]{CLD25}, gives an integral PL germ $a$
with $(V,A)$ a model of $L_a(x)$.  Shrinking once, regard $a$ and all the
$h_i$ as functions on the same neighbourhood.

For each $i$, the line bundles
\[
                         A^{\otimes N}\otimes M_i,
             \qquad     A^{\otimes N}\otimes M_i^{-1}
\]
are ample for all sufficiently large $N$.  Indeed, apply the ample-twist
criterion \cite[Tag~0892]{stacks-project} to the identity morphism of $V$,
the line bundles $M_i^{\pm1}$, and $A$.  Since the family is finite, fix one
positive integer $N$ that works for every $i$.  Thus
\begin{equation}
\label{eq:p01-ample-domination-residues}
                       NA+M_i\quad\text{and}\quad NA-M_i
\end{equation}
are nef on $V$.  By residue linearity \cite[Lemma~6.18]{GJR26} and the
residue criterion \cite[Proposition~7.1]{GJR26}, $a$ and each $Na\pm h_i$
are semipositive at $x$.  The semipositive locus is open by
\cite[Remark~7.3(4)]{GJR26}; after intersecting the finitely many such loci,
all these functions are semipositive on one open neighbourhood $U$ of $x$.
Consequently, \cite[Theorem~7.14]{GJR26} gives positive currents on $U$
\begin{equation}
\label{eq:p01-current-domination}
 T_a:=\dprimeddprime[a]\geq0,
 \qquad NT_a+T_i\geq0,
 \qquad NT_a-T_i\geq0,
 \qquad T_i:=\dprimeddprime[h_i].
\end{equation}

Set $\varepsilon_m=1/m$.  Choose $q_i^{(m)}\in\mathbf Q$ so close to
$\lambda_i$ that, with $r_i^{(m)}=q_i^{(m)}-\lambda_i$,
\begin{equation}
\label{eq:p01-rational-error-budget}
                    N\sum_{i=1}^{s}|r_i^{(m)}|\leq\varepsilon_m.
\end{equation}
Define the $\mathbf Q$-PL function
\begin{equation}
\label{eq:p01-rational-positive-approximants}
                       g_m:=\sum_{i=1}^{s}q_i^{(m)}h_i
                                  +\varepsilon_m a.
\end{equation}
For a nonzero real number $r$, write $\sigma(r)=r/|r|$; when $r=0$, omit
the corresponding summand.  Linearity of residues and currents, the
assumed positivity of $T_h:=\dprimeddprime[h]$, and
\eqref{eq:p01-current-domination} give
\begin{align}
\label{eq:p01-positive-decomposition}
 \dprimeddprime[g_m]
  ={}&T_h
    +\left(\varepsilon_m-N\sum_i|r_i^{(m)}|\right)T_a \\
   &+\sum_{r_i^{(m)}\ne0}|r_i^{(m)}|
       \left(NT_a+\sigma(r_i^{(m)})T_i\right)\geq0
\end{align}
on $U$.  Notice that this argument tests positivity against every positive
compactly supported form on $U$; no choice of a tropical chart or skeleton
is involved.

Choose a positive integer $d_m$ such that $d_mg_m$ is integral PL.  Apply
\cite[Corollary~18.8.8]{CLD25} to the trivial line bundle with the integral
PL metric $\lVert1\rVert=\mathrm e^{-d_mg_m}$.  Its first Chern current is
$\dprimeddprime[d_mg_m]$.  The neighbourhood $U$ is boundaryless,
topologically separated, and pure-dimensional, so the cited corollary
applies.  In view of
\eqref{eq:p01-positive-decomposition}, its residue at $x$ is nef.  Dividing
by $d_m$ shows that
\begin{equation}
\label{eq:p01-nef-approximating-residues}
             L_{g_m}(x)=\sum_iq_i^{(m)}L_{h_i}(x)
                              +\varepsilon_mL_a(x)
\end{equation}
is nef.

On the fixed projective model $V$, the class in
\eqref{eq:p01-nef-approximating-residues} is represented by
$\sum_iq_i^{(m)}M_i+\varepsilon_mA$.  By
\cite[Lemma~6.11]{GJR26}, this representative is nef.  Hence for every
proper integral curve $C\subset V$,
\[
  0\leq
  \sum_iq_i^{(m)}\deg_C(M_i)+\varepsilon_m\deg_C(A).
\]
Here $q_i^{(m)}\to\lambda_i$ by
\eqref{eq:p01-rational-error-budget}, and $\varepsilon_m\to0$.  Taking the
limit yields
\[
                       \sum_i\lambda_i\deg_C(M_i)\geq0.
\]
This holds for every proper integral curve on $V$, so
$\sum_i\lambda_iM_i$ is nef.  It represents
$L_h(x)=\sum_i\lambda_iL_{h_i}(x)$ by
\cite[Lemma~6.18]{GJR26}; hence $L_h(x)$ is nef.  Since $x$ was arbitrary,
(2) holds.

For a pure-dimensional analytic space with boundary, its interior is a
boundaryless pure-dimensional open analytic subspace.  Apply the statement
already proved to $\operatorname{Int}(X)$.  This proves the final assertion.
\end{proof}

This is a book-level deduction from the cited results; no novelty or priority
claim is made.  The argument uses the ample class only as a finite-dimensional
order unit.  It does not assert a pullback theorem for arbitrary positive CLD
currents.

The vanishing case admits a second proof that exposes the rational lattice
calibration along codimension-one faces.  Keeping the analytic cells of a
skeletal current separate, rather than passing only to their possibly
self-overlapping tropical image, recovers precisely that information.

\begin{proposition}
\label{prop:ch18-real-pl-zero-current-harmonic}
Let $K$ be a complete nontrivially valued non-Archimedean field, and let
$X$ be a good, separated, boundaryless, strictly $K$-analytic space that is
pure of dimension $n\geq1$.  If $h:X\to\mathbf R$ is an $\mathbf R$-PL
function and its Chambert--Loir--Ducros current satisfies
\[
                         \dprimeddprime[h]=0,
\]
then $h$ is harmonic in the sense of Gubler--Jell--Rabinoff.  Equivalently,
the residue class
\[
 L_h(x)\in\operatorname{Pic}(\operatorname{red}(X,x))
                  \otimes_{\mathbf Z}\mathbf R
\]
is numerically trivial for every $x\in X$.  Consequently, $h$ is
classically pluriharmonic in the sense of Gubler--Rabinoff.
\end{proposition}

\begin{proof}
Fix $x\in X$.  By \cite[Proposition~5.8 and Lemma~5.10]{GJR26}, after
shrinking to an analytic neighbourhood $W_0$ of $x$, one can write
\begin{equation}
\label{eq:ch18-zero-current-independent}
                         h=\sum_{i=1}^{s}\lambda_i h_i,
\end{equation}
where the $h_i$ are integral PL functions on $W_0$ and the real numbers
$\lambda_1,\ldots,\lambda_s$ are linearly independent over $\mathbf Q$.
By \cite[Proposition~8.3 and \S8.4]{GJR26}, after replacing $W_0$ by a
smaller analytic neighbourhood of $x$, choose a smooth tropicalization
\[
                         F:W_0\longrightarrow\mathbf R^N
\]
simultaneously tropicalizing $h_1,\ldots,h_s$ on $W_0$.  Thus
$h_i=\varphi_i\circ F$, with $\varphi_i$ piecewise
$(\mathbf Z,\Gamma)$-affine.  Choose a compact strictly affinoid
neighbourhood $A\subset W_0$ of $x$ with $x\in\operatorname{Int}(A)$, and
then choose an open neighbourhood $W$ of $x$ with
$W\subset\operatorname{Int}(A)$.  Restrict the $h_i$ and $F$ to $W$.
Since boundarylessness and purity are local here, we may suppose that $W$
is boundaryless and pure of dimension $n$.  We prove that
\begin{equation}
\label{eq:ch18-zero-current-integral}
                         \dprimeddprime[h_i]=0
                         \qquad(1\leq i\leq s)
\end{equation}
on a neighbourhood of $x$.

Let $T_0=[W]$ be the integration current.  By
\cite[Example~7.5.3]{CLD25}, $T_0$ is a closed skeletal current of dimension
$n$, described in each moment map by the canonical tropical calibration.
For each $i$, put
\[
                         T_i:=\dprimeddprime(h_iT_0)
                             =\dprimeddprime[h_i].
\]
By \cite[Theorem~8.10.3]{CLD25}, $T_i$ is a skeletal current of dimension
$n-1$.

We recall the arithmetic in the construction in the proof of that theorem,
using the fixed tropicalization $F$ defined on all of $W$.  Let
$P_F^{T_0}\subset W$ be the minimal
canonically calibrated skeleton describing $T_0$ relative to $F$, as in
\cite[Propositions~7.5.7--7.5.8]{CLD25}.  In particular, $F$ is an immersion
on each cell of $P_F^{T_0}$; distinct analytic cells are not identified in
this description, even if their images overlap in $F(W)$.  Choose a common
polyhedral subdivision $\mathscr C$ of $P_F^{T_0}$ on which all the
$\varphi_i\circ F$ are affine.  A top-dimensional cell $C$ carries its
canonical calibration.  By
\cite[Proposition~14.3.1 and \S14.3.2]{CLD25}, it is rational with respect to
the integral tangent lattice, so it has the form
\[
                         \mu_C=q_C\mu_C^{\mathrm{prim}},
                         \qquad q_C\in\mathbf Q_{>0},
\]
where $\mu_C^{\mathrm{prim}}$ is the primitive lattice volume transported
through $F|_C$.  If $D$ is a codimension-one cell and $\nu_D$ is its
primitive lattice volume, the formula in the proof of
\cite[Theorem~8.10.3]{CLD25} describes $T_i$ along $D$ by
\begin{equation}
\label{eq:ch18-zero-current-calibration}
                         \mu_{i,D}=a_{i,D}\nu_D,
                         \qquad a_{i,D}\in\mathbf Q.
\end{equation}
Indeed, the analytic degree entering the canonical tropical multiplicity is
an integer, while its comparison with the primitive tangent lattice is
rational by the cited propositions.  The derivative of the integral-affine
function $\varphi_i\circ F|_C$ in a primitive normal direction is an integer,
and every lattice transition factor is rational.  With the orientation sign
used in the boundary integral, the contribution from each cell $C$ adjacent
to $D$ is therefore rational.  Summing the finitely many contributions gives
$a_{i,D}\in\mathbf Q$.  This is the coefficient of the calibrated tropical
divisor of $\varphi_i\circ F$ on the weighted skeleton.  A further common
refinement subdivides these calibrated cycles and preserves rationality.

By linearity and \eqref{eq:ch18-zero-current-independent}, the skeletal
current $\sum_i\lambda_iT_i=\dprimeddprime[h]$ is described on the same
subdivision along $D$ by
\begin{equation}
\label{eq:ch18-zero-current-sum-calibration}
                         \left(\sum_i\lambda_i a_{i,D}\right)\nu_D.
\end{equation}
The assumed current identity says that this skeletal current is zero.  Let
$|\mathscr Q|\subset W$ denote the source-side support of the resulting
calibrated codimension-one skeleton.  Consequently the calibration in
\eqref{eq:ch18-zero-current-sum-calibration} vanishes on every
codimension-one cell.  Indeed, suppose that it were nonzero on the relative
interior of some $D$.  Choose a compact polyhedral patch
$K_D\Subset D^\circ$ with nonempty relative interior.  In the source-side
skeleton, choose an ambient open neighbourhood $V\subset W$ of $K_D$ such
that
\[
                         V\cap|\mathscr Q|\subset D^\circ.
\]
This is possible because the source-side subdivision is locally finite and
$K_D$ is contained in the interior of its cell.  In particular, overlapping
target images create no additional source-side incidence.  By
\cite[Corollary~6.2.4]{CLD25}, there is a nonnegative smooth function $\chi$
which is equal to $1$ on a neighbourhood of $K_D$ and satisfies
$\operatorname{Supp}(\chi)\subset V$.
Choose a constant Lagerberg $(n-1,n-1)$-form whose evaluation on $\nu_D$ is
nonzero, and change its sign if necessary.  Pairing the current with the
test form obtained by multiplying its pullback by $\chi$, the source-side
formula of \cite[Definition~7.5.1]{CLD25} gives a strictly positive value, a
contradiction.  Notice that distinct analytic cells remain distinct in this
formula even if their images under $F$ overlap.  This source-side separation
is part of the skeletal-current description and its uniqueness in
\cite[Definition~7.5.1 and Propositions~7.5.7--7.5.8]{CLD25}.  Thus
\[
                         \sum_i\lambda_i a_{i,D}=0.
\]
Since $a_{i,D}\in\mathbf Q$ and the $\lambda_i$ are $\mathbf Q$-linearly
independent, $a_{i,D}=0$ for every $i$ and $D$.  Hence the calibrated
skeletal description of each $T_i$ is zero in the fixed chart $F$ and its
polyhedral refinements.

It remains to verify the full current identity, not merely a fixed-chart
pushforward.  Let $\alpha$ be any smooth compactly supported
$(n-1,n-1)$-test form in $W$.  By the compact-support localized partition of
unity of \cite[Corollary~6.2.8]{CLD25}, together with the determination of
smooth forms by localized terms in \cite[Lemma~7.2.5]{CLD25}, it is a
finite sum of terms $\rho\omega$, where $\rho$ is a smooth compactly
supported function and $\omega$ is tropicalized by a moment map $H$ on an
open set $V\subset W$ containing
$\operatorname{Supp}(\rho)$.  Since the fixed $F$ is defined on all of $W$,
the common refinement $(F|_V,H)$ is defined on $V$ and still simultaneously
tropicalizes $h_1,\ldots,h_s$.  The same computation on a common subdivision
again has rational coefficients.  Since the calibrated coefficients of
$\sum_i\lambda_iT_i=0$ vanish and the $\lambda_i$ are $\mathbf Q$-linearly
independent, the coefficient of every $T_i$ on every source-side cell of
this common refinement is zero.  The defining identity in
\cite[Definition~7.5.1]{CLD25} then gives
\[
                         \langle T_i,\rho\omega\rangle=0.
\]
Summing the localized terms proves \eqref{eq:ch18-zero-current-integral}.
Notice that no analytic current is pulled back and no target representative
is required to avoid the image of the boundary of a compact analytic
domain.  Compact support is carried by $\rho$ on the open boundaryless
space, exactly as in the definition of a skeletal current.

For each $i$, apply \cite[Theorem~18.8.4]{CLD25} with $d=1$ to the trivial
line bundle with the integral PL metric determined by $h_i$.  The separated
and boundaryless hypotheses are satisfied, and
\eqref{eq:ch18-zero-current-integral} implies that $L_{h_i}(x)$ is
numerically trivial.  The residue construction is linear, so
\[
                         L_h(x)=\sum_i\lambda_iL_{h_i}(x)
\]
is numerically trivial.  Since $x$ was arbitrary,
\cite[Proposition~7.1 and Definition~7.2]{GJR26} imply that $h$ is harmonic.

Finally, harmonicity means that $h$ and $-h$ are semipositive.  The field is
nontrivially valued, $X$ is good and strict, and both functions are
$\mathbf R$-PL.  Hence \cite[Theorem~4.3]{GR25Classical} makes both functions
classically psh, so $h$ is classically pluriharmonic.
\end{proof}

The proof of \Cref{prop:ch18-real-pl-zero-current-harmonic} is a book-level
consequence of the cited CLD, Gubler--Jell--Rabinoff, and
Gubler--Rabinoff results; no novelty or priority claim is made.  The
rational calibrated coefficients and the refinement by the chart of an
arbitrary test form are essential.  In particular, the proof neither pulls
back a CLD current nor uses a target-side cutoff to avoid the image of the
boundary.

\begin{warning}{The continuous zero-current problem}
\Cref{prop:ch18-real-pl-zero-current-harmonic} settles the implication for
$\mathbf R$-PL functions under its hypotheses.  For an arbitrary continuous
function $u$ in higher dimension, however, it is not known in this
generality whether $\dprimeddprime[u]=0$ forces harmonicity or classical
pluriharmonicity: there is no residual line bundle or integral skeletal
calibration from which to recover the face-balancing law, and arbitrary CLD
currents cannot simply be pulled back to all curve probes.  CLD further note
in \cite[Remark~7.8.3]{CLD25} that it is not even known whether local
integrability plus $\dprimeddprime[u]=0$ forces continuity.  The published
continuous result on smooth proper algebraic curves is recorded in
\Cref{cor:ch19-continuous-zero-current-harmonic}; see also
\Cref{def:ch16-harmonic-function,sec:ch16-harmonic-ambiguity}.
\end{warning}

\subsection{Harmonic tropicalizations and their weights}
\label{sec:ch18-harmonic-trop-def}

\begin{definition}[Harmonic tropicalization]
On a compact strict analytic domain \(W\), a map
\[
                              h:W\to\mathbf R^r
\]
is a \emph{harmonic tropicalization} if all its coordinate functions are
integral harmonic.  Smooth tropicalizations by invertible analytic functions
are special cases, since \(-\log|f|\) is harmonic when \(f\) is invertible.
\end{definition}

The image is a finite polyhedral set.  It carries canonical positive weights
obtained locally from ordinary tropicalizations and then descended through
the graph of \(h\).

We describe the weights explicitly, since they are what balancing is about.
Cover $W$ by finitely many affinoid domains $W_j$ on each of which an
ordinary moment map $g_j$ is defined, and consider the graph map
$(g_j,h):W_j\to\mathbf R^{m_j}\times\mathbf R^r$.  Its image is a
polyhedral set with the canonical calibration of
\Cref{def:trop-canonical-calibration}, by
\Cref{thm:trop-calibration-canonical}.  Let $p:\mathbf R^{m_j}\times
\mathbf R^r\to\mathbf R^r$ be the second projection.  For a $d$-dimensional
face $\tau$ of $h(W_j)$ set
\begin{equation}
\label{eq:ch18-harmonic-weight}
 m_\tau\bigl(h(W_j)\bigr)
 =\sum_{\substack{\sigma\ \text{face of }(g_j,h)(W_j)\\ p(\sigma)=\tau,\
 \dim\sigma=d}}
 m_\sigma\,\bigl[\Lambda_\tau:p(\Lambda_\sigma)\bigr] ,
\end{equation}
the pushforward weight of \Cref{thm:ch11-pushforward-balanced}, where
$\Lambda$ denotes the tangent lattice.  Faces $\sigma$ of dimension $d$
collapsing to a face of smaller dimension contribute nothing, by
\Cref{lem:ch11-cell-change-of-variables}\,\cref{itm:ch11-cov:collapse}.  The
resulting weights are nonnegative because the canonical calibration is, by
\Cref{cor:trop-weights-integral}; and they are independent of $j$ and of
$g_j$ by \Cref{thm:trop-calibration-chart-independent}.

\begin{theorem}\label{thm:harmonic-balancing}
Let \(W\) be a compact good strict analytic space of pure dimension \(d\),
and let \(h:W\to\mathbf R^r\) be a harmonic tropicalization.  The canonical
weights on the \(d\)-dimensional part of \(h(W)\) satisfy the balancing
condition at each codimension-one face not contained in
\(h(S_h(W)_d\cap\partial W)\), where \(S_h(W)_d\) is the \(d\)-skeleton of
the tropicalization.  In particular, balancing holds at every such face not
contained in \(h(\partial W)\).
\end{theorem}
\begin{proof}
Choose an ordinary moment map $g$ on a finite affinoid cover of $W$ and form
the graph tropicalization $(g,h)$.  Its top-dimensional part has the usual
analytic calibration and is balanced away from the analytic boundary by
\Cref{thm:analytic-tropical-balancing}.  Projecting the graph to the
$h$-coordinates defines the canonical weights on $h(W)$ through the
lattice-index pushforward formula.

Fix a codimension-one face $\tau$ of $h(W)$ outside the set named in the
statement and refine so that all graph facets over $\tau$ map linearly.
Apply Stokes to the graph over a small transversal neighborhood of $\tau$.
The contributions from ordinary moment coordinates cancel by analytic
balancing.  For a coordinate $h_j$, the remaining coefficient is the degree
of its residual line bundle on the corresponding residual curve.  That
degree is zero because $h_j$ is harmonic.  Thus, after pushforward, the sum
of the weighted primitive normal vectors at $\tau$ is zero.

If a graph face lies over $S_h(W)_d\cap\partial W$, Stokes leaves its outward
boundary contribution and the preceding cancellation need not hold.  The
excluded image is exactly the set of faces for which this can occur.  This
proves the stated qualified balancing and is the graph-and-Stokes proof of
\cite[Theorem~9.7]{GJR26}.
\end{proof}

\begin{remark}[Why the boundary exclusion is needed]
\label{rem:ch18-balancing-boundary}
The qualification in \Cref{thm:harmonic-balancing} is not a technicality.
Balancing is equivalent to the vanishing of the Stokes defect
(\Cref{thm:ch11-balancing-iff-stokes}), and the defect at a face has two
possible sources: a nonzero residual degree, killed here by harmonicity, and
an outward boundary flux, which is not.  \Cref{ex:ch11-tripod-unbalanced}
exhibits the second phenomenon in the simplest case: a fan whose weights do
not balance because one ray has been removed, which is exactly what
restricting to a domain with boundary does.  In particular a harmonic
tropicalization of a \emph{closed disc} produces an unbalanced image, and
the theory of \Cref{sec:weak-smooth} is stated for spaces without boundary
or with the boundary faces excluded.
\end{remark}

\subsection{Examples of harmonic tropicalizations}
\label{sec:ch18-harmonic-examples}

\begin{example}[Ordinary tropicalizations]
\label{ex:ch18-harmonic-ordinary}
If $f_1,\ldots,f_r\in\cO(W)^\times$ then
$h=(-\log|f_1|,\ldots,-\log|f_r|)$ is a harmonic tropicalization by
\Cref{prop:ch18-harmonic-basic}\,\cref{itm:ch18-hb:units}, and it is the
moment map of \Cref{def:trop-moment-invariant}.  The canonical weights
\eqref{eq:ch18-harmonic-weight} then agree with the tropical multiplicities
of \eqref{eq:trop-multiplicity-formula}, because the graph of $h$ over an
ordinary moment map projects isomorphically onto a refinement.  So the
harmonic theory contains the ordinary one, and
\Cref{thm:harmonic-balancing} contains
\Cref{thm:analytic-tropical-balancing}.
\end{example}

\begin{example}[The cross-fan as a chart]
\label{ex:ch18-cross-chart}
Let $C$ be a rig-smooth $K$-analytic curve whose skeleton near a point $x$
is the cross-fan $\Sigma$ of \Cref{ex:ch18-harmonic-not-smooth}, with four
branches of multiplicity $1$; concretely, take the analytification of a
curve with a node-type degeneration whose special fibre at $x$ has four
residual branches.  The function $h=|x_1|-|x_2|$ on $\Sigma$ pulls back to
an $\mathbf R$-PL function on a neighbourhood of $x$ in $C$, and it is
harmonic by
\Cref{prop:ch18-harmonic-basic}\,\cref{itm:ch18-hb:curve}, since its four
outgoing slopes $1,1,-1,-1$ sum to zero.  It is a harmonic tropicalization
$W\to\mathbf R$ of rank one.  Its image is an interval, and the canonical
weight \eqref{eq:ch18-harmonic-weight} of that interval is
\[
 m=\sum_{\text{branches with slope}>0}1\cdot1=2 ,
\]
the two branches along $e_1$ and $-e_1$ each mapping with lattice index $1$
onto the positive half-line and the other two onto the negative half-line;
so the weight of the positive half is $2$ and of the negative half is $2$.
The image is therefore the line with the constant weight $2$, which is
balanced at $0$: $2\cdot(+1)+2\cdot(-1)=0$, as
\Cref{thm:harmonic-balancing} requires.  By contrast, the same $h$ is not
an ordinary moment map: it is not of the form $-\log|f|$ near $x$, since
that would make it smooth, contradicting
\Cref{ex:ch18-harmonic-not-smooth}.
\end{example}

\begin{example}[Failure of balancing on a disc]
\label{ex:ch18-harmonic-disc}
Let $W=\{|T|\leq1\}$ be the closed unit disc and $h=-\log|T|$, an ordinary
and hence harmonic tropicalization.  Its image is $[0,+\infty]$, of
dimension one, with weight $1$; the face $\{0\}$ has a single adjacent
maximal face, with primitive normal $+1$, so the balancing sum is $1\neq0$.
The point $0$ is $h(\partial W)$, the image of the Shilov boundary, so it is
excluded by \Cref{thm:harmonic-balancing}.  Every interior point of
$(0,+\infty)$ has two adjacent faces after subdividing, with opposite
normals and weight $1$, and balances.
\end{example}

\begin{example}[An abelian variety]
\label{ex:ch18-harmonic-abelian}
On a totally degenerate abelian variety $A^{\An}=(\mathbf G_m^{g})^{\An}/\Lambda$
the retraction to the real torus $\mathbf R^g/\Lambda$ is not given by
global units, since the coordinates $-\log|z_i|$ are only defined modulo
$\Lambda$; but locally they are, and the local harmonic functions patch
into a harmonic tropicalization of $A^{\An}$ onto a fundamental domain,
which is the content of \Cref{thm:abelian-harmonic-tropicalization}.  The
canonical weight is $1$ everywhere and balancing is automatic, the image
being an open subset of $\mathbf R^g$.  This is the example for which the
harmonic theory was designed.  There is a canonical injection
$\Lambda^{p,q}M_{\mathbf R}\to H^{p,q}$ of the invariant classes into the
Dolbeault cohomology of the weakly smooth complex.  The unrestricted
isomorphism question posed in \cite[\S1.15]{GJR26} has a negative answer by
\Cref{prop:ch18-gjr-cohomology-counterexample}.  That counterexample has good
reduction, so it does not decide the sharper totally degenerate case.
\end{example}

\begin{proposition}\label{prop:ch18-gjr-cohomology-counterexample}
The canonical injection of \cite[\S1.15]{GJR26} is not an isomorphism for
all abelian varieties and all bidegrees.  More precisely, put
\[
 K=\widehat{\overline{\mathbf C((t))}}
\]
with the extended $t$-adic absolute value, and let
\[
 E_0/\mathbf C:\qquad y^2z=x^3-xz^2,
 \qquad A=E_0\times_{\mathbf C}K .
\]
Then $A$ has good reduction and, in bidegree $(1,1)$, the canonical
injection is
\begin{equation}
 0=\Lambda^{1,1}M_{\mathbf R}
 \longrightarrow H^{1,1}(A^{\An}),
 \qquad
\dim_{\mathbf R}H^{1,1}(A^{\An})=1.
\label{eq:ch18-gjr-cohomology-counterexample}
\end{equation}
In particular it is not surjective.
\end{proposition}

\begin{proof}
The field $K$ is complete, nontrivially valued, and non-Archimedean by
construction.  The displayed Weierstrass equation has coefficients
$a=-1$, $b=0$, and discriminant
\[
 -16(4a^3+27b^2)=64\ne0.
\]
Thus $E_0$, with origin $[0:1:0]$, is an elliptic curve.  The constants
$\mathbf C\subset K$ lie in the valuation ring $K^\circ$, and
\[
 \mathcal A=E_0\times_{\mathbf C}\operatorname{Spec}K^\circ
\]
is a smooth proper group scheme over $K^\circ$, since smoothness and
properness are preserved by base change.  Its generic fibre is $A$, so
$A$ has good reduction.

For an abelian variety with good reduction, the torus in its Raynaud
uniformization is trivial; see \cite[\S16]{GJR26}.  Hence its cocharacter
lattice $N$ is zero and
\[
 M_{\mathbf R}=\operatorname{Hom}(N,\mathbf R)=0,
 \qquad
 \Lambda^{1,1}M_{\mathbf R}
 =\Lambda^1M_{\mathbf R}\otimes_{\mathbf R}\Lambda^1M_{\mathbf R}=0.
\]
On the other hand, \cite[Theorem~1.14]{GJR26} states that for a smooth
projective curve $Y$, the weakly smooth Dolbeault group
$H^{p,q}(Y^{\An})$ has real dimension one when $p+q\in\{0,2\}$.
Applying this with $Y=A$ and $(p,q)=(1,1)$ gives the second assertion in
\eqref{eq:ch18-gjr-cohomology-counterexample}.  The injection from
\cite[\S1.15]{GJR26} therefore has zero source and a one-dimensional target,
and cannot be surjective.

The toric rank of this example is zero, whereas a totally degenerate elliptic
curve has toric rank one.  Thus the argument answers the literal unrestricted
question but makes no assertion about the totally degenerate refinement.
\end{proof}

\section{Weakly smooth forms and strong currents}
\label{sec:weak-smooth}

\begin{definition}[Weakly smooth form]
A \emph{weakly smooth} \((p,q)\)-form is locally an expression
\[
                                h^*\alpha,
\]
where \(h\) is a harmonic tropicalization and \(\alpha\) is a smooth
Lagerberg form on a neighbourhood of \(h(W)\).  Two expressions are
identified after passage to common harmonic refinements.
\end{definition}

Balancing makes the construction independent of representatives and gives
a bigraded differential algebra with \(\dprime\), \(\ddprime\), wedge product, and
functorial pullback.

We indicate why balancing is the operative hypothesis, since this is the
only point at which \Cref{thm:harmonic-balancing} is used.

\begin{proposition}
\label{prop:ch18-weakly-smooth-welldefined}
Let $h:W\to\mathbf R^r$ and $h':W\to\mathbf R^{r'}$ be harmonic
tropicalizations and $\alpha,\alpha'$ smooth Lagerberg forms near $h(W)$,
$h'(W)$.  Suppose $h^*\alpha$ and $h'^*\alpha'$ agree pointwise on $W$ in
the sense that they have the same restriction to every common harmonic
refinement.  Then they define the same weakly smooth form, and the
assignment
\[
 (h,\alpha)\longmapsto h^*\alpha
\]
is compatible with wedge products, with $\dprime$ and $\ddprime$, and with
pullback along a morphism of good strict analytic spaces.
\end{proposition}

\begin{proof}
Form the product harmonic tropicalization $(h,h'):W\to\mathbf
R^{r+r'}$, which is harmonic by
\Cref{prop:ch18-harmonic-functorial}\,\cref{itm:ch18-hf:compose}, and let
$p,p'$ be the two projections.  Both $h^*\alpha=(h,h')^*p^*\alpha$ and
$h'^*\alpha'=(h,h')^*p'^*\alpha'$ are pullbacks along the same map, so it
suffices to know that $(h,h')^*$ is injective on smooth forms defined near
the image, that is, that the image determines the form.  This is where
balancing enters: by \Cref{thm:harmonic-balancing} the $d$-dimensional part
of the image is a balanced weighted polyhedral set away from the excluded
boundary faces, so by
\Cref{thm:ch10-forms-presentation-independent} a smooth form on it is
determined by its restrictions to the cells, which are exactly the values of
the pullback on tangent vectors.  Without balancing the image could fail to
be a tropical cycle and the quotient of \eqref{eq:ch10-form-kernel} would not
be the right one.

Compatibility with wedge and with the two differentials is then formal:
these operations are computed on a common refinement and commute with
integral affine pullback by
\Cref{prop:ch09-pullback-functorial}.  Compatibility with pullback along
$\varphi:Y\to X$ is
\Cref{prop:ch18-harmonic-functorial}\,\cref{itm:ch18-hf:pullback} applied to
each coordinate of $h$.
\end{proof}

The principal comparison is
\begin{equation}
              A^{p,q}_{\mathrm{CLD}}\hookrightarrow
              A^{p,q}_{\mathrm{weak}};
\label{eq:cld-weak-inclusion}
\end{equation}
the map is injective by \cite[Proposition~10.14]{GJR26}.  For a continuous
function \(u\),
\[
 u\text{ is }\mathbf R\text{-harmonic}
 \quad\Longleftrightarrow\quad
 u\in A^{0,0}_{\mathrm{weak}}\text{ and }\dprimeddprime u=0
\]
by \cite[Proposition~10.17]{GJR26}.
For an $\mathbf R$-PL function on a separated boundaryless pure-dimensional
space, \Cref{prop:ch18-real-pl-zero-current-harmonic} removes the weakly
smooth hypothesis when the equality is interpreted as an equality of CLD
currents.  This does not identify the weakly smooth and CLD calculi, and it
does not extend the conclusion to an arbitrary continuous function.
The inclusion \eqref{eq:cld-weak-inclusion} is strict already for functions
on smooth curves, and monotone approximation by ordinary smooth harmonic
functions does not close the gap:
\Cref{prop:ch19-cld-smooth-harmonic-nonregularization} gives an integral PL
weakly smooth harmonic function which is not even a locally uniform limit of
ordinary CLD-smooth harmonic functions.  That statement does not concern
ordinary CLD-smooth psh approximation.

Weakly smooth forms admit integration, Stokes and Green formulas.  The duals
of compactly supported weakly smooth forms are called \emph{strong
currents}.  On a Hausdorff strict analytic space of pure dimension without
boundary, Poincar\'e--Lelong holds for strong currents
\cite[Theorem~12.7]{GJR26}.  The weakly smooth Dolbeault complex satisfies a
local Poincar\'e lemma \cite[Theorem~13.1]{GJR26} and has better behavior on
curves than the original smooth complex.

\begin{warning}{Two independent enlargements}
Weakly smooth forms allow more coordinate maps but still have smooth
Lagerberg coefficients.  Delta-forms allow polyhedral current coefficients
in ordinary tropical charts.  Thus
\eqref{eq:cld-weak-inclusion} does not identify weakly smooth forms with
delta-forms, and neither theory as currently stated contains every object of
the other.  A calculus of delta-forms based on harmonic tropicalizations
would be a further construction.
\end{warning}

\section{Semipositive PL and classically psh functions}
\label{sec:ch18-classical}

The harmonic theory also clarifies local positivity.  Let \(h\) be an
\(\mathbf R\)-PL function on a good strict analytic space.  It is
\emph{semipositive at \(x\)} if the residual class
\[
                    L_h(x)\in\operatorname{Pic}(\operatorname{red}(X,x))
                                  \otimes_{\mathbf Z}\mathbf R
\]
is nef.

\begin{theorem}\label{thm:ch18-semipositive-integral-residual-curve}
Let $K$ be a complete nontrivially valued non-Archimedean field, let $X$ be
a good strict $K$-analytic space, and let $x\in X$ satisfy
\[
 \operatorname{trdeg}_{\widetilde K}\widetilde{\mathscr H(x)}\leq1.
\]
Let
$\Lambda\subset\mathbf R$ be either $\mathbf Z$ or a divisible additive
subgroup.  If a $\Lambda$-PL function $h$ is semipositive at $x$, then there
are a strict analytic neighbourhood $W$ of $x$, finitely many coefficients
$a_j\in\Lambda_{\geq0}$, and integral PL functions $g_j$ on $W$, each
semipositive at every point of $W$, such that
\[
                         h|_W=\sum_j a_jg_j.
\]
All but at most one of the $g_j$ may be chosen integral harmonic.  The
statement allows boundary points and requires neither algebraic closedness
of $K$ nor any hypothesis relating $\Lambda$ to the value group of $K$.
In particular, it applies at every point of a good strict analytic curve.
\end{theorem}

\begin{proof}
If $\Lambda=\mathbf Z$, openness of the semipositivity locus
\cite[Remark~7.3(4)]{GJR26} lets us shrink around $x$ until $h$ is
semipositive on one strict analytic neighbourhood; then the one-term
decomposition $h=1\cdot h$ has all the asserted properties.  We may
therefore suppose that $\Lambda$ is divisible.

By \cite[Proposition~5.8]{GJR26}, after shrinking to a strict analytic
neighbourhood $U$ of $x$ we can write
\[
                         h=\sum_{i=1}^m\lambda_i h_i,
 \qquad \lambda_i\in\Lambda,\quad h_i\in\operatorname{PL}_{\mathbf Z}(U).
\]
Linearity of residues \cite[Lemma~6.18]{GJR26} gives
\[
                         L_h(x)=\sum_{i=1}^m\lambda_iL_{h_i}(x).
\]
The system of models of the reduction of the germ is filtered.  By
\cite[Lemma~6.5 and \S6.6]{GJR26}, we may therefore choose one model $V$ of
$\operatorname{red}(X,x)$ and line bundles $L_i$ on $V$ representing all
the residue classes $L_{h_i}(x)$.  By the hypothesis on $x$,
\[
        \dim V=\operatorname{trdeg}_{\widetilde K}
                    \widetilde{\mathscr H(x)}\leq1.
\]

Suppose first that $\dim V=0$, or that $\dim V=1$ and $V$ is not proper.
Then $V$ contains no integral proper closed curve.  Nefness is consequently
vacuous for every $L_i$ and its dual.  Thus every $h_i$ is harmonic at $x$
by \cite[Proposition~7.1]{GJR26}.  Writing
\[
          \lambda_i=\lambda_i^+-\lambda_i^-,\qquad
          \lambda_i^+=\max\{\lambda_i,0\},\quad
          \lambda_i^-=\max\{-\lambda_i,0\},
\]
expresses $h$ as a finite $\Lambda_{\geq0}$-linear combination of the
integral harmonic functions $h_i$ and $-h_i$.

It remains to treat the case in which $V$ is a proper integral curve.  Put
\[
                 d_i=\deg L_i\in\mathbf Z,\qquad
                 d=(d_1,\ldots,d_m),\qquad
                 D=d\cdot\lambda=\deg L_h(x).
\]
Semipositivity of $h$ says exactly that $D\geq0$.  If $d=0$, every $h_i$
is harmonic at $x$, and the preceding sign-splitting argument applies.
Assume henceforth that $d\ne0$.  Let
\[
                 \delta=\gcd(d_1,\ldots,d_m)>0,
                 \qquad \ell=\delta^{-1}d:\mathbf Z^m\to\mathbf Z.
\]
The homomorphism $\ell$ is primitive.  Choose $u\in\mathbf Z^m$ with
$\ell(u)=1$ and a $\mathbf Z$-basis
$v_1,\ldots,v_{m-1}$ of $\ker\ell$.  Since
$\mathbf Z^m=\mathbf Zu\oplus\ker\ell$, the same integral change of basis
over $\Lambda$ gives a unique decomposition
\[
            \lambda=t u+\sum_{j=1}^{m-1}c_jv_j,
            \qquad
            t=\ell(\lambda)=D/\delta\in\Lambda_{\geq0},
            \quad c_j\in\Lambda.
\]
Define integral PL functions on $U$ by
\[
                    q=\sum_i u_i h_i,
             \qquad r_j=\sum_i(v_j)_i h_i.
\]
Their residue line bundles on $V$ have degrees
\[
              \deg L_q(x)=d\cdot u=\delta>0,
              \qquad \deg L_{r_j}(x)=d\cdot v_j=0.
\]
On a proper integral curve, a line bundle is nef precisely when its degree
is nonnegative and is numerically trivial precisely when its degree is zero.
The residue criterion \cite[Proposition~7.1]{GJR26} therefore shows that
$q$ is semipositive at $x$ and that each $r_j$ and $-r_j$ is harmonic at
$x$.  Now
\[
 h=tq+\sum_j c_j^+r_j+\sum_j c_j^-(-r_j),
 \qquad
 c_j^+=\max\{c_j,0\},\quad c_j^-=\max\{-c_j,0\},
\]
is the required positive-coefficient decomposition at the level of germs.
All summands other than the possibly nonharmonic $q$ are harmonic.

Finally, the semipositivity and harmonicity loci of PL functions are open
by \cite[Remark~7.3(4)]{GJR26}.  Intersecting the finitely many such loci
with $U$ produces one strict analytic neighbourhood $W$ of $x$ on which
every displayed summand has the asserted property.  Restricting the identity
to $W$ completes the proof.
\end{proof}

The sentence immediately before \cite[Proposition~7.5]{GJR26} asks whether
the corresponding integral-decomposition characterization of harmonic PL
functions has a semipositive analogue.  Since semipositive germs form only a
monoid, rather than a group, by \cite[Remark~7.3(5)]{GJR26}, the nonnegative
coefficients in the theorem are the natural formulation.  The theorem
settles that formulation when the residual dimension is at most one.  In
higher residual dimension the nef cone of a common model need not be a
single rational half-space.  The following two complementary results settle
the unrestricted question.  If signs are allowed, semipositive integral
germs span every PL germ, so the condition does not detect semipositivity.
For the natural nonnegative-coefficient formulation, an irrational extremal
nef ray on an abelian surface gives a counterexample already in residual
dimension two.

\begin{theorem}
\label{thm:p09-signed-semipositive-span}
Let \(K\) be a complete nontrivially valued non-Archimedean field, let
\(X\) be a good strictly \(K\)-analytic space, and let \(x\in X\).
Let \(\Lambda\subset\mathbf R\) be either \(\mathbf Z\) or a divisible
additive subgroup.

For every finite family of integral PL germs
\(h_1,\ldots,h_m\) at \(x\), there exist an integral PL germ \(a\), an
integer \(N\geq1\), and a strict analytic neighbourhood \(W\) of \(x\)
on which all these germs are represented, such that
\[
             a,\qquad Na+h_i,\qquad Na-h_i
\]
are semipositive at every point of \(W\), for every
\(i\in\{1,\ldots,m\}\).

Consequently, every \(\Lambda\)-PL germ \(h\) at \(x\) admits, on one
strict analytic neighbourhood \(W\), an identity
\begin{equation}
\label{eq:p09-signed-semipositive-span}
                         h|_W=\sum_{j=1}^{r}c_jg_j-c\,a,
\end{equation}
where
\[
                  c_1,\ldots,c_r,c\in\Lambda_{\geq0}
\]
and \(a,g_1,\ldots,g_r\) are integral PL functions semipositive at every
point of \(W\).  In particular, every \(\Lambda\)-PL germ is a finite
signed combination, with coefficients in \(\Lambda\), of semipositive
integral PL germs, and every such finite combination is \(\Lambda\)-PL.
\end{theorem}

% Frozen source-aware proof candidate.  Do not insert into the live chapter
% before production Rethlas verification.  No novelty or priority claim is
% intended.

\paragraph{Imported inputs.}
The following primary-source statements are used with their stated
hypotheses.
\begin{enumerate}
\item Gubler--Jell--Rabinoff, Proposition~5.8: if
\(\Lambda=\mathbf Z\), or if \(\Lambda\subset\mathbf R\) is divisible,
then every \(\Lambda\)-PL function is analytically locally a finite
\(\Lambda\)-linear combination of integral PL functions; the converse
holds as well.
\item Gubler--Jell--Rabinoff, \S6.6 and Lemma~6.5: invertible sheaves on
the reduction of a good strict germ are represented on algebraic models,
the models form a filtered system under proper domination, and a finite
family of residue bundles can be represented on one common model.
Proper domination preserves the associated open subset of the
Riemann--Zariski space by the model formalism in their \S6.2.
\item Gubler--Jell--Rabinoff, Lemma~6.18: the residue map on PL germs is
linear.  Their Proposition~7.1 identifies semipositivity of a PL germ at
\(x\) with nefness of its residue class, and their Remark~7.3(4) says that
the semipositivity locus is open.
\item Chambert--Loir--Ducros, Theorem~17.5.7(b), in the version of
25~July~2025: on an almost-good strict germ the strict integral residual
functor is essentially surjective already for analytically trivial
underlying bundles.  In rank one, a residual line bundle is therefore the
residue of an integral PL metric on a locally trivial analytic line bundle.
A good strict germ satisfies the required almost-good hypothesis.  This is
the lifting input used in the proof of Gubler--Jell--Rabinoff,
Proposition~7.10.
\item Chow's lemma, in the form of Stacks Project, Tag~02O2, gives a
proper birational quasi-projective model dominating an integral separated
finite-type model over the residue field.  The ample-twist criterion of
Stacks Project, Tag~0892, says that if \(A\) is ample and \(M\) is an
invertible sheaf on such a model, then \(A^{\otimes N}\otimes M\) is ample
for all sufficiently large \(N\).
\end{enumerate}

\begin{proof}
All assertions are local at \(x\).  Since \(X\) is good and strict, we
may first replace it by a strictly affinoid neighbourhood of \(x\).
This neighbourhood is separated, and its germ at \(x\) is good, hence
almost good in the sense required by
\cite[Theorem~17.5.7(b)]{CLD25}.

Put
\[
                         \mathscr R=\operatorname{red}(X,x).
\]
For each \(i\), the residue
\[
                         L_{h_i}(x)\in\operatorname{Pic}(\mathscr R)
\]
is represented by a line bundle on a model of \(\mathscr R\).  The
description of invertible sheaves in \cite[\S6.6]{GJR26}, together with
the filteredness statement \cite[Lemma~6.5]{GJR26}, therefore gives one
model \(V_0\) of \(\mathscr R\) and line bundles \(M_{i,0}\) on \(V_0\)
representing all the \(L_{h_i}(x)\).

Apply Chow's lemma \cite[Tag~02O2]{stacks-project} to \(V_0\) over
\(\widetilde K\).  It gives a proper surjective morphism
\(Z\to V_0\), with \(Z\) quasi-projective, which is an isomorphism over
a dense open subset of \(V_0\).  Replacing \(Z\) by the reduced closure
of the inverse image of that dense open subset, we obtain an integral
quasi-projective variety \(V\) and a proper birational morphism
\[
                             \pi:V\longrightarrow V_0.
\]
By \cite[\S6.2]{GJR26}, proper domination does not change the associated
open subset of the Riemann--Zariski space.  Hence \(V\) is again a model
of \(\mathscr R\).  Put
\[
                             M_i=\pi^*M_{i,0}.
\]

Choose an ample line bundle \(A\) on \(V\).  Apply the ample-twist criterion
\cite[Tag~0892]{stacks-project} with \(f=\operatorname{id}_V\), with
\(\mathcal M=A\), and successively with
\(\mathcal L=M_i,M_i^{-1}\).  Every invertible sheaf is
\(\operatorname{id}_V\)-ample, since its restriction to an affine open is
ample.  Hence, after taking the maximum of finitely many thresholds, there
is one integer \(N\geq1\) such that
\begin{equation}
\label{eq:p09-ample-order-interval}
                A^{\otimes N}\otimes M_i,
                \qquad
                A^{\otimes N}\otimes M_i^{-1}
\end{equation}
are ample for every \(i\).  In particular, \(A\) and all the line
bundles in \eqref{eq:p09-ample-order-interval} are nef: their restrictions
to every proper integral curve have nonnegative degree.

The line bundle \(A\) determines an invertible sheaf on
\(\mathscr R\).  The strict lifting theorem
\cite[Theorem~17.5.7(b)]{CLD25}, used for line bundles in precisely this
way in the proof of \cite[Proposition~7.10]{GJR26}, gives a strictly
integral PL metric on an analytically trivial line bundle over the germ
\((X,x)\) whose residual line bundle is represented by \((V,A)\).
Choosing a trivializing frame and taking the negative logarithm of its
norm gives an integral PL germ \(a\) satisfying
\[
                             L_a(x)\simeq[A].
\]
After shrinking, regard \(a,h_1,\ldots,h_m\) as actual functions on one
common neighbourhood of \(x\).

Residue linearity \cite[Lemma~6.18]{GJR26} now gives
\[
 \begin{aligned}
 L_{Na+h_i}(x)
   &\simeq [A^{\otimes N}\otimes M_i],\\
 L_{Na-h_i}(x)
   &\simeq [A^{\otimes N}\otimes M_i^{-1}].
 \end{aligned}
\]
These residue classes, as well as \(L_a(x)\), are nef.  The residue
criterion \cite[Proposition~7.1]{GJR26} shows that \(a\) and all
\(Na\pm h_i\) are semipositive at \(x\).  Their semipositivity loci are
open by \cite[Remark~7.3(4)]{GJR26}.  Intersecting these finitely many
open loci and then choosing a strict analytic neighbourhood \(W\) of
\(x\) inside the intersection proves the first assertion.

Let now \(h\) be a \(\Lambda\)-PL germ.  By
\cite[Proposition~5.8]{GJR26}, after shrinking around \(x\), there is an
identity of actual functions
\[
                              h=\sum_{i=1}^{m}\lambda_i h_i,
 \qquad
 \lambda_i\in\Lambda,\quad
 h_i\in\operatorname{PL}_{\mathbf Z}.
\]
Apply the first assertion to \(h_1,\ldots,h_m\), and restrict everything
to the resulting common neighbourhood \(W\).  Omit the indices for
which \(\lambda_i=0\).  For every remaining index, put
\[
 \epsilon_i=
 \begin{cases}
  1,&\lambda_i>0,\\
 -1,&\lambda_i<0,
 \end{cases}
 \qquad
 c_i=|\lambda_i|,
 \qquad
 g_i=Na+\epsilon_i h_i,
\]
and set
\[
                              c=N\sum_i|\lambda_i|.
\]
Every \(g_i\), as well as \(a\), is an integral PL function semipositive
on \(W\).  Moreover,
\[
 \begin{aligned}
 \sum_i c_ig_i-ca
 &=\sum_i|\lambda_i|
       \bigl(Na+\epsilon_i h_i\bigr)
       -N\sum_i|\lambda_i|a\\
 &=\sum_i\lambda_i h_i
 =h.
 \end{aligned}
\]
Since \(\Lambda\) is an additive subgroup of \(\mathbf R\), one has
\[
                 |\lambda_i|\in\Lambda_{\geq0},
 \qquad
                 N\sum_i|\lambda_i|\in\Lambda_{\geq0}.
\]
Thus \eqref{eq:p09-signed-semipositive-span} has all the asserted
coefficient properties.  If all \(\lambda_i\) vanish, take the sum to
be empty and \(c=0\).

Conversely, every finite \(\Lambda\)-linear combination of integral PL
functions is \(\Lambda\)-PL, by the reverse implication in
\cite[Proposition~5.8]{GJR26}.  This proves the final assertion.
\end{proof}

\begin{remark}
\label{rem:p09-signed-semipositive-span-scope}
The proof uses the common residual model only to establish nefness of the
residues of the actual functions \(Na\pm h_i\); it never replaces the
original \(h_i\) by unrelated lifts.  Thus the displayed identity is an
identity of functions, not merely of residue classes.  If \(a'\) is another
integral PL lift of \([A]\), then \(L_{a-a'}(x)=0\) by residue linearity,
and \cite[Lemma~6.18 and Proposition~7.1]{GJR26} identify \(a-a'\) as a
harmonic germ.  In any case, replacing \(a\) by \(a'\) in both occurrences
of \eqref{eq:p09-signed-semipositive-span}, and redefining
\(g_i'=Na'+\epsilon_i h_i\), produces exact cancellation, so there is no
hidden kernel correction.

No interior-point assumption is used.  At a boundary point the first
common residual model need not be proper; Chow's lemma supplies a proper
dominating quasi-projective model, and quasi-projectivity is sufficient for
the ample order-unit argument.

Finally, the theorem is not a positive-cone decomposition theorem.  Its
proof never uses semipositivity of \(h\), and the coefficient \(-c\) in
\eqref{eq:p09-signed-semipositive-span} is generally negative.  It
identifies the signed \(\Lambda\)-span of the semipositive integral germs;
when \(\Lambda=\mathbf Z\), this is their group completion.  It gives no
positive-coefficient decomposition of a semipositive \(\Lambda\)-PL germ.
\end{remark}

\begin{proposition}
\label{prop:p09-irrational-nef-ray-obstruction}
Let \(K\) be a complete nontrivially valued non-Archimedean field whose
residue field \(k=\widetilde K\) is algebraically closed of characteristic
zero.  Let \(X\) be a good strict \(K\)-analytic space and let \(x\in X\)
satisfy
\[
 \widetilde{\mathscr H(x)}=k(S),
 \qquad
 \operatorname{red}(X,x)=P_{k(S)/k},
\]
where \(S\) is a smooth projective surface over \(k\).  Let \(h\) be an
\(\mathbf R\)-PL germ at \(x\), semipositive at \(x\), and suppose that
\(L_h(x)\) is represented on \(S\) by a nef class
\(\alpha\in\operatorname{Pic}(S)\otimes_{\mathbf Z}\mathbf R\) with
nonzero numerical class.  Assume that this class spans an extremal
ray of \(\operatorname{Nef}(S)\) containing no nonzero rational class:
\begin{equation}
\label{eq:p09-irrational-extremal-ray}
 \mathbf R_{\geq0}[\alpha]\cap N^1(S)_{\mathbf Q}=\{0\}.
\end{equation}
Then \(h\) admits no identity
\begin{equation}
\label{eq:p09-forbidden-positive-decomposition}
                         h=\sum_{j=1}^{m}a_jg_j
\end{equation}
of germs at \(x\), where \(a_j\in\mathbf R_{\geq0}\) and the \(g_j\) are
integral PL germs semipositive at \(x\).

Moreover, if \(\Lambda\subset\mathbf R\) is a divisible additive subgroup
and
\[
                 \alpha\in
                 \operatorname{Pic}(S)\otimes_{\mathbf Z}\Lambda,
\]
then there exists a semipositive \(\Lambda\)-PL germ at \(x\) having
residue \(\alpha\).  Thus the obstruction is nonvacuous for every such
\(\Lambda\)-rational residual class.
\end{proposition}

\begin{corollary}
\label{cor:p09-good-reduction-positive-cone-counterexample}
Put \(K=\mathbf C((t))\), with its \(t\)-adic absolute value normalized by
\(-\log|t|=1\), and let
\[
 E=\mathbf C/(\mathbf Z+\pi\mathrm i\,\mathbf Z),
 \qquad S=E\times E,
 \qquad X=(S_K)^{\mathrm{an}}.
\]
Let \(x\) be the divisorial point reducing to the generic point of the
special fibre of the constant model
\(S\times_{\mathbf C}\operatorname{Spec}\mathbf C[[t]]\), and put
\[
                         \Lambda=\mathbf Q+\mathbf Q\sqrt2.
\]
There are a strict analytic neighbourhood \(U\) of \(x\) and a
semipositive \(\Lambda\)-PL function \(h\) on \(U\) such that the germ of
\(h\) at \(x\) has no decomposition
\eqref{eq:p09-forbidden-positive-decomposition}.  Thus the natural
positive-cone analogue of the local integral decomposition for harmonic PL
functions is false in residual dimension two, already at an interior
good-reduction point of a smooth proper analytic surface.
\end{corollary}

% Frozen source-aware proof candidate.  Do not insert into the live chapter
% before production Rethlas verification.  No novelty or priority claim is
% intended.

\paragraph{Imported inputs.}
The proof uses the following primary-source statements with their stated
hypotheses.
\begin{enumerate}
\item Chambert--Loir--Ducros, Theorem~17.5.7(b), in the version of
25~July~2025, says that the strict integral residual functor on an
almost-good strict germ is essentially surjective already in the subcategory
whose underlying analytic bundle is trivial.  In rank one, every line bundle
on a model of the reduction of a good strict germ is therefore represented
by the residue of an integral PL germ.
\item Gubler--Jell--Rabinoff, \S\S6.2 and~6.6 and Lemma~6.5, identify the
Picard group of the reduction of a germ with the filtered colimit over its
models and provide common dominating models for finite collections of
residual line bundles.  Lemma~6.11 says that every model representing a nef
residual class is nef.  Lemma~6.18 gives linearity of residues, and
Proposition~7.1 identifies semipositivity at a point with nefness of the
residual class.  Remark~7.3(4) says that the semipositivity locus is open.
\item For a product principally polarized abelian variety, the Rosati map
identifies its real N\'eron--Severi space with the symmetric subspace of its
real endomorphism algebra and identifies the ample cone with the cone of
positive-definite matrices.  This is Prendergast-Smith, Theorem~4.3.  The
nef cone is its closure, the cone of positive-semidefinite matrices.
For \(S=E\times E\) with \(\operatorname{End}(E)=\mathbf Z\), the three
fibre and diagonal classes used below form an integral basis of
\(\operatorname{NS}(S)\); compare Bauer--Borntr\"ager, Lemma~3.1 and
Remark~3.4.
\item Chow's lemma \cite[Tag~02O2]{stacks-project} and resolution of
surface singularities \cite[Tag~0BGP]{stacks-project} give a smooth
projective surface dominating any proper surface model of the same function
field.  A proper birational morphism between smooth surfaces factors into
blowups in closed points \cite[Tag~0C5Q]{stacks-project}.  For such a
blowup \(b:Y'\to Y\), with exceptional curve \(E\), the sequence
\[
 0\longrightarrow\operatorname{Pic}(Y)
  \xrightarrow{b^*}\operatorname{Pic}(Y')
  \xrightarrow{\deg_E}\mathbf Z\longrightarrow0
\]
is exact \cite[Tag~0C5L]{stacks-project}.
\end{enumerate}

\begin{proof}[Proof of \cref{prop:p09-irrational-nef-ray-obstruction}]
Suppose, for a contradiction, that
\eqref{eq:p09-forbidden-positive-decomposition} holds.  Delete every term
with \(a_j=0\); thus all remaining \(a_j\) are strictly positive.  Taking
residues gives
\begin{equation}
\label{eq:p09-residual-positive-decomposition}
                 \alpha=\sum_{j=1}^{m}a_jL_{g_j}(x)
 \quad\text{in}\quad
 \operatorname{Pic}(P_{k(S)/k})\otimes_{\mathbf Z}\mathbf R.
\end{equation}
By the filtered common-model formalism
\cite[\S\S6.2, 6.6 and Lemma~6.5]{GJR26}, there is one model \(V\) of
the full Riemann--Zariski space \(P_{k(S)/k}\) dominating \(S\) and all
models involved in \eqref{eq:p09-residual-positive-decomposition}, far
enough that the equality itself is represented on \(V\).  Because it is a
model of the full Riemann--Zariski space, the morphism \(V\to S\) is proper.
Both models have function field \(k(S)\), so this morphism is birational.
Moreover, every integral line bundle on \(V\) representing
\(L_{g_j}(x)\) is nef by \cite[Lemma~6.11]{GJR26}.

Chow's lemma \cite[Tag~02O2]{stacks-project} and surface resolution
\cite[Tag~0BGP]{stacks-project} give a smooth projective surface \(Y\)
dominating \(V\).  Denote the composite morphism by
\[
                              f:Y\longrightarrow S.
\]
Pulling the represented equality to \(Y\), and then passing to numerical
classes, yields
\begin{equation}
\label{eq:p09-upstairs-nef-decomposition}
                       f^*[\alpha]=\sum_{j=1}^{m}a_j[M_j]
              \quad\text{in }N^1(Y)_{\mathbf R},
\end{equation}
where every \(M_j\) is an integral nef line bundle on \(Y\).

Let \(C\subset Y\) be an irreducible \(f\)-exceptional curve.  The
projection formula and \eqref{eq:p09-upstairs-nef-decomposition} give
\[
 0=f^*[\alpha]\cdot C
   =\sum_{j=1}^{m}a_j(M_j\cdot C).
\]
Every intersection number on the right is nonnegative, and every \(a_j\)
is positive.  Hence
\begin{equation}
\label{eq:p09-exceptional-orthogonality}
                              M_j\cdot C=0
\end{equation}
for every \(j\) and every \(f\)-exceptional curve \(C\).

By \cite[Tag~0C5Q]{stacks-project}, the morphism \(f\) factors as
\[
 Y=Y_r\xrightarrow{b_r}Y_{r-1}\xrightarrow{b_{r-1}}\cdots
   \xrightarrow{b_1}Y_0=S,
\]
where every \(b_\ell\) is the blowup of a closed point.  Fix \(j\), set
\(M_{j,r}=M_j\), and consider the exceptional curve \(E_r\subset Y_r\)
of \(b_r\).  It is \(f\)-exceptional, so
\eqref{eq:p09-exceptional-orthogonality} gives
\(\deg_{E_r}M_{j,r}=0\).  The exact Picard sequence for a point blowup
\cite[Tag~0C5L]{stacks-project} therefore supplies a line bundle
\(M_{j,r-1}\) on \(Y_{r-1}\) such that
\[
                              M_{j,r}\simeq b_r^*M_{j,r-1}.
\]
The line bundle \(M_{j,r-1}\) is nef: if \(D\subset Y_{r-1}\) is an
irreducible curve and \(\widetilde D\subset Y_r\) its strict transform,
then
\[
                     M_{j,r-1}\cdot D
                   = M_{j,r}\cdot\widetilde D\geq0.
\]
It also has degree zero on every curve contracted by
\(Y_{r-1}\to S\), since the strict transform of such a curve is
\(f\)-exceptional.  Iterating the same argument gives a nef line bundle
\(B_j=M_{j,0}\) on \(S\) with
\[
                         M_j\simeq f^*B_j.
\]
Consequently, for the integral nef class
\(\beta_j=c_1(B_j)\in N^1(S)_{\mathbf Z}\), one has
\[
                         [M_j]=f^*\beta_j
                  \quad\text{in }N^1(Y)_{\mathbf R}.
\]

Substituting into \eqref{eq:p09-upstairs-nef-decomposition} and using the
injectivity of \(f^*:N^1(S)_{\mathbf R}\to N^1(Y)_{\mathbf R}\), we obtain
\begin{equation}
\label{eq:p09-downstairs-nef-decomposition}
                              [\alpha]=\sum_{j=1}^{m}a_j\beta_j
                       \quad\text{in }N^1(S)_{\mathbf R}.
\end{equation}
Each summand is nef.  Since \([\alpha]\) spans an extremal nef ray,
every nonzero \(\beta_j\) belongs to \(\mathbf R_{\geq0}[\alpha]\).  At
least one \(\beta_j\) is nonzero because \([\alpha]\ne0\).  This contradicts
\eqref{eq:p09-irrational-extremal-ray}, since
\(\beta_j\in N^1(S)_{\mathbf Z}\subset N^1(S)_{\mathbf Q}\).
Therefore the positive decomposition cannot exist.

For the last assertion, write the residual class as a finite sum
\[
                         \alpha=\sum_{i=1}^{q}\lambda_i[L_i],
 \qquad
             \lambda_i\in\Lambda,\quad L_i\in\operatorname{Pic}(S).
\]
Replace \(X\) by a strictly affinoid neighbourhood of \(x\).  It is
separated, and its good germ at \(x\) is almost good.
The strict residual lifting theorem
\cite[Theorem~17.5.7(b)]{CLD25} supplies a strictly integral PL metric on
an analytically trivial line bundle for each residual bundle \([L_i]\).
Choose a local trivializing frame and let \(u_i\) be the negative logarithm
of its norm.  Compatibility of the strict residual functor with the residue
line-bundle convention, exactly as in the realization step in the proof of
\cite[Proposition~7.10]{GJR26}, gives
\[
                              L_{u_i}(x)=[L_i].
\]
After shrinking once so that all these germs are represented by functions,
put
\[
                              h_\alpha=\sum_{i=1}^{q}\lambda_i u_i.
\]
This is a \(\Lambda\)-PL germ.  Residue linearity
\cite[Lemma~6.18]{GJR26} gives \(L_{h_\alpha}(x)=\alpha\), and the nefness
of \(\alpha\), together with the residue criterion
\cite[Proposition~7.1]{GJR26}, shows that \(h_\alpha\) is semipositive at
\(x\).
\end{proof}

\begin{proof}[Proof of
\cref{cor:p09-good-reduction-positive-cone-counterexample}]
Set \(\tau=\pi\mathrm i\).  We first check that \(E\) has no complex
multiplication.  If multiplication by \(z\in\mathbf C\) preserves
\(\mathbf Z+\tau\mathbf Z\), then for some integers \(a,b,c,d\) one has
\[
                              z=a+b\tau,
 \qquad
                              z\tau=c+d\tau.
\]
Thus
\[
                       b\tau^2+(a-d)\tau-c=0.
\]
The number \(\tau\) is transcendental, so \(b=c=0\) and \(a=d\).  Hence
\(z=a\in\mathbf Z\), proving
\[
                              \operatorname{End}(E)=\mathbf Z.
\]

Let \(p_1,p_2:S\to E\) be the projections, and put
\[
 H_1=p_1^*(0)=\{0\}\times E,
 \qquad
 H_2=p_2^*(0)=E\times\{0\},
 \qquad
 \Delta=(p_1-p_2)^*(0).
\]
The classes \(H_1,H_2,\Delta\) form a \(\mathbf Z\)-basis of
\(\operatorname{NS}(S)\) by
\cite[Remark~3.4]{BauerBorntraeger18}.  For completeness, the Rosati
identification is
\[
\rho:\operatorname{NS}(S)\simeq
 \left\{
   \begin{pmatrix}r&s\\s&u\end{pmatrix}:r,s,u\in\mathbf Z
 \right\}
\]
for the product principal polarization \(\Theta=H_1+H_2\), where
\(\rho([D])=\phi_\Theta^{-1}\circ\phi_D\).  Under this identification,
\[
 H_1\longmapsto
 \begin{pmatrix}1&0\\0&0\end{pmatrix},
 \qquad
 H_2\longmapsto
 \begin{pmatrix}0&0\\0&1\end{pmatrix},
 \qquad
 \Delta\longmapsto
 \begin{pmatrix}1&-1\\-1&1\end{pmatrix}.
\]
Indeed, if \(\Theta_E=(0)\) is the principal polarization on \(E\), the
functorial formula for the polarization homomorphism of \(q^*\Theta_E\)
gives the symmetric matrix \(q^{\mathsf T}q\).  Taking
\(q=p_1,p_2,p_1-p_2\) gives exactly the three displayed matrices and fixes
the sign of the off-diagonal entries.
The ample cone corresponds to the positive-definite cone by
\cite[Theorem~4.3]{PrendergastSmith12}, and therefore the nef cone
corresponds to the positive-semidefinite cone.

Put \(s=\sqrt2\) and consider the Picard class
\[
 \widehat\xi=(2+s)[\mathcal O_S(H_1)]
        +(1+s)[\mathcal O_S(H_2)]-s[\mathcal O_S(\Delta)]
        \in\operatorname{Pic}(S)\otimes_{\mathbf Z}\Lambda
\]
and its numerical image
\begin{equation}
\label{eq:p09-explicit-irrational-nef-class}
               \xi=(2+s)H_1+(1+s)H_2-s\Delta
                 \in\operatorname{NS}(S)\otimes_{\mathbf Z}\Lambda.
\end{equation}
Its Rosati matrix is
\begin{equation}
\label{eq:p09-rank-one-rosati-matrix}
 Q=
 \begin{pmatrix}2&s\\s&1\end{pmatrix}
 =
 \begin{pmatrix}s\\1\end{pmatrix}
 \begin{pmatrix}s&1\end{pmatrix}.
\end{equation}
Thus \(\xi\) is nonzero and nef.  Its ray is extremal.  Indeed, if
\(Q=A+B\) with \(A,B\) positive semidefinite and if \(w\in\ker Q\), then
\[
                    0=w^{\mathsf T}Aw+w^{\mathsf T}Bw.
\]
Both summands are nonnegative, so both vanish.  For a positive-semidefinite
matrix, \(w^{\mathsf T}Aw=0\) implies \(Aw=0\); hence
\(\ker Q\subset\ker A\cap\ker B\).  Since \(Q\) has rank one, the images
of \(A\) and \(B\) lie in the one-dimensional image of \(Q\), and
\(A,B\) are nonnegative scalar multiples of \(Q\).

The ray is irrational.  If \(r\xi\in N^1(S)_{\mathbf Q}\), the difference
between its \(H_1\)- and \(H_2\)-coefficients is \(r\), whereas its
\(\Delta\)-coefficient is \(-r\sqrt2\).  Both are rational only when
\(r=0\).  Thus
\[
                   \mathbf R_{\geq0}\xi\cap N^1(S)_{\mathbf Q}=\{0\}.
\]

For the divisorial point \(x\) of the constant good-reduction model, the
residue of its completed residue field is the function field of the closure
of its reduction, namely
\[
 \widetilde{\mathscr H(x)}=\mathbf C(S).
\]
The construction in \cite[\S6.13]{GJR26} identifies the reduction of the
germ with
\[
 \operatorname{red}(X,x)=P_{\mathbf C(S)/\mathbf C}\{S\}.
\]
Since \(S\) is proper, the valuative criterion gives
\(P_{\mathbf C(S)/\mathbf C}\{S\}=P_{\mathbf C(S)/\mathbf C}\).
The last assertion of
\cref{prop:p09-irrational-nef-ray-obstruction}, applied to
\(\alpha=\widehat\xi\),
gives a semipositive \(\Lambda\)-PL germ with residue \(\widehat\xi\); the first
assertion excludes every nonnegative-real decomposition of this germ into
semipositive integral PL germs.  By the openness of its
semipositivity locus \cite[Remark~7.3(4)]{GJR26}, the germ is represented by
a semipositive function on one strict analytic neighbourhood \(U\) of
\(x\), as asserted.
\end{proof}

\begin{remark}
\label{rem:p09-positive-coefficients-essential}
The sign restriction in
\eqref{eq:p09-forbidden-positive-decomposition} is essential.  Each of
\(H_1,H_2,\Delta\) is nef, so their integral PL lifts may be chosen
semipositive at \(x\).  Formula
\eqref{eq:p09-explicit-irrational-nef-class} therefore writes the same
counterexample as a signed \(\Lambda\)-linear combination of semipositive
integral PL germs.  With negative coefficients, intersections against
exceptional curves can cancel and the descent argument above no longer
applies.
\end{remark}


Thus there is a sharp distinction between the two meanings of linear
combination.  Signed generation is automatic in every residual dimension
by \cref{thm:p09-signed-semipositive-span}; positive generation holds in
residual dimension at most one by
\cref{thm:ch18-semipositive-integral-residual-curve}, but fails in residual
dimension two by
\cref{cor:p09-good-reduction-positive-cone-counterexample}.  No
higher-dimensional positive-cone assertion survives without additional
rational-polyhedral hypotheses on the relevant nef cone.

\begin{theorem}\label{thm:gr-pl-curve-test}
Let \(K\) be complete and nontrivially valued, \(X\) a good strict
\(K\)-analytic space, \(x\in X\), and \(h\) an \(\mathbf R\)-PL function.
Then \(h\) is semipositive at \(x\) if and only if on one neighbourhood
\(W\) of \(x\), for every sufficiently small strict analytic neighbourhood
\(U\subset W\) of \(x\) and every Zariski-closed analytic curve
\(C\subset U\), the restriction \(h|_C\) is semipositive.
\end{theorem}
\begin{proof}
If $h$ is semipositive at $x$, shrink to a neighborhood on which its residual
line bundles are nef.  Restriction to a closed analytic curve pulls those
bundles back to nef residual line bundles, so $h|_C$ is semipositive.

Conversely, suppose the curve test holds but $h$ is not semipositive at
$x$.  Its residual line bundle on the reduction of the germ then has
negative degree on a proper residual curve.  The local lifting theorem
produces, after shrinking a strict neighborhood and if necessary making an
analytic scalar extension, a Zariski-closed analytic curve whose reduction
contains that residual curve with positive generic degree.  The projection
formula identifies the degree of the restricted residual metric with that
negative degree, multiplied by the positive lifting degree.  Thus $h$ is not
semipositive on the lifted analytic curve, contradicting the hypothesis.
This proves the equivalence; the lifting construction and degree comparison
are \cite[Theorem~5.3]{GR25PL}.

The pointwise-closedness theorem
\cite[Theorem~5.6]{GR25PL} is a separate consequence: a pointwise limit of
semipositive $\mathbf R$-PL functions that remains $\mathbf R$-PL restricts
to such a limit on every curve, where the one-dimensional slope measures are
closed; the equivalence just proved then returns semipositivity.  Both
results are local and extend beyond the proper algebraic scope of
\cite{GM19}.
\end{proof}

Gubler--Rabinoff then introduce a larger, curve-tested class.

\begin{definition}[Classically psh]
An upper semicontinuous function \(u:X\to[-\infty,\infty)\) is
\emph{classically psh} if, after every analytic extension \(K'/K\), its
pullback along every morphism from a \(K'\)-analytic curve is subharmonic in
the functorial extension of Thuillier's sense.
\end{definition}

This is a sheaf-theoretic definition; it is stable under pullback and base
change by construction.

\begin{theorem}\label{thm:gr-classical-pl}
If \(K\) is nontrivially valued, \(X\) is good and strict, and \(h\) is an
\(\mathbf R\)-PL function, then
\[
                    h\text{ is semipositive}
       \quad\Longleftrightarrow\quad
                    h\text{ is classically psh}.
\]
\end{theorem}
\begin{proof}
This is the same equivalence as \Cref{thm:gr-pl-comparison}.  For
completeness, semipositivity is preserved by scalar extension and analytic
pullback, and on a smooth strict analytic curve it is equivalent to
subharmonicity; hence a semipositive PL function is classically psh.  In the
reverse direction, classical psh restricts to a subharmonic, hence
semipositive, function on every lifted analytic curve.  The curve test in
\Cref{thm:gr-pl-curve-test} then makes every residual line bundle nef.
Therefore the function is semipositive.  This is
\cite[Theorem~4.3]{GR25Classical}.

For arbitrary classically psh functions, the cited paper proves local and
affinoid maximum principles but does not construct Monge--Amp\`ere measures.
\end{proof}

CLD v2 psh functions are classically psh: functoriality and base change
reduce the statement to curves, where their currents are positive and
Wanner's theorem identifies \emph{current positivity} with
Thuillier subharmonicity.  The argument is recorded in
\cite[\S1.2]{GR25Classical}, using \cite[Theorem~4.7]{Wanner19}.  Wanner's
terminology does not by itself give the reverse inclusion into the
approximation-generated CLD hierarchy.  For real PL functions on a
boundaryless good strict space, however, the missing reverse arrow is proved in
\Cref{thm:verified-pl-psh-equivalence} by combining local LFS regularization
with a decreasing CLD log-sum-exp approximation.  At a relative-boundary
point that argument is not available.  In particular, a curve-test definition
alone does not supply Bedford--Taylor products or finite-energy theory.

\section{Model functions and piecewise affine potentials}
\label{sec:ch18-model-comparison}

\Cref{chap:metrics} constructed model functions from vertical divisors and
computed their Chambert--Loir measures by intersection theory on a model.
The present chapter has constructed corner cycles and Chern delta-forms from
piecewise affine potentials on skeleta.  This section proves that the two
constructions are the same, term by term.  It is the promised return to
\Cref{chap:metrics}, and it is the local half of the master dictionary of
\Cref{chap:master-dictionary}.

\subsection{The three categories of potential}
\label{sec:ch18-three-potentials}

Fix a complete nontrivially valued field $K$ and a proper $K$-variety $X$ of
dimension $n$.  Three classes of function on $X^{\An}$ have appeared, and
they must not be confused.

\begin{itemize}
\item \emph{Model functions.}  For a normal model $\cX$ and a vertical
 $\mathbf Q$-Cartier divisor $D$, the function $\varphi_D$ of
 \Cref{prop:metrics-divisor-model-function}.  These form the group
 $\mathcal D(X)_{\mathbf Q}$, which is dense in $C^0(X^{\An})$ by
 \Cref{prop:metrics-model-function-density}.
\item \emph{$\mathbf Q$-PL functions.}  Functions which are locally of the
 form $\frac1m\log|f|$ with $f$ invertible, glued on a finite cover by
 strict analytic domains; equivalently, by
 \Cref{thm:gm-metric-comparison}\,(1), the weights of formal
 $\mathbf Q$-metrics on $\cO_X$.
\item \emph{Piecewise affine potentials on a skeleton.}  Functions on a
 skeleton $\Gamma\subseteq X^{\An}$ which are affine on each cell of some
 polyhedral structure, extended to $X^{\An}$ by the retraction
 $\rho_\Gamma$.
\end{itemize}

\begin{proposition}
\label{prop:ch18-three-classes}
Let $K$ be discretely valued with $v_K(\pi)=1$ and let $X$ be a proper
geometrically connected $K$-variety.
\begin{enumerate}
\item\label{itm:ch18-tc:modelpl}
 Every model function is a $\mathbf Q$-PL function, and conversely every
 $\mathbf Q$-PL function on all of $X^{\An}$ is a model function.
\item\label{itm:ch18-tc:skeletal}
 If $\cX$ is an SNC model with dual complex $\Delta_\cX$ and retraction
 $\rho_\cX$, then $\varphi_D=\varphi_D\circ\rho_\cX$ for every vertical
 divisor $D$ supported on $\cX_s$, and $\varphi_D|_{\Delta_\cX}$ is affine
 on each face, with the values $a_E/b_E$ at the vertices.
\item\label{itm:ch18-tc:notall}
 Not every piecewise affine function on $\Delta_\cX$ is the restriction of a
 model function on $\cX$: the values at vertices must lie in
 $\frac1{b_E}\mathbf Z$, and the function must be affine with respect to the
 integral affine structure of $\Delta_\cX$.
\end{enumerate}
\end{proposition}

\begin{proof}
\cref{itm:ch18-tc:modelpl} A model function is the weight of a model metric
on $\cO_X$, and model metrics are formal, hence PL, by
\Cref{thm:gm-metric-comparison}\,(2); the passage between $\mathbf Z$ and
$\mathbf Q$ coefficients is the $\mathbf Q$-linearization in the same
theorem.  Conversely a $\mathbf Q$-PL function on the whole proper
$X^{\An}$ is the weight of a formal $\mathbf Q$-metric on $\cO_X$, which by
the same theorem is an algebraic model metric; its model line bundle differs
from $\cO_\cX$ by a vertical divisor, whose model function is the given one.

\cref{itm:ch18-tc:skeletal} Invariance under $\rho_\cX$ for a divisor
determined on $\cX$ is the last assertion of
\Cref{cor:pl-model-function}; affinity on faces is
\Cref{prop:skeleton-model-pa}, and the vertex values are
\Cref{prop:metrics-divisor-model-function}\,(2) together with the
description of $\Delta_\cX$ by monomial valuations in
\Cref{chap:skeleta}.

\cref{itm:ch18-tc:notall} The vertex values are $a_E/b_E$ with
$a_E\in\mathbf Z$ by \cref{itm:ch18-tc:skeletal}, which is the stated
integrality; and affinity is part of \cref{itm:ch18-tc:skeletal}.  A function
violating either is not a model function on that model, although by
\Cref{cor:pl-model-function} every rational PL function on
$\Delta_\cX$ does arise from a model function after passing to a dominating
model, and model functions are dense in $C^0(X^{\An})$ by
\Cref{prop:metrics-model-function-density}.
\end{proof}

\subsection{The curvature of a model function}
\label{sec:ch18-model-curvature}

The following theorem is the main comparison.  It says that the corner cycle
of the skeletal restriction, the Chern delta-form, and the vertical
intersection numbers are the same object written three ways.

\begin{theorem}
\label{thm:ch18-model-corner}
Let $K$ be discretely valued with $v_K(\pi)=1$, let $X$ be a proper
$K$-variety of dimension $n$, let $\cX$ be a normal model, and let
$L$ be a line bundle with a model metric induced by a model line bundle
$\cL$ on $\cX$.  Then:
\begin{enumerate}
\item\label{itm:ch18-mc:delta}
 the metric is a delta-metric, and its Chern delta-form
 $c_1^\delta(\overline L)$ is, in every tropical chart, the sum of the
 facewise Hessian and the corner cycle
 \eqref{eq:ch18-corner-weight} of the local metric weight;
\item\label{itm:ch18-mc:PL}
 if the weight is piecewise affine---in particular for a model metric---the
 facewise Hessian vanishes and $c_1^\delta(\overline L)$ is the corner
 cycle alone;
\item\label{itm:ch18-mc:numbers}
 in top degree, $c_1^\delta(\overline L)^{\wedge n}$ is the measure
 \eqref{eq:appJ-cl-measure},
 \[
  \sum_Eb_E\bigl(\cL|_E\bigr)^{n}\,\delta_{x_E},
 \]
 supported on the divisorial points of $\cX_s$, of total mass $(L^n)_X$;
\item\label{itm:ch18-mc:curve}
 if $n=1$, the corner cycle is the $0$-cycle $\sum_E\deg(\cL|_E)\{x_E\}$,
 and the atom \eqref{eq:ch18-atom} of the metric weight at $x_E$ equals
 $\deg(\cL|_E)-\deg(\cL_0|_E)$ for any reference model bundle $\cL_0$.
\end{enumerate}
\end{theorem}

\begin{proof}
\cref{itm:ch18-mc:delta} By \Cref{thm:gm-metric-comparison} the model metric
is a formal metric, hence PL, so on each tropical chart the weight is
piecewise affine and \eqref{eq:tropical-delta-pl} applies; that the local
preforms are compatible under refinement is
\Cref{lem:ch18-delta-metric-gluing}\,\cref{itm:ch18-dm:delta}, which for
formal metrics is \cite[Proposition~8.11]{GK17}.

\cref{itm:ch18-mc:PL} A piecewise affine function has vanishing second
derivative on the relative interior of each cell, so
$\dprime_P\ddprime_P\phi=0$.

\cref{itm:ch18-mc:numbers} This is \Cref{thm:gk-delta-cl-measure} applied
with $\overline L_1=\cdots=\overline L_n=\overline L$, together with the
normalization \eqref{eq:appJ-cl-measure}; the total mass is
\Cref{cor:ch18-delta-total-mass}.

\cref{itm:ch18-mc:curve} For $n=1$ the corner cycle is a $0$-cycle on the
skeleton, supported at the vertices, and by \cref{itm:ch18-mc:numbers} its
weight at $x_E$ is $b_E\deg(\cL|_E)$; since $\cX$ is a model of a curve and
$b_E$ is the multiplicity of $E$ in $\cX_s$, the normalization
$x_E=b_E^{-1}\ord_E$ absorbs $b_E$, giving $\deg(\cL|_E)$ in the
normalization of \eqref{eq:appJ-cl-measure}.  For the last clause, write
$\cL=\cL_0\otimes\cO(D)$; the weights differ by $\varphi_D$, whose corner
cycle at $x_E$ is $\operatorname{at}_{x_E}(\varphi_D)$ by
\eqref{eq:ch18-corner-weight}, and the degrees differ by
$\deg(\cO(D)|_E)=(D\cdot E)$.  Equality of these two numbers is
\Cref{prop:skeletal-curvature-curve}, and it is the identity
\eqref{eq:cld-curve-graph-laplacian}.
\end{proof}

\begin{corollary}
\label{cor:ch18-model-psh}
In the situation of \Cref{thm:ch18-model-corner}, the following are
equivalent:
\begin{enumerate}
\item $\cL$ is nef on $\cX_s$;
\item $c_1^\delta(\overline L)$ is a positive delta-form;
\item on every skeleton, the metric weight is convex on each face and its
 corner cycle is effective;
\item the metric is semipositive in the sense of
 \Cref{sec:metrics-semipositive}.
\end{enumerate}
\end{corollary}

\begin{proof}
(1)$\Leftrightarrow$(2) and (2)$\Leftrightarrow$(3) are
\Cref{thm:gk-formal-positivity-comparison} and
\Cref{thm:gk-tropical-frame-positivity}, applied after base change to an
algebraically closed complete extension, which is harmless because both
nefness and effectivity of the corner cycle are preserved and reflected by
such a base change: nefness by
\Cref{prop:semipositive-stability}, and the corner cycle by
\Cref{prop:trop-calibration-base-change}.  (1)$\Leftrightarrow$(4) is the
definition of semipositivity of a model metric together with
\Cref{thm:model-semipositive-closed}.
\end{proof}

\subsection{Worked comparisons}
\label{sec:ch18-model-worked}

\begin{example}[The blowup of $\mathbf P^1$, all three languages]
\label{ex:ch18-compare-p1}
Take $X=\mathbf P^1_K$, $\cX'$ the blowup of one closed point of the special
fibre, with $\cX'_s=\widetilde C+E$ and $b_{\widetilde C}=b_E=1$, as in
\Cref{ex:metrics-p1-blowup}.
\begin{center}
\small
\begin{tabular}{p{.27\textwidth}|p{.32\textwidth}|p{.32\textwidth}}
\textbf{Language}&\textbf{Datum}&\textbf{Value}\\ \hline
model
 &$\cL'=\rho^*\cO(1)\otimes\cO(-E)$
 &$\deg(\cL'|_{\widetilde C})=0$, $\deg(\cL'|_E)=1$\\
potential
 &weight $\max\{-1,\log|T|\}$ on the apartment
 &slopes $0$ and $1$; atom $1$ at $x_E$, $0$ at $\zeta_{0,1}$\\
delta-form
 &$c_1^\delta=\{x_E\}$
 &effective, mass $1=\deg\cO(1)$\\
\end{tabular}
\end{center}
All three rows agree, and the metric is semipositive by
\Cref{cor:ch18-model-psh} because both degrees are $\geq0$.  If instead one
takes $\cL''=\rho^*\cO(1)\otimes\cO(E)$, the degrees are
$\deg(\cL''|_{\widetilde C})=2$ and $\deg(\cL''|_E)=-1$; the corner cycle is
$2\{\zeta_{0,1}\}-\{x_E\}$, of mass $1$ but not effective, and the metric is
not semipositive.  Again mass alone decides nothing.
\end{example}

\begin{example}[The chain, all three languages]
\label{ex:ch18-compare-chain}
For the chain of \Cref{ex:ch18-chain-delta} with
$D=a_0C_0+a_1E+a_2C_2$ the three languages read as follows.
\begin{center}
\small
\begin{tabular}{p{.24\textwidth}|p{.36\textwidth}|p{.32\textwidth}}
\textbf{Language}&\textbf{Datum}&\textbf{Value}\\ \hline
model
 &$\deg(\cL|_{C_0}),\deg(\cL|_E),$\newline
   $\deg(\cL|_{C_2})$
 &$1+a_1-a_0$, $a_0+a_2-2a_1$, $1+a_1-a_2$\\
potential
 &PL function on $[0,2]$ with values $a_0,a_1,a_2$
 &atoms \eqref{eq:ch18-chain-atoms}\\
delta-form
 &corner cycle of that potential plus $\{v_0\}+\{v_2\}$
 &\eqref{eq:metrics-cl-chain-measure}, mass $2$\\
\end{tabular}
\end{center}
For $a=(0,1,0)$ the numbers are $2,-2,2$: total $2$, not effective.  For
$a=(0,-1,0)$ they are $0,2,0$: total $2$, effective, and the metric is
semipositive with Monge--Amp\`ere measure $2\delta_{v_1}$, concentrated at
the middle vertex.  This is the standard example of a semipositive model
metric whose measure is a single atom, and it is the smallest instance of
the phenomenon studied in \Cref{chap:bfj-continuity}.
\end{example}

\begin{example}[Toric potentials]
\label{ex:ch18-compare-toric}
Let $X$ be a proper toric variety with dense torus $T$ and let $\phi$ be a
toric model metric on an ample $L$, so that its weight on the skeleton
$N_{\mathbf R}$ is a piecewise affine concave function
$\psi$ whose Legendre transform is the support function of the polytope
$P_L$; see \eqref{eq:toric-metric-potential}.  The corner cycle of $-\psi$
is the codimension-one part of the corner locus of $-\psi$, a balanced
one-dimensional cycle whose weights are lattice lengths of edges dual to the
walls; the top delta-form product is by \eqref{eq:ch18-stable-weight} the
real Monge--Amp\`ere measure of $-\psi$, which by
\eqref{eq:appJ-alexandrov} is $n!\vol(P_L)$ in total mass, that is,
$(L^n)$.  This recovers \eqref{eq:metrics-toric-measure} and
\Cref{thm:toric-positivity}, and it is the tropical form of the comparison
\eqref{eq:cld-toric-ma-comparison}.  In the smallest case, $X=\mathbf P^2$
and $\psi=-\max\{0,x_1,x_2\}$, the corner cycle is the standard tropical
line of \Cref{ex:ch18-corner-max} and the top product is $\delta_0$, of mass
$1=(\cO(1)^2)$, as computed in
\Cref{ex:ch18-two-dim-product}.
\end{example}

\subsection{What the comparison does not say}
\label{sec:ch18-model-limits}

\begin{warning}{Three boundaries of the comparison}
\Cref{thm:ch18-model-corner} is a statement about model metrics.  Three
extensions of it are \emph{not} proved here and should not be assumed.
\begin{enumerate}
\item \emph{Continuous semipositive metrics.}  A uniform limit of model
 metrics has a curvature current and a Monge--Amp\`ere measure, by
 \Cref{thm:cl-continuous-extension} and
 \Cref{chap:cld-bedford-taylor}; it need not have a corner cycle, and the
 statement ``its Chern delta-form is positive'' is not defined.  The
 correct comparison in that generality is
 \Cref{chap:compare-ma}.
\item \emph{Arbitrary psh metrics.}  For a psh metric of the CLD hierarchy
 of \Cref{def:ch16-psh-filtration} beyond the continuous stage, no
 polyhedral presentation exists and none of this section applies.
\item \emph{Trivially valued base.}  The whole section assumes a nontrivial
 valuation; the homogeneous theory of \Cref{part:bj} has its own notion of
 PL function, tied to the scaling action, and the two are compared in
 \Cref{chap:master-dictionary}, not here.
\end{enumerate}
\end{warning}

\section{Chapter dictionary}
\label{sec:ch18-dictionary}

\subsection{From the objects of this chapter to their use}
\label{sec:ch18-dict-forward}

\begin{center}
\small
\begin{tabular}{p{.27\textwidth}|p{.35\textwidth}|p{.30\textwidth}}
\textbf{Object}&\textbf{Established property}&\textbf{First use}\\ \hline
atom $\operatorname{at}_v$
 &vanishes for smooth functions
  (\Cref{prop:ch18-smooth-zero-atom})
 &\Cref{ex:ch18-gauge-atom},
  \Cref{sec:ch18-harmonic-motivation}\\
corner cycle $\varphi\cdot C$
 &balanced, subdivision independent, affine invariant
  (\Cref{lem:ch18-corner-welldefined})
 &\Cref{thm:gk-poincare-lelong},
  \Cref{thm:ch18-model-corner}\\
delta-preform
 &polyhedral current presented by $\alpha\wedge\delta_C$
  (\eqref{eq:delta-preform})
 &\Cref{thm:delta-preform-algebra}\\
associated current
 &determines the collected coefficients in a fixed ambient space
  (\Cref{prop:ch18-current-faithful})
 &\Cref{def:ch18-delta-form-euclidean}\\
delta-form
 &the global subalgebra $B_{\mathrm{GK}}$ is differential
  (\Cref{prop:ch18-global-pb-structure})
 &\Cref{sec:ch18-delta-analytic}\\
delta-current
 &restricts canonically to a CLD current; no injectivity claim
  (\Cref{prop:ch18-delta-vs-cld-currents})
 &\Cref{thm:gk-poincare-lelong}\\
positivity of a delta-form
 &cellwise in fixed ambient; functorial under algebraic pullback
  (\Cref{lem:ch18-positivity-consistent})
 &\Cref{thm:gk-tropical-frame-positivity}\\
Chern delta-form $c_1^\delta$
 &exists for delta-metrics
  (\Cref{lem:ch18-delta-metric-gluing})
 &\Cref{thm:gk-delta-cl-measure}\\
mixed delta-form measure
 &is the Chambert--Loir measure, mass $(L_1\cdots L_n)$
  (\Cref{cor:ch18-delta-total-mass})
 &\Cref{chap:compare-ma}\\
harmonic PL function
 &functorial under pullback and base change
  (\Cref{prop:ch18-harmonic-functorial})
 &\Cref{thm:harmonic-balancing}\\
harmonic tropicalization
 &canonical weights, balanced off the boundary
  (\Cref{thm:harmonic-balancing})
 &\Cref{prop:ch18-weakly-smooth-welldefined}\\
weakly smooth form
 &bigraded differential algebra, contains CLD forms
  (\eqref{eq:cld-weak-inclusion})
 &\Cref{thm:abelian-harmonic-tropicalization}\\
model function
 &Chern delta-form is its corner cycle
  (\Cref{thm:ch18-model-corner})
 &\Cref{chap:master-dictionary}\\
\end{tabular}
\end{center}

\subsection{The smallest instance of each phenomenon}
\label{sec:ch18-dict-backward}

\begin{center}
\small
\begin{tabular}{p{.30\textwidth}|p{.34\textwidth}|p{.28\textwidth}}
\textbf{Phenomenon}&\textbf{Smallest instance}&\textbf{Computed in}\\
\hline
curvature invisible to smooth forms
 &gauge on the tropical line, atom $3$
 &\Cref{ex:ch18-gauge-atom}\\
curvature of a max-plus polynomial
 &$\max\{0,x_1,x_2\}$, corner cycle $=$ tropical line
 &\Cref{ex:ch18-corner-max}\\
stable self-intersection
 &$L\cdot L=1\cdot\{0\}$; $2L\cdot2L=4\cdot\{0\}$
 &\Cref{ex:ch18-two-lines-numbers}\\
a lattice index $\neq1$
 &$\max\{0,x_1+x_2\}$, $\max\{0,x_1-x_2\}$: product $2\{0\}$
 &\Cref{ex:ch18-nonunimodular}\\
noneffective corner cycle
 &$\min\{0,t\}$: weight $-1$
 &\Cref{ex:ch18-corner-dim-one}\\
positive mass, not semipositive
 &chain with $a=(0,1,0)$: $2,-2,2$
 &\Cref{ex:ch18-zhang-chain}\\
Poincar\'e--Lelong on a curve
 &$(T-a)/(T-b)$: $\delta_a-\delta_b$, mass $0$
 &\Cref{ex:ch18-pl-p1}\\
harmonic but not smooth
 &$|x_1|-|x_2|$ on the cross-fan, slopes $1,1,-1,-1$
 &\Cref{ex:ch18-harmonic-not-smooth}\\
harmonic chart weights
 &image line with weight $2$
 &\Cref{ex:ch18-cross-chart}\\
balancing failure at a boundary
 &closed disc, $h=-\log|T|$, defect $1$
 &\Cref{ex:ch18-harmonic-disc}\\
model $=$ potential $=$ delta-form
 &blowup of $\mathbf P^1$: $0,1$; $c_1^\delta=\{x_E\}$
 &\Cref{ex:ch18-compare-p1}\\
single-atom Monge--Amp\`ere
 &chain with $a=(0,-1,0)$: $2\delta_{v_1}$
 &\Cref{ex:ch18-compare-chain}\\
\end{tabular}
\end{center}

\subsection{Hypothesis ledger}
\label{sec:ch18-hypothesis-ledger}

\begin{center}
\small
\begin{tabular}{p{.30\textwidth}|p{.62\textwidth}}
\textbf{Hypothesis}&\textbf{Where it is used, and what fails without it}\\
\hline
$C$ balanced
 &\Cref{prop:ch18-smooth-zero-atom},
  \Cref{lem:ch18-differential-delta},
  \Cref{prop:ch18-weakly-smooth-welldefined}; without it, smooth functions
  have nonzero atoms and $\delta_C$ is not closed\\
weights constant
 &\Cref{lem:ch18-differential-delta}; with smooth weights an extra
  $\rd m_\sigma$ term appears (\Cref{rem:ch18-two-weights})\\
representatives meet properly
 &\Cref{thm:delta-preform-algebra}; otherwise the product is defined only
  after a generic translation, and the limit must be shown to exist\\
$K$ algebraically closed, nontrivially valued
 &\Cref{thm:gk-formal-positivity-comparison}: curve lifting; also the very
  affine chart system of \Cref{lem:ch18-very-affine-charts}\\
metric piecewise linear
 &\Cref{thm:gk-formal-positivity-comparison},
  \Cref{thm:ch18-model-corner}; a continuous psh metric has no residual
  bundle and no corner cycle\\
metric a delta-metric
 &\Cref{lem:ch18-delta-metric-gluing}; the piecewise smooth part and the
  residual codimension-one form always exist, while the full ordinary Chern
  delta-form requires the ambient Hessian lift\\
$h$ harmonic
 &\Cref{thm:harmonic-balancing}; without harmonicity the residual degree
  survives Stokes and the image is unbalanced\\
boundary excluded
 &\Cref{thm:harmonic-balancing},
  \Cref{prop:ch18-delta-integration}\,\cref{itm:ch18-di:stokes};
  \Cref{ex:ch18-harmonic-disc} shows the defect\\
$X$ good and strict
 &\Crefrange{sec:harmonic-trop}{sec:ch18-classical}; residual line bundles
  are defined through reductions of germs\\
$K$ discretely valued
 &\Cref{thm:ch18-model-corner}: integer multiplicities $b_E$ and the
  normalization $\varphi_{\cX_s}=1$ of
  \eqref{eq:appJ-special-fibre-test}\\
real-PL zero-current converse
 &\Cref{prop:ch18-real-pl-zero-current-harmonic}: $X$ is separated,
  boundaryless, strict, and pure-dimensional; these are the hypotheses used
  for the closed canonical skeletal current and for
  \cite[Theorem~18.8.4]{CLD25}; no arbitrary continuous function is covered\\
\end{tabular}
\end{center}

\subsection{Comparison ledger}
\label{sec:ch18-ledger}

The proved implications needed later can be summarized as follows.
\begin{center}
\begin{tabular}{p{0.24\textwidth}p{0.31\textwidth}p{0.31\textwidth}}
\toprule
Object & Proven comparison & Scope or warning \\
\midrule
CLD smooth form & injects into a weakly smooth form & good strict analytic,
nontrivial valuation; \cite[Proposition~10.14]{GJR26} \\
Formal Chern current & is a Chern delta-form & algebraic variety over an
algebraically closed nontrivially valued field; \cite{GK17} \\
Semipositive formal metric & positive Chern delta-form, functorially psh, and
curve-positive & proper algebraic formal locus; \Cref{thm:gk-formal-positivity-comparison} \\
Semipositive \(\mathbf R\)-PL function & classically psh; boundaryless also
CLD-v2 psh & good strict analytic; nontrivial valuation;
\Cref{thm:gr-classical-pl,thm:verified-pl-psh-equivalence} \\
CLD v2 psh function & classically psh & inclusion only; no converse asserted
for arbitrary singular functions \\
$\mathbf R$-PL function with zero CLD current & harmonic and classically
pluriharmonic & good separated boundaryless strict pure-dimensional space;
\Cref{prop:ch18-real-pl-zero-current-harmonic}; arbitrary continuous
functions remain outside the conclusion \\
Weakly smooth form vs. delta-form & no general equality used & different
coordinate and coefficient enlargements \\
Classically psh function & maximum principles and curve functoriality & no
automatic Monge--Amp\`ere product \\
\bottomrule
\end{tabular}
\end{center}

\section*{Source notes}

\paragraph{Summary of provenance.}
Delta-forms, tropical intersections, Poincar\'e--Lelong, Chern delta-forms,
model measures, and local heights come from \cite{GK17}.  The formal
semipositivity equivalences and tropical-frame positivity criterion are
\cite{GK19}.  Harmonic tropicalizations and weakly smooth forms are from
\cite{GJR26}.  The residue-nef curve test and its stability theorem are from
\cite{GR25PL}; the classically psh comparison is \cite{GR25Classical}.
The real-PL zero-current converse of
\Cref{prop:ch18-real-pl-zero-current-harmonic} is a book-level deduction from
the calibrated skeletal-current results of \cite{CLD25}, the residue and
common-tropicalization results of \cite{GJR26}, and the classical comparison
of \cite{GR25Classical}; no novelty or priority claim is made.
These sources use different ambient categories, so every theorem above keeps
its original hypotheses rather than presenting a universal equivalence.

\paragraph{What belongs to the CLD core and what does not.}
The plan of \Cref{part:cld} requires that the reader be able to see the
boundary of the original theory.  It is this:
\Crefrange{chap:cld-forms}{chap:cld-metrized-bundles} contain the
Chambert-Loir--Ducros theory of \cite{CLD25}---smooth forms, currents,
Poincar\'e--Lelong, Bedford--Taylor products, Chern currents---and nothing in
those chapters uses a delta-form or a harmonic chart, with the single
exception of the harmonic-chart compatibility of
\Cref{sec:ch16-harmonic-charts}, which is \cite[\S7.8]{CLD25} and belongs to
the core.  Everything in the present chapter is later work of other authors.
Concretely:
\begin{itemize}
\item \Crefrange{sec:delta-tropical}{sec:delta-positivity} is
 Gubler--K\"unnemann; the objects there do not occur in \cite{CLD25} at all,
 and the tropical intersection product is imported from tropical geometry
 (\cite{MR18Tropical} and the references of \Cref{app:tropical-algebra}),
 not from the CLD calculus.
\item \Crefrange{sec:harmonic-trop}{sec:weak-smooth} is
 Gubler--Jell--Rabinoff; here the objects of \cite{CLD25} are contained in
 the new ones by \eqref{eq:cld-weak-inclusion}, so the relation is an
 enlargement rather than a parallel construction, but the enlargement is not
 in \cite{CLD25}.
\item \Cref{sec:ch18-classical} is Gubler--Rabinoff; the notion of
 classically psh function is defined by curve tests, which is a different
 mechanism from the approximation-generated sheaf of \cite{CLD25}, and the
 two are compared only where a theorem says so.
\item \Cref{sec:ch18-model-comparison} is a comparison written for this
 book.  Its inputs are \Cref{chap:metrics}, \Cref{chap:skeleta} and the
 Gubler--K\"unnemann results of this chapter; the statement
 \Cref{thm:ch18-model-corner} assembles them, and its individual
 ingredients are attributed in its proof.
\end{itemize}

\paragraph{Order of discovery.}
The chronology is worth stating because it explains the shape of the
theory.  Lagerberg's superforms and the CLD calculus came first, and were
designed for a differential-geometric purpose.  Gubler and K\"unnemann's
delta-forms \cite{GK17} were designed for an Arakelov-theoretic purpose:
the local heights of \Cref{sec:ch18-green} are the target, and the
enlargement of coefficients is what makes a $*$-product available.  Their
positivity paper \cite{GK19} then compared the resulting notion with
Zhang's semipositivity \cite{Zhang1995} and with the metric-theoretic
notions of \cite{GM19}.  Gubler, Jell and Rabinoff \cite{GJR26} started from
a cohomological question---the failure of Poincar\'e duality for the smooth
Dolbeault complex, whose local theory is Jell's \cite{Jell16}---and enlarged
the charts instead.  The curve-test papers \cite{GR25PL,GR25Classical} are
later still, and they belong to a different tradition, that of Thuillier's
potential theory on curves and its comparison with the CLD sheaf
\cite{Wanner19}.  A reader who keeps these four purposes apart will not be
tempted to look for a single theory containing all four constructions.

\paragraph{The motivating computation.}
\Cref{prop:ch18-smooth-zero-atom} and \Cref{ex:ch18-gauge-atom} are
elementary and are recorded for orientation, as
\Cref{rem:ch18-gauge-provenance} states; they are not theorems of this book
and carry no priority claim.  The vanishing of the atom is
\Cref{ex:ch10-smooth-at-vertex} restated; the gauge is the standard
Minkowski functional of a triangle; the concentration statement is
\Cref{thm:cld-bt-continuity}.  The pair is included because the sentence
``smooth forms do not contain all model-theoretic curvature'' is otherwise
an assertion without a number attached to it.

\paragraph{Normalization audit.}
The Poincar\'e--Lelong constants of this chapter were checked against
\Cref{app:normalizations} in \Cref{cor:ch18-pl-restricts}, against
\Cref{thm:cld-poincare-lelong} in the same corollary, and against
\Cref{cor:ch15-pl-section} for the section form.  The Chambert--Loir
normalization was checked in \Cref{cor:ch18-delta-total-mass} against
\eqref{eq:appJ-cl-total-mass}, and in numbers in
\Cref{ex:ch18-chain-delta,ex:ch18-p1-blowup-delta,ex:ch18-compare-toric}.
The lattice indices, which are the one place where a hidden factor could
enter, were tested in \Cref{ex:ch18-two-lines-numbers} and
\Cref{ex:ch18-nonunimodular}.

\paragraph{Statements quoted without proof.}
Four statements in this chapter are used with a citation and are not proved
here: the analytic tropical fibre-calibration theorem
\cite[Proposition~4.6.6]{CLD25} inside
\Cref{lem:ch18-fibre-cylinder}; that
formal metrics are delta-metrics, \cite[Proposition~8.11]{GK17}, inside
\Cref{thm:ch18-model-corner}\,\cref{itm:ch18-mc:delta}; the residual
mixed-degree comparison \cite[Theorem~10.5]{GK17} inside
\Cref{thm:gk-delta-cl-measure}; and the curve-lifting theorem used in
\Cref{thm:gk-formal-positivity-comparison} and
\Cref{thm:gr-pl-curve-test}.  Each is a deep input whose proof would require
machinery not developed in this book, and in each case the hypotheses under
which we invoke it are exactly those of the source.  The remaining citations
in the chapter are attributions of statements which are proved here, or
pointers to the literature made in prose.
